%%%%%%%%%%%%%%%%%%%%%%%%%%%%%%%%%%%%%%%%%%%%%%%%%%%%%%%%%%%%%%%
%  "Magic Rabbit" Lecture Notes / Book Template
%  Updated for MATH 61 Structure (Weeks 1-6 + Solutions)
%%%%%%%%%%%%%%%%%%%%%%%%%%%%%%%%%%%%%%%%%%%%%%%%%%%%%%%%%%%%%%%

\documentclass[11pt, letterpaper, openany]{book} 
% 'openany' prevents blank pages between chapters

%% === PACKAGES ===
\usepackage[utf8]{inputenc}
\usepackage[T1]{fontenc}
\usepackage[margin=1in, top=1.2in, bottom=1in]{geometry}
\usepackage{amsmath, amssymb, amsthm}
\usepackage{enumitem}
\usepackage{graphicx}
\usepackage{fancyhdr}
\usepackage{etoolbox}
\usepackage{tocloft} 
\usepackage{titlesec}
\usepackage{tikz}
\usepackage{tikz-cd}
\usetikzlibrary{arrows.meta, positioning, calc} 

%% === TYPOGRAPHY ===
\usepackage[osf,sc]{mathpazo} % Palatino (Serif)
\linespread{1.05}
\usepackage[scale=0.95]{cabin} % Cabin (Sans-Serif)
\usepackage[varqu,varl]{inconsolata} % Monospace

%% === COLOR PALETTE ===
\usepackage[dvipsnames, svgnames, table]{xcolor}
\definecolor{MagicTeal}{HTML}{1F3A4D}
\definecolor{MagicPink}{HTML}{FCAEAC}
\definecolor{MagicPeach}{HTML}{FDD0B1}
\definecolor{MagicMint}{HTML}{B5EAD7}
\definecolor{MagicBlue}{HTML}{A2D2FF}
\definecolor{MagicCream}{HTML}{FAFAF7}

\pagecolor{MagicCream} % Off-white background

%% === HYPERLINKS ===
\usepackage[colorlinks=true, linkcolor=MagicTeal, urlcolor=MagicBlue, citecolor=MagicPink]{hyperref}

%% === HEADER & FOOTER DESIGN ===
\pagestyle{fancy}
\fancyhf{}
% Left Header: Current Chapter/Week
\lhead{\sffamily\color{MagicTeal}\small \leftmark}
% Right Header: Page Number
\rhead{\sffamily\color{MagicTeal}\bfseries\thepage}
\renewcommand{\headrulewidth}{0pt} 

% Custom colorful header underline
\newcommand{\colorrule}{%
\hbox to\headwidth{%
  \color{MagicPink}\leaders\hrule height 1.5pt\hskip .25\headwidth\hfill%
  \color{MagicMint}\leaders\hrule height 1.5pt\hskip .25\headwidth\hfill%
  \color{MagicPeach}\leaders\hrule height 1.5pt\hskip .25\headwidth\hfill%
  \color{MagicBlue}\leaders\hrule height 1.5pt\hskip .25\headwidth\hfill%
}}
\renewcommand{\headrule}{\colorrule}

%% === TITLE & SECTION STYLING ===

% -- PART FORMAT --
\titleformat{\part}[display]
{\centering\sffamily\bfseries\color{MagicTeal}}
{\Huge\color{MagicPink} \partname\ \thepart}
{0pt}
{\Huge}

% -- CHAPTER FORMAT (Used for Weeks) --
\titleformat{\chapter}[display]
  {\sffamily\bfseries\color{MagicTeal}}
  {\flushright\fontsize{40}{40}\selectfont\color{MagicPink!50}Week \thechapter} 
  {-10pt}
  {\Huge}
  [\vspace{1ex}{\color{MagicMint}\titlerule[3pt]}] 

% -- SECTION FORMAT --
\titleformat{\section}
  {\Large\sffamily\bfseries\color{MagicTeal}}
  {\thesection}{1em}{}

%% === BOXES (Theorems, etc.) ===
\usepackage{tcolorbox}
\tcbuselibrary{skins, breakable, theorems}

% 1. MAIN THEOREM (Pink)
\newtcbtheorem[number within=chapter]{thm}{Theorem}{%
    enhanced, breakable,
    colframe=MagicPink, colback=MagicPink!10!MagicCream,
    coltitle=MagicTeal!90!black, fonttitle=\bfseries\sffamily,
    attach boxed title to top left={yshift=-2mm, xshift=5mm},
    boxed title style={colback=MagicPink, frame hidden, arc=1mm},
    separator sign={\ $\cdot$\ }, fontupper=\itshape
}{thm}

% 1b. LEMMA (Same Pink style as Theorem)
\newtcbtheorem[use counter from=thm]{lem}{Lemma}{%
    enhanced, breakable,
    colframe=MagicPink, colback=MagicPink!10!MagicCream,
    coltitle=MagicTeal!90!black, fonttitle=\bfseries\sffamily,
    attach boxed title to top left={yshift=-2mm, xshift=5mm},
    boxed title style={colback=MagicPink, frame hidden, arc=1mm},
    separator sign={\ $\cdot$\ }, fontupper=\itshape
}{lem}

% 2. DEFINITION (Mint)
\newtcbtheorem[number within=chapter]{defn}{Definition}{%
    enhanced, breakable,
    colframe=MagicMint!80!MagicTeal, colback=MagicMint!15!MagicCream,
    coltitle=MagicTeal!90!black, fonttitle=\bfseries\sffamily,
    attach boxed title to top left={yshift=-2mm, xshift=5mm},
    boxed title style={colback=MagicMint, frame hidden, arc=1mm},
    separator sign={\ $\cdot$\ },
}{def}

% 3. EXAMPLE/REMARK (Blue side bar)
\newtcolorbox{example}[1][]{
    blanker, breakable, left=1em,
    borderline west={3pt}{0pt}{MagicBlue},
    before upper={\textcolor{MagicBlue!80!black}{\sffamily\bfseries Example:}\space},
    #1
}

% 4. PROOF (Clean Left-bar)
\renewenvironment{proof}[1][\proofname]{%
  \tcolorbox[blanker, breakable, left=1em,
    before skip=1em, after skip=1em,
    borderline west={2pt}{0pt}{MagicTeal!50},
    colback=MagicCream,
    before upper={\textcolor{MagicTeal}{\sffamily\bfseries #1.}\space}
  ]%
}{%
  \hfill\textcolor{MagicTeal}{\ensuremath{\blacksquare}}\endtcolorbox%
}

% 5. REMARK (Peach, Light Background, Run-in Title)
\newtcolorbox{remark}[1][]{
    enhanced, 
    breakable,
    frame hidden,
    colback=MagicPeach!15!MagicCream,
    borderline west={3pt}{0pt}{MagicPeach},
    coltitle=MagicTeal,
    fonttitle=\bfseries\sffamily,
    attach title to upper={\ },
    title={Remark.},
    after title={\ifstrempty{#1}{}{ \textbf{(#1)}}}, 
    sharp corners,
    left=1em, right=1em, top=1ex, bottom=1ex
}


% 5b. COROLLARY (Same Pink style)
\newtcbtheorem[use counter from=thm]{cor}{Corollary}{%
    enhanced, breakable,
    colframe=MagicPink, colback=MagicPink!10!MagicCream,
    coltitle=MagicTeal!90!black, fonttitle=\bfseries\sffamily,
    attach boxed title to top left={yshift=-2mm, xshift=5mm},
    boxed title style={colback=MagicPink, frame hidden, arc=1mm},
    separator sign={\ $\cdot$\ }, fontupper=\itshape
}{cor}

% 6. EXERCISES (Final Polished Version)
\newtcolorbox{exercises}{enhanced,breakable,colback=MagicTeal!5!MagicCream,colframe=MagicTeal!5!MagicCream,coltitle=MagicTeal,fonttitle=\bfseries\Large\sffamily,title={Exercises},sharp corners,boxrule=0pt,toptitle=1ex,bottomtitle=0.5ex,left=1.5em,right=1.5em,top=0.5em,bottom=1.5em,before skip=2em,after skip=2em}

% 6b. NUMBERED EXERCISE ENVIRONMENT (for individual exercises)
\newcounter{exercise}[chapter]
\renewcommand{\theexercise}{\thechapter.\arabic{exercise}}
\newenvironment{exercise}{%
    \refstepcounter{exercise}%
    \par\noindent\textbf{\sffamily\textcolor{MagicTeal}{Exercise \theexercise.}}\space
}{\par\medskip}

% 7. SOLUTION (Blue Left-bar, hollow square at end)
\newenvironment{solution}[1][Solution]{%
  \tcolorbox[blanker, breakable, left=1em,
    before skip=1em, after skip=1em,
    borderline west={2pt}{0pt}{MagicBlue},
    colback=MagicCream,
    before upper={\textcolor{MagicBlue}{\sffamily\bfseries #1.}\space}
  ]%
}{%
  \hfill\textcolor{MagicBlue}{\ensuremath{\square}}\endtcolorbox%
}

%% === COVER PAGE MACRO ===
\newcommand{\makeMagicCover}[2]{
    \begin{titlepage}
        \pagecolor{MagicTeal} 
        \color{MagicCream}
        \vspace*{2in}
        \begin{center}
            {\fontsize{50}{60}\selectfont \sffamily\bfseries #1} \\[1em]
            {\fontsize{20}{25}\selectfont \sffamily #2} \\[3in]
            {\color{MagicPink}\Huge $\bullet$} \quad 
            {\color{MagicMint}\Huge $\bullet$} \quad 
            {\color{MagicPeach}\Huge $\bullet$} \quad 
            {\color{MagicBlue}\Huge $\bullet$}
        \end{center}
    \end{titlepage}
    \pagecolor{MagicCream} 
}


% Limit TOC depth to Sections only (hides 1.1.1)
\setcounter{tocdepth}{1}

%%%%%%%%%%%%%%%%%%%%%%%%%%%%%%%%%%%%%%%%%%%%%%%%%%%%%%%%%%%%%%%
% DOCUMENT CONTENT
%%%%%%%%%%%%%%%%%%%%%%%%%%%%%%%%%%%%%%%%%%%%%%%%%%%%%%%%%%%%%%%
\begin{document}

% --- COVER PAGE ---
\makeMagicCover{Lecture Notes}{Damian Burgess}

% --- TABLE OF CONTENTS ---
{
  \hypersetup{linkcolor=MagicTeal}
  \tableofcontents
}

% =======================================================
% PART I: MATH 61 - WEEKS 1-6
% =======================================================
\part{MATH 61CM: Course Notes}

% ==================== WEEK 1 ====================
\chapter{The Real Numbers}

\section{Analysis: Mon/}

\subsection{Introduction to $\mathbb{R}$}
Real analysis is built on the structure of fields — specifically the real numbers $\mathbb{R}$.

\begin{defn}{The Three Axioms for $\mathbb{R}$}{axioms_R}
The real numbers are the \textbf{unique} field that satisfy all three of the following axioms:
\begin{enumerate}
    \item \textbf{Arithmetic:} Field axioms.
    \item \textbf{Order:} Order axioms.
    \item \textbf{No Holes:} Completeness axiom.
\end{enumerate}
\end{defn}

\subsection{Field Axioms}
\begin{defn}{Field}{field_def}
A set combined with two binary operations $(\mathbb{F}, +, \cdot)$ is a \textbf{field} if it satisfies the following properties: 
\begin{enumerate}
    \item F1 (Commutativity): For all $a,b\in\mathbb{F}$, $$a+b=b+a\text{ and }ab=ba$$
    \item F2 (Associativity): For all $a,b,c\in\mathbb{F}$, $$(a+b)+c=a+(b+c)\text{ and } (ab)c=a(bc)$$
    \item F3 (Distributivity): For all $a,b,c\in\mathbb{F}$, $$a(b+c)=ab+ac$$
    \item F4 (Identities Exist): There are special elements, $0,1\in\mathbb{F}$ s.t. $$a+0=0+a=a\text{ and }a\cdot1=a\text{ for all }a\in\mathbb{F}$$
    \item F5 (Additive Inverse): For every $a\in\mathbb{F}$ there exists $-a$ s.t. $a+(-a)=0$
    \item F6 (Multiplicative Inverse): For every $a\neq{0}$, there exists $a^{-1}$ s.t. $a\cdot{a^{-1}}=1$
\end{enumerate}
\end{defn}
\begin{remark}
    In F4, it is assumed that $1\neq{0}$. 
\end{remark}
\begin{remark}
    A set combined with one binary operation that fulfills F1-F5 is called a commutative group, $(G,\circ)$. 
\end{remark}

\begin{lem}{Inverses are Unique}{inverses_unique}
For all $a,b,c\in\mathbb{F}$, if $a+b=0$ and $a+c=0$, then $b=c$. 
\end{lem}
\begin{proof}
    We start with $b=0+b$ (F4). This is equal to $b+0$ (F2) which is equal to $b+a+c$. Now we can use Associativity which states that $b+a+c=(b+a)+c\,(=a+b+c)$ (F1). $a+b=0$ so this is simply $0+c=c$ (F4). We have now proved our claim that $b=c$. 
\end{proof}

\begin{lem}{"Multiplying" by the zero element}{multiplicationbyzero}
    For all $a\in{\mathbb{F}}, 0\cdot{a}=0$
\end{lem}
\begin{proof}
    The idea behind this proof is quite simple. Note that $0\cdot{a}=(0+0)\cdot{a}=a\cdot(0+0)=a\cdot0+a\cdot0=0\cdot{a}+0\cdot{a}$. If we add $-(0\cdot{a})$ to both sides, we get that $0\cdot{a}-0\cdot{a}=0\cdot{a}+0\cdot{a}-0\cdot{a}$. This implies that $0=0\cdot{a}$
\end{proof}
\begin{lem}{The relationship between $-1$ and the additive inverse. }{-1andinverse}
    For all $a\in{\mathbb{F}}, (-1)\cdot{a}=-a$
\end{lem}
\begin{proof}
    Note that $a+(-1)\cdot{a}=1\cdot{a}+(-1)\cdot{a}=(1+(-1))\cdot{a}=0\cdot{a}$. Using lemma 1.2, we know that $0\cdot{a}=0$ which means that $a+(-1)\cdot{a}=0$. Since inverses are unique (lemma 1.1), we know that $(-1)\cdot{a}=-a$. 
\end{proof}

\subsection{Order Axioms}
\begin{defn}{Order Axioms}{order_axioms}
An ordered field has a positive set $\mathbb{P}$ satisfying:
\begin{enumerate}
    \item \textbf{Trichotomy:} For $a \in \mathbb{F}$, exactly one is true: $a=0$, $a \in \mathbb{P}$, $-a \in \mathbb{P}$.
    \item \textbf{Closure:} $\mathbb{P}$ is closed under addition and multiplication.
\end{enumerate}
\end{defn}
\begin{remark}
    Here is some more common notation: 
    \begin{enumerate}
        \item Write $a>0$ for $a\in\mathbb{P}$
        \item Write $a>b$ or $b<a$ if $a-b=a+(-b)>0$
        \item Write $a\geq{b}$ or $b\leq{a}$ if $(a>b\text{ or }a=b)$
    \end{enumerate}
\end{remark}
\begin{lem}{1>0}{1greaterthan0}
    $1$ is greater than $0$
\end{lem}
\begin{proof}
    Assume for the sake of contradiction that $1<0$ (We know that $1\neq0$ by definition). This implies that $(-1)>0$ by F4 and Trichotomy. By O2, $(-1)(-1)>0$. By Lemma 1.3, $(-1)(-1)=1$ and thus $1>0$ which contradicts our initial claim. Therefore $1>0$. 
\end{proof}

\subsection{Completeness}
\begin{defn}{Completeness}{completeness}
An ordered field $(F,+,\cdot)$is \textbf{complete} if any non-empty subset $S\subseteq\mathbb{F}$ bounded above has a \textbf{least upper bound} (supremum).
\end{defn}
\begin{remark}
    For now, take $\mathbb{R}$ to be the ONLY complete ordered field.
\end{remark}
\begin{defn}{Bounds}{lowerandupperbound}
    $S\subseteq\mathbb{F}$ is $\left\{\begin{array}{c}
        \text{bdd above} \\
        \text{bdd below}
    \end{array}\right\}$ if $\exists\alpha\in\mathbb{F}$ s.t. $\left\{\begin{array}{c}
        x\leq{\alpha} \\
        x \geq \alpha 
    \end{array}\right\}$, $\forall{x}\in\mathbb{F}$. \\ 
    Such $\alpha$ is called a $\left\{\begin{array}{c}
        \text{an upper bound} \\
        \text{a lower bound}
    \end{array}\right\}$
\end{defn}
\begin{remark}
    $S\subseteq\mathbb{F}$ is bounded if it is bounded above and below. 
\end{remark}
\begin{defn}{Supremum}{Sup}
    Let $S\subseteq{\mathbb{F}}$. $\beta\in\mathbb{F}$ is a least upper bound if the following both hold: \begin{enumerate}
        \item $\beta$ is an upper bound of $S$
        \item If $\alpha$ is an upper bound of $S$ then $\beta\leq\alpha$. Similarly for the greatest lower bound. 
    \end{enumerate}
    We denote the least upper bound as $\sup{S}$ and the greatest lower bound as $\inf{S}$. 
\end{defn}
\begin{remark}
    $\sup{S}$ need not be in $S$. 
\end{remark}
\begin{defn}{$\mathbb{N,Q,Z}$}{natural,rational,integer}
    \begin{enumerate}
        \item The \textbf{natural numbers} are defined as $\mathbb{N}=\{1,1+1,1+1+1,...\}$
        \item The \textbf{integers} are defined as $\mathbb{Z}=\mathbb{N}\cup\{0\}\cup\{-n:n\in\mathbb{N}\}\subseteq{\mathbb{R}}$ 
        \item The \textbf{rational numbers} are defined as $\mathbb{Q}=\left\{\frac{a}{b}:a,b\in\mathbb{Z},b\neq{0}\right\}$
    \end{enumerate}
\end{defn}
\begin{thm}{Archimedean Principle}{arch_principle}
$\mathbb{N}$ is not bounded above. Equivalently, for any two positive real numbers $a,b\in\mathbb{R}$, there exists a natural number $n$ s.t. $a<n\cdot{b}$. 
\end{thm}
\begin{proof}
    Assume for the sake of contradiction that $\mathbb{N}$ was bounded above. If it were, since $\mathbb{N}$ is a subset of $\mathbb{R}$, we know that there exists $\alpha=\sup{\mathbb{N}}$ because of completeness. $\alpha=\sup{\mathbb{N}}\implies{n\leq{\alpha}}\,\forall{n}\in\mathbb{N}\implies{n+1}\leq\alpha\implies{n\leq{\alpha-1}}$ which contradicts our assumption that $\alpha$ was the $\sup$. 
\end{proof}

\subsection{Sequences}
\begin{defn}{Sequence}{sequence_def}
    A \textbf{sequence} of real numbers is a map $\mathbb{N}\to\mathbb{R}$. We denote this by $\{a_n\}^\infty_{n=1}$, where $a_n=a(n)$
\end{defn}
\begin{defn}{Limit of a Sequence}{limit_def}
Let $\{a_n\}^\infty_{n=1}\subseteq{\mathbb{R}}$. We say $\{a_n\}_{n=1}^\infty$ has a limit $L\in\mathbb{R}$, written $\lim_{n\to\infty}a_n=L$ if for all $\varepsilon>0$, there exists $N\in\mathbb{N}$ s.t. $|a_n-L|<\varepsilon$ for all $n\geq{N}$. 
\end{defn}

% Epsilon-delta visualization for sequences
\begin{center}
\begin{tikzpicture}[scale=0.8]
    % Axes
    \draw[->] (0,0) -- (7,0) node[right] {$n$};
    \draw[->] (0,0) -- (0,3.5) node[above] {};
    
    % The limit L
    \draw[dashed, MagicTeal] (0,2) -- (7,2) node[right] {$L$};
    
    % Epsilon band
    \draw[dashed, MagicPink] (0,2.6) -- (7,2.6);
    \draw[dashed, MagicPink] (0,1.4) -- (7,1.4);
    \draw[<->] (6.5,1.4) -- (6.5,2) node[midway, right] {$\varepsilon$};
    \draw[<->] (6.5,2) -- (6.5,2.6);
    
    % N marker
    \draw[thick] (3,-0.1) -- (3,0.1);
    \node at (3,-0.4) {$N$};
    
    % Sequence points
    \foreach \x/\y in {0.5/0.5, 1/1.2, 1.5/2.8, 2/1.6, 2.5/2.4}{
        \fill (\x,\y) circle (2pt);
    }
    % Points after N (inside epsilon band)
    \foreach \x/\y in {3.5/2.3, 4/1.8, 4.5/2.2, 5/1.9, 5.5/2.1, 6/2.05}{
        \fill[MagicTeal] (\x,\y) circle (2pt);
    }
    
    % Label
    \node at (1.5,3.2) {$n < N$};
    \node[MagicTeal] at (5,3.2) {$n \geq N$: $|a_n - L| < \varepsilon$};
\end{tikzpicture}
\end{center}

\begin{remark}
    We say $\{a_n\}^\infty_{n=1}\subseteq\mathbb{R}$ is $\left\{\begin{array}{c}\text{convergent if it has a limit}\\\text{divergent otherwise}\end{array}\right\}$
\end{remark}

\begin{thm}{Monotone Convergence}{monotone_conv}
If $\{a_n\}$ is bounded above and increasing ($a_n\leq{a_{n+1}}$), it converges to $\sup \{a_n\}$.
\end{thm}
\begin{proof}
    Let $\alpha=\sup{a_n}$, which we know exists because of the completeness property of the reals. Now we will prove there exists $a_j$ s.t. $\alpha>a_j>\alpha-\frac{1}{j}$. If such $a_j$ didn't exist, then $\alpha-\frac{1}{j}$ would be the supremum, which is a contradiction. Therefore such $a_j$ exists. Since $\alpha-\frac{1}{j}<a_j<\alpha$, we know that $|\alpha-a_j|<|\alpha-\frac{1}{j}|$. For all $\varepsilon>0$, we know there exists $\frac{1}{J}$ s.t. $\frac{1}{J}<\varepsilon$ due to the Archimedean property of the real numbers. This means that for all $j\geq{J}$ (or $\frac{1}{j}<\frac{1}{J}$), $|a_j-\alpha|\leq\frac{1}{j}\leq\frac{1}{J}<\varepsilon$
\end{proof}

\section{Linear Algebra: Thu/Fri}

\subsection{Vector Spaces}

\begin{defn}{Vector Space}{vector_space}
A \textbf{vector space} $V$ over a field $\mathbb{F}$ is a non-empty set of vectors with two operations:
\begin{align*}
+&: V \times V \to V \\
\cdot&: \mathbb{F} \times V \to V
\end{align*}
such that:
\begin{enumerate}
    \item $(V, +)$ forms a commutative group:
    \begin{itemize}
        \item[(A1)] $\forall v_1, v_2, v_3$: $(v_1 + v_2) + v_3 = v_1 + (v_2 + v_3)$
        \item[(A2)] $\exists$ zero vector $\vec{0}$: $v + \vec{0} = v$
        \item[(A3)] $\forall v$, $\exists (-v)$ s.t. $v + (-v) = \vec{0}$
        \item[(A4)] $v + w = w + v$
    \end{itemize}
    \item $\forall v \in V$: $1 \cdot v = v$
    \item $(\alpha \beta) \cdot v = \alpha \cdot (\beta \cdot v)$
    \item Distributive laws: $(\alpha + \beta) \cdot v = \alpha \cdot v + \beta \cdot v$ and $\alpha \cdot (v + w) = \alpha \cdot v + \alpha \cdot w$
\end{enumerate}
\end{defn}

\begin{lem}{Zero Times Anything}{zero_times}
$0 \cdot v = \vec{0}$ for all $v \in V$. (Same proof as in $\mathbb{F}$!)
\end{lem}

\begin{example}
\begin{enumerate}
    \item $\mathbb{F}^n = \{(x_1, \ldots, x_n) : x_i \in \mathbb{F}\}$ with componentwise operations.
    \item $\mathbb{F}^S$ = space of functions from $S \to \mathbb{F}$, with $(f+g)(x) = f(x) + g(x)$ and $(\alpha \cdot f)(x) = \alpha \cdot f(x)$.
    \item $P(\mathbb{R})$ = vector space of polynomials over $\mathbb{R}$.
\end{enumerate}
\end{example}

\subsection{Subspaces}

\begin{defn}{Subspace}{subspace}
$W$ is a \textbf{subspace} of $V$ if $W \subseteq V$ with the same operations and field, which forms a vector space on its own. I.e., $W$ is closed under the operations.
\end{defn}

\begin{lem}{Subspace Test}{subspace_test}
A subset $W$ of a vector space $V$ forms a subspace iff $\forall \alpha, \beta \in \mathbb{F}$, $\forall v, w \in W$: $\alpha v + \beta w \in W$.
\end{lem}

\begin{example}
Is $W = \{(x,y) : 2x - y + 3 = 0\}$ a subspace? \textbf{No}, because $W$ doesn't contain $\vec{0}$.

Is $W = \{(x,y,z) : 5x - 2y + 5z = 0\}$ a subspace? \textbf{Yes}, because for any two solutions, their linear combination is also a solution.
\end{example}

\subsection{Span and Linear Independence}

\begin{defn}{Span}{span_def}
The \textbf{span} of vectors $v_1, \ldots, v_n \in V$ is
\[\text{span}(v_1, \ldots, v_n) = \left\{\sum_{i=1}^n \alpha_i v_i : \alpha_1, \ldots, \alpha_n \in \mathbb{F}\right\}\]
This is the set of all linear combinations.
\end{defn}

\begin{lem}{Span is a Subspace}{span_subspace}
Any set of vectors $v_1, \ldots, v_n$: $\text{span}(v_1, \ldots, v_n)$ is a subspace of $V$.
\end{lem}

\begin{defn}{Linear Independence}{lin_indep}
Vectors $v_1, \ldots, v_n \in V$ are \textbf{linearly independent} if the only linear combination $c_1 v_1 + c_2 v_2 + \cdots + c_n v_n = \vec{0}$ is the ``trivial'' one where $c_1 = \cdots = c_n = 0$.

Conversely, we call $v_1, \ldots, v_n$ \textbf{linearly dependent} if $\exists c_1, \ldots, c_n \in \mathbb{F}$, at least one nonzero, s.t. $c_1 v_1 + \cdots + c_n v_n = \vec{0}$.
\end{defn}

\begin{lem}{Dependence Means Redundancy}{dep_redundancy}
$v_1, \ldots, v_n \in V$ are linearly dependent iff $\exists k \in \{1, \ldots, n\}$ such that $v_k = c_1 v_1 + \cdots + c_{k-1} v_{k-1}$ for some $c_1, \ldots, c_{k-1} \in \mathbb{F}$.
\end{lem}

\begin{proof}
$(\Leftarrow)$: Rearrange $v_k = c_1 v_1 + \cdots + c_{k-1} v_{k-1}$ to get $c_1 v_1 + \cdots + c_{k-1} v_{k-1} + (-1) v_k + 0 \cdot v_{k+1} + \cdots + 0 \cdot v_n = \vec{0}$.

$(\Rightarrow)$: $c_1 v_1 + \cdots + c_n v_n = \vec{0}$ with not all zero. Consider the maximum $k$ s.t. $c_k \neq 0$. Then $c_k v_k = -c_1 v_1 - \cdots - c_{k-1} v_{k-1}$, so $v_k = -\frac{c_1}{c_k} v_1 - \cdots - \frac{c_{k-1}}{c_k} v_{k-1}$.
\end{proof}

\begin{exercises}
\begin{enumerate}
    \item \textbf{(1.1)} Using only the field axioms (F1--F6), prove: (i) $a > 0 \Rightarrow -a < 0$, (ii) $a > 0 \Rightarrow \frac{1}{a} > 0$, (iii) $a > b > 0 \Rightarrow \frac{1}{a} < \frac{1}{b}$.
    
    \item \textbf{(1.2)} If $S \subset \mathbb{R}$ is nonempty and bounded above, prove there exists a sequence $\{a_n\} \subset S$ with $\lim a_n = \sup S$.
    
    \item \textbf{(1.3)} Prove that every positive real number has a positive square root. \textit{Hint: Consider $S = \{x : x^2 < a, x > 0\}$.}
    
    \item \textbf{(1.4)} Let $\mathbb{Z}_2 = \{0, 1\}$ with $1 + 1 = 0$. Show that $\mathbb{Z}_2$ is a field but cannot be ordered.
    
    \item \textbf{(1.5)} Use the Archimedean property to prove that $\lim \frac{1}{n} = 0$.
    
    \item \textbf{(1.6)} Decide whether each is a subspace of $\mathbb{R}^3$: (a) $\{(x_1,x_2,x_3) : x_1 + x_2 + x_3 = 0\}$, (b) $\{(x_1,x_2,x_3) : x_1 + x_2 + x_3 = 1\}$.
    
    \item \textbf{(1.7)} If $v_1, v_2, v_3, v_4$ are linearly independent, are $v_1 - v_2, v_2 - v_3, v_3 - v_4, v_4 - v_1$ linearly independent?
    
    \item \textbf{(1.8)} Prove that for subspaces $U_1, U_2 \subset V$, $U_1 \cup U_2$ is a subspace iff $U_1 \subseteq U_2$ or $U_2 \subseteq U_1$.
\end{enumerate}
\end{exercises}

% ==================== WEEK 2 ====================
\chapter{Subsequences and Continuity}

\section{Analysis: Mon/Tue}

\subsection{Subsequences and Bolzano-Weierstrass}

\begin{thm}{Convergent $\implies$ bounded}{convergent->bounded}
    If $\{a_n\}_{n=1}^\infty\subseteq\mathbb{R}$ is a convergent sequence then $\{a_n\}_{n=1}^\infty$ is bounded. 
\end{thm}
\begin{proof}
    We WTS there exists $c>0$ s.t. $|a_n|<c$ for all $n\geq{N}$. Using the definition of convergence with $\varepsilon=1$, we know that there exists $L\in\mathbb{R}$ and $N\in{\mathbb{N}}$ s.t. $|a_n-L|<1$ for all $n>N$. Hence, $a_n\leq\max\{|L|+1,|a_1|,...,|a_N|\}=c$. 
\end{proof}
\begin{remark}
    The concept of the subsequence is very important in real analysis. Before we get to them, consider the sequence $-1,1,-1,1,-1,...$ that is bounded but not convergent. Observe that if we remove a lot of terms, the resulting sequence is convergent. (It becomes $1,1,1,...$ or $-1,-1,-1,...$)
\end{remark}


\begin{defn}{Subsequence}{subseq}
If $\{a_n\}^\infty_{n=1}$ is a subsequence, $\{a_{f(k)}\}^\infty_{k=1}$ if a subsequence of $a_n$ if $f:\mathbb{N}\to\mathbb{N}$ is strictly increasing ($f(k)<f(k+1))$. 
\end{defn}
\begin{remark}
    Common notation for a subsequence is $\{a_{n(k)}\}^\infty_{n=1}$ or $\{a_{n_k}\}_{k=1}^\infty$
\end{remark}

\begin{thm}{Bolzano-Weierstrass}{bw_thm}
Any bounded sequence in $\mathbb{R}$ has a convergent subsequence.
\end{thm}

\begin{proof}
Let $\{a_n\}$ be bounded, say $a_n \in [c, d]$ for all $n$.

\textbf{Step 1: Construct nested intervals.}

Let $I_0 = [c, d]$. Divide $I_0$ into two halves. At least one half contains infinitely many terms of $\{a_n\}$. Call it $I_1 = [c_1, d_1]$.

Continue: given $I_k = [c_k, d_k]$ containing infinitely many terms, divide in half and let $I_{k+1}$ be the half containing infinitely many terms. Note:
\begin{itemize}
    \item $I_0 \supseteq I_1 \supseteq I_2 \supseteq \cdots$ (nested)
    \item $|I_k| = d_k - c_k = \frac{d - c}{2^k} \to 0$
\end{itemize}

\textbf{Step 2: Find the limit.}

Since $\{c_k\}$ is increasing and bounded above (by $d$), it converges to some $L$. Similarly $\{d_k\}$ is decreasing and bounded below, converging to some $L'$.

But $d_k - c_k \to 0$, so $L = L'$. Thus $\bigcap_{k=0}^\infty I_k = \{L\}$.

\textbf{Step 3: Construct the convergent subsequence.}

Choose $n_1$ such that $a_{n_1} \in I_1$.

Given $n_k$ with $a_{n_k} \in I_k$, choose $n_{k+1} > n_k$ such that $a_{n_{k+1}} \in I_{k+1}$. (Possible since $I_{k+1}$ has infinitely many terms.)

\textbf{Step 4: Verify convergence.}

For any $\varepsilon > 0$, choose $K$ such that $|I_K| < \varepsilon$. Then for all $k \geq K$:
\[|a_{n_k} - L| \leq |I_k| \leq |I_K| < \varepsilon\]
since $a_{n_k} \in I_k$ and $L \in I_k$.

Hence $\lim_{k \to \infty} a_{n_k} = L$.
\end{proof}

\begin{cor}{Sequential Compactness}{seq_compact}
A subset $K \subseteq \mathbb{R}$ is compact (closed and bounded) iff every sequence in $K$ has a subsequence converging to a point in $K$.
\end{cor}

\subsection{Limit Theorems}

\begin{thm}{Algebraic Limit Theorems}{limit_algebra}
Let $\lim_{n \to \infty} a_n = A$ and $\lim_{n \to \infty} b_n = B$. Then:
\begin{enumerate}
    \item $\lim(a_n + b_n) = A + B$
    \item $\lim(a_n \cdot b_n) = A \cdot B$
    \item $\lim(a_n / b_n) = A/B$ (if $B \neq 0$ and $b_n \neq 0$ for all $n$)
\end{enumerate}
\end{thm}

\begin{proof}
We prove (1); (2) and (3) are similar but more technical.

Let $\varepsilon > 0$. By definition of limits:
\begin{itemize}
    \item $\exists N_1$ s.t. $|a_n - A| < \varepsilon/2$ for $n \geq N_1$
    \item $\exists N_2$ s.t. $|b_n - B| < \varepsilon/2$ for $n \geq N_2$
\end{itemize}

Let $N = \max(N_1, N_2)$. For $n \geq N$:
\[|(a_n + b_n) - (A + B)| \leq |a_n - A| + |b_n - B| < \frac{\varepsilon}{2} + \frac{\varepsilon}{2} = \varepsilon\]
\end{proof}

\begin{lem}{Limits Preserve Inequalities}{limit_ineq}
If $a_n \leq b_n$ for all $n$ and both limits exist, then $\lim a_n \leq \lim b_n$.
\end{lem}

\begin{proof}
Suppose $A = \lim a_n > B = \lim b_n$. Let $\varepsilon = (A - B)/2 > 0$.

Then $\exists N$ s.t. $|a_n - A| < \varepsilon$ and $|b_n - B| < \varepsilon$ for $n \geq N$.

This gives $a_n > A - \varepsilon = (A + B)/2$ and $b_n < B + \varepsilon = (A + B)/2$.

So $a_n > b_n$, contradicting $a_n \leq b_n$.
\end{proof}

\begin{thm}{Squeeze Theorem}{squeeze}
If $a_n \leq b_n \leq c_n$ and $\lim a_n = \lim c_n = L$, then $\lim b_n = L$.
\end{thm}

\begin{proof}
Let $\varepsilon > 0$. $\exists N$ s.t. for $n \geq N$:
\[L - \varepsilon < a_n \leq b_n \leq c_n < L + \varepsilon\]
Hence $|b_n - L| < \varepsilon$ for $n \geq N$.
\end{proof}

\subsection{Continuity}
\begin{defn}{Continuity}{continuous_func}
$f: [a,b] \to \mathbb{R}$ is \textbf{continuous} at $c$ if:
\[ \lim_{x \to c} f(x) = f(c) \]
Equivalently (Sequential Criterion), for any sequence $x_n \to c$, $f(x_n) \to f(c)$.
\end{defn}

\begin{thm}{Extreme Value Theorem (EVT)}{evt}
If $f$ is continuous on a closed interval $[a,b]$, then $f$ is bounded and achieves its maximum and minimum.
\end{thm}

\begin{proof}
\textbf{Part 1: $f$ is bounded above.}

Suppose not. Then for each $n$, $\exists x_n \in [a,b]$ with $f(x_n) > n$.

By Bolzano-Weierstrass, $\{x_n\}$ has a convergent subsequence $x_{n_k} \to c \in [a,b]$.

By continuity, $f(x_{n_k}) \to f(c)$. But $f(x_{n_k}) > n_k \to \infty$. Contradiction.

\textbf{Part 2: $f$ achieves its maximum.}

Let $M = \sup\{f(x) : x \in [a,b]\}$. (Exists since $f$ is bounded.)

By definition of sup, $\exists x_n \in [a,b]$ with $f(x_n) > M - 1/n$.

By Bolzano-Weierstrass, $\{x_n\}$ has a convergent subsequence $x_{n_k} \to c \in [a,b]$.

By continuity, $f(x_{n_k}) \to f(c)$.

Since $f(x_{n_k}) > M - 1/n_k \to M$, we have $f(c) \geq M$.

Since $f(c) \leq M$ by definition of $M$, we conclude $f(c) = M$.

The minimum case is analogous.
\end{proof}

\begin{thm}{Intermediate Value Theorem (IVT)}{ivt}
If $f$ is continuous on $[a,b]$ and $A$ lies between $f(a)$ and $f(b)$, then there exists $c \in [a,b]$ such that $f(c) = A$.
\end{thm}

\begin{proof}
WLOG assume $f(a) < A < f(b)$. (Otherwise replace $f$ with $-f$.)

Define $S = \{x \in [a,b] : f(x) < A\}$.

\textbf{Claim:} $c = \sup S$ satisfies $f(c) = A$.

$S$ is nonempty (contains $a$) and bounded above (by $b$), so $c = \sup S$ exists.

\textbf{Step 1:} $f(c) \geq A$.

By definition of sup, $\exists x_n \in S$ with $x_n \to c$.

Since $f(x_n) < A$ and $f$ is continuous: $f(c) = \lim f(x_n) \leq A$.

\textbf{Step 2:} $f(c) \leq A$.

Suppose $f(c) > A$. By continuity, $\exists \delta > 0$ s.t. $f(x) > A$ for $x \in (c - \delta, c + \delta) \cap [a,b]$.

But this means no $x \in (c - \delta, c]$ can be in $S$, contradicting $c = \sup S$.

Combining: $f(c) = A$.
\end{proof}

\begin{cor}{Roots of Polynomials}{poly_roots}
Any odd-degree polynomial has at least one real root.
\end{cor}

\begin{proof}
Let $p(x) = x^{2n+1} + a_{2n}x^{2n} + \cdots + a_0$.

As $x \to +\infty$, $p(x) \to +\infty$. As $x \to -\infty$, $p(x) \to -\infty$.

So $\exists a < 0 < b$ with $p(a) < 0 < p(b)$. By IVT, $\exists c$ with $p(c) = 0$.
\end{proof}

\section{Linear Algebra: Thu/Fri}

\subsection{Basis and Dimension}

\begin{lem}{Replacement Lemma}{replacement}
If $W = \text{span}(v_1, \ldots, v_n)$ and $w_1, \ldots, w_k \in W$ are linearly independent, then we can replace $k$ of the $v_i$'s with the $w_i$'s while preserving $\text{span}(\cdots) = W$.
\end{lem}

\begin{proof}[Sketch]
We start with $(v_1, \ldots, v_n)$ and replace them one by one with $w_1, \ldots, w_k$, preserving that $\text{span}(\cdots) = W$.

Since $w_1 \in W$, we have $w_1 = c_1 v_1 + \cdots + c_n v_n$. There is some $v_j$ which is a linear combination of the others (including $w_1$), so we can replace $v_j$ with $w_1$ and keep the same span. Repeat.
\end{proof}

\begin{lem}{Bound on Independent Vectors}{bound_indep}
If $V = \text{span}(v_1, \ldots, v_n)$ and $w_1, \ldots, w_k \in V$ are linearly independent, then $k \leq n$.
\end{lem}

\begin{proof}
Assume not. By the replacement lemma, $V = \text{span}(w_1, \ldots, w_n)$. Then $w_{n+1} \in V = \text{span}(w_1, \ldots, w_n)$, so $\{w_1, \ldots, w_n, w_{n+1}\}$ is linearly dependent, contradicting independence of $w_1, \ldots, w_k$.
\end{proof}

\begin{defn}{Basis}{basis_def}
$\{v_1, \ldots, v_n\} \subseteq V$ is called a \textbf{basis} of $V$ if:
\begin{enumerate}
    \item[(B1)] $\text{span}(v_1, \ldots, v_n) = V$
    \item[(B2)] $\{v_1, \ldots, v_n\}$ is linearly independent
\end{enumerate}
\end{defn}

\begin{example}
$\{e_1, e_2, e_3\}$ is a basis for $\mathbb{R}^3$, where $e_1 = (1,0,0)$, $e_2 = (0,1,0)$, $e_3 = (0,0,1)$.
\end{example}

\begin{lem}{Unique Coordinates}{unique_coords}
If $\{v_1, \ldots, v_n\}$ is a basis, then any $v \in V$ can be written \textbf{uniquely} as $v = \sum_{i=1}^n \alpha_i v_i$.
\end{lem}

\begin{proof}
That any $v$ can be written this way follows from (B1). For uniqueness, suppose $v = \sum \alpha_i v_i = \sum \beta_i v_i$. Then $\sum(\alpha_i - \beta_i) v_i = \vec{0}$. By (B2), $\alpha_i - \beta_i = 0$ for all $i$, so $\alpha_i = \beta_i$.
\end{proof}

\begin{defn}{Dimension}{dimension_def}
Any two bases of a vector space $V$ have the same number of elements. We call this number the \textbf{dimension} of $V$, written $\dim V$.
\end{defn}

\begin{proof}[Proof that dimension is well-defined]
Suppose $\{v_1, \ldots, v_n\}$ and $\{w_1, \ldots, w_m\}$ are both bases. Whichever is larger would give linearly independent vectors in the span of the smaller, contradicting the bound lemma.
\end{proof}

\subsection{Inner Products}

\begin{defn}{Dot Product}{dot_product}
For $v, w \in \mathbb{R}^n$ with $v = (v_1, \ldots, v_n)$ and $w = (w_1, \ldots, w_n)$, define
\[v \cdot w = \sum_{i=1}^n v_i w_i\]
We say $v$ is \textbf{orthogonal} to $w$ (write $v \perp w$) if $v \cdot w = 0$.
\end{defn}

\begin{thm}{Properties of Dot Product}{dot_props}
\begin{enumerate}
    \item $v \cdot w = w \cdot v$
    \item $(\alpha v + \beta w) \cdot u = \alpha(v \cdot u) + \beta(w \cdot u)$
    \item $v \cdot v \geq 0$, with equality iff $v = \vec{0}$
\end{enumerate}
\end{thm}

\begin{defn}{Norm}{norm_def}
Given $v \in \mathbb{R}^n$, define $\|v\| = \sqrt{v \cdot v}$.
\end{defn}

\begin{thm}{Pythagorean Theorem}{pythag}
If $v \cdot w = 0$, then $\|v\|^2 + \|w\|^2 = \|v + w\|^2$.
\end{thm}

\begin{proof}
$\|v + w\|^2 = (v + w) \cdot (v + w) = v \cdot v + 2(v \cdot w) + w \cdot w = \|v\|^2 + \|w\|^2$.
\end{proof}

\begin{thm}{Cauchy-Schwarz Inequality}{cauchy_schwarz}
$|v \cdot w| \leq \|v\| \|w\|$ for all $v, w \in \mathbb{R}^n$.
\end{thm}

\begin{proof}
If $w = 0$, both sides are 0. Assume $w \neq 0$.

Decompose $v = v_\parallel + v_\perp$ where $v_\parallel = \alpha w$ for some $\alpha$ and $v_\perp \cdot w = 0$.

From $v \cdot w = (v_\parallel + v_\perp) \cdot w = \alpha w \cdot w = \alpha \|w\|^2$, we get $\alpha = \frac{v \cdot w}{\|w\|^2}$.

By Pythagorean theorem: $\|v\|^2 = \|v_\parallel\|^2 + \|v_\perp\|^2 \geq \|v_\parallel\|^2 = \frac{(v \cdot w)^2}{\|w\|^2}$.

Hence $(v \cdot w)^2 \leq \|v\|^2 \|w\|^2$.
\end{proof}

\begin{thm}{Triangle Inequality}{triangle_ineq}
$\|v + w\| \leq \|v\| + \|w\|$.
\end{thm}

\begin{exercises}
\begin{enumerate}
    \item \textbf{(2.1)} (\textit{Sandwich Theorem}) If $\lim a_n = \lim b_n = L$ and $a_n \leq c_n \leq b_n$ for all $n$, prove $\lim c_n = L$.
    
    \item \textbf{(2.2)} (\textit{Heine-Borel}) Let $\mathcal{F}$ be a collection of open intervals covering $[a,b]$. Prove there is a finite subcollection that still covers $[a,b]$.
    
    \item \textbf{(2.3)} If $\lim f(x_n) = f(c)$ for every sequence $\{x_n\} \to c$, prove $f$ is continuous at $c$.
    
    \item \textbf{(2.4)} Let $f(x) = 0$ if $x$ is irrational, and $f(p/q) = 1/q$ for $p/q$ in lowest terms. Prove $f$ is continuous at every irrational.
    
    \item \textbf{(2.5)} Let $V$ be finite-dimensional and $U \subseteq V$ a subspace. Prove $U = V$ iff $\dim U = \dim V$.
    
    \item \textbf{(2.6)} If $\dim U + \dim W > \dim V$ for subspaces $U, W \subseteq V$, prove $U \cap W \neq \{0\}$.
    
    \item \textbf{(2.7)} Find a basis for $V = \{p \in \mathcal{P}_4(\mathbb{R}) : p(1) = 0\}$ and extend it to a basis of $\mathcal{P}_4(\mathbb{R})$.
    
    \item \textbf{(2.8)} If $\{v_1, \ldots, v_m\}$ is linearly independent and $w \in V$, prove $\dim(\text{span}\{v_1+w, \ldots, v_m+w\}) \geq m-1$.
\end{enumerate}
\end{exercises}


% ==================== WEEK 3 ====================
\chapter{Differentiation and Topology}

\section{Analysis: Mon/Tue}

\subsection{Differentiable Functions in $\mathbb{R}$}

\begin{defn}{Differentiability}{diff_R}
Given $f:[a,b]\to\mathbb{R}$, $c\in[a,b]$, we say $f$ is \textbf{differentiable} at $c$ if 
\[\lim_{h\to0}\frac{f(c+h)-f(c)}{h}\]
exists. We denote the limit, if it exists, by $f'(c)$ or $\frac{df}{dx}(c)$.
\end{defn}

\begin{remark}
Some basic derivative laws:
\begin{enumerate}
    \item If $f,g$ are differentiable at $c$, then $\alpha f + \beta g$ is differentiable at $c$ with $(\alpha f + \beta g)'(c) = \alpha f'(c) + \beta g'(c)$.
    \item If $f:[a,b]\to\mathbb{R}$ is constant, then $f'(c) = 0$ for all $c \in (a,b)$.
    \item If $f(x) = x$, then $f'(c) = 1$ for all $c \in (a,b)$.
\end{enumerate}
\end{remark}

\begin{thm}{Mean Value Theorem}{mvt}
Let $f:[a,b]\to\mathbb{R}$ be continuous on $[a,b]$ and differentiable on $(a,b)$. Then there exists $c \in (a,b)$ such that
\[f'(c) = \frac{f(b) - f(a)}{b - a}\]
\end{thm}

\begin{thm}{Rolle's Theorem}{rolle}
Let $f:[a,b]\to\mathbb{R}$ be continuous on $[a,b]$ and differentiable on $(a,b)$. Suppose $f(a) = 0 = f(b)$. Then there exists $c \in (a,b)$ such that $f'(c) = 0$.
\end{thm}

\begin{proof}
We first prove Rolle's theorem and use it to prove MVT.

If $f$ is constant, then $f'(c) = 0$ and we are done. Otherwise, either $\exists x_0 \in (a,b)$ with $f(x_0) > 0$ or $\exists x_0 \in (a,b)$ with $f(x_0) < 0$.

In the former case, there exists $c \in [a,b]$ such that $f(x) \leq f(c)$ for all $x \in [a,b]$. Since $f(c) \geq f(x_0) > 0$, we know $c \in (a,b)$.

\textbf{Claim:} $f'(c) = 0$.

We know $f'(c) = \lim_{x \to c} \frac{f(x) - f(c)}{x - c}$.

Take $\{x_n\}$ with $\lim x_n = c$ and $c < x_n < b$. Then 
\[f'(c) = \lim_{n\to\infty} \frac{f(x_n) - f(c)}{x_n - c} \leq 0\]
since the numerator is $\leq 0$ and denominator is $> 0$.

On the other hand, take $\{x_n\}$ with $\lim x_n = c$ and $a < x_n < c$. Then
\[f'(c) = \lim_{n\to\infty} \frac{f(x_n) - f(c)}{x_n - c} \geq 0\]
since both numerator and denominator are $\leq 0$.

Hence $f'(c) = 0$.

\textbf{MVT from Rolle's:} Given $f$, define $g:[a,b]\to\mathbb{R}$ by
\[g(x) = f(x) - f(a) - \frac{f(b) - f(a)}{b - a}(x - a)\]

Note that $g(a) = 0$ and $g(b) = 0$. By Rolle's theorem, there exists $c \in (a,b)$ such that
\[0 = g'(c) = f'(c) - \frac{f(b) - f(a)}{b - a}\]
which gives $f'(c) = \frac{f(b) - f(a)}{b - a}$.
\end{proof}

\subsection{Sequences in $\mathbb{R}^n$}

\begin{defn}{Sequence in $\mathbb{R}^n$}{seq_Rn}
A sequence in $\mathbb{R}^n$ is a map $\mathbb{N} \to \mathbb{R}^n$. We denote by $\{a^{(k)}\}_{k=1}^\infty$, where $a^{(k)} = (a_1^{(k)}, \ldots, a_n^{(k)})$.
\end{defn}

\begin{defn}{Limit in $\mathbb{R}^n$}{limit_Rn}
Given $\{a^{(k)}\}_{k=1}^\infty \subseteq \mathbb{R}^n$ and $a \in \mathbb{R}^n$, we say $\lim_{k\to\infty} a^{(k)} = a$ if for all $\varepsilon > 0$, there exists $\mathcal{K} \in \mathbb{N}$ such that $\|a^{(k)} - a\| < \varepsilon$ whenever $k \geq \mathcal{K}$.

Here $\|b\| = \sqrt{b \cdot b} = \sqrt{\sum_{i=1}^n b_i^2}$.
\end{defn}

\begin{lem}{Component-wise Convergence}{componentwise}
$\{a^{(k)}\}_{k=1}^\infty \subseteq \mathbb{R}^n$, $a \in \mathbb{R}^n$. Then $\lim a^{(k)} = a$ if and only if $\lim a_i^{(k)} = a_i$ for all $i = 1, \ldots, n$.
\end{lem}

\begin{proof}
$(\Leftarrow)$: Let $\varepsilon > 0$. Since $\lim a_i^{(k)} = a_i$ for every $i$, there exists $\mathcal{K}_i \in \mathbb{N}$ such that $|a_i^{(k)} - a_i| < \frac{\varepsilon}{\sqrt{n}}$ whenever $k \geq \mathcal{K}_i$.

Let $\mathcal{K} = \max\{\mathcal{K}_1, \ldots, \mathcal{K}_n\}$. Then
\[\|a^{(k)} - a\| = \left(\sum_{i=1}^n |a_i^{(k)} - a_i|^2\right)^{1/2} < \left(\sum_{i=1}^n \frac{\varepsilon^2}{n}\right)^{1/2} = \varepsilon\]

$(\Rightarrow)$: Exercise.
\end{proof}

\begin{defn}{Bounded Set}{bounded_Rn}
$S \subseteq \mathbb{R}^n$ is \textbf{bounded} if there exists $c > 0$ such that $\|x\| < c$ for all $x \in S$.
\end{defn}

\begin{thm}{Bolzano-Weierstrass in $\mathbb{R}^n$}{bw_Rn}
Any bounded sequence in $\mathbb{R}^n$ has a convergent subsequence.
\end{thm}

\begin{proof}[Sketch]
Prove by induction on $n$. The base case was proven in Week 1. Assume the result up to dimension $n-1$.

Let $\{a^{(k)}\}_{k=1}^\infty \subseteq \mathbb{R}^n$. Write $a^{(k)} = (a^{1(k)}, a_n^{(k)})$ where $a^{1(k)} = (a_1^{(k)}, \ldots, a_{n-1}^{(k)})$.

By induction, there exists a subsequence $\{a^{(k_j)}\}$ such that $\lim_{j\to\infty} a^{1(k_j)}$ exists.

By BW in $\mathbb{R}$, there exists a further subsequence $\{a_n^{(k_{j_\ell})}\}$ such that $\lim_{\ell\to\infty} a_n^{(k_{j_\ell})}$ exists.

Then $\{a^{(k_{j_\ell})}\}_{\ell=1}^\infty \subseteq \mathbb{R}^n$ is convergent.
\end{proof}

\subsection{Open and Closed Sets}

\begin{defn}{Ball}{ball_def}
For $x \in \mathbb{R}^n$ and $\rho > 0$, define $B_\rho(x) = \{y \in \mathbb{R}^n : \|x - y\| < \rho\}$.
\end{defn}

\begin{defn}{Open Set}{open_def}
$U \subseteq \mathbb{R}^n$ is \textbf{open} if for all $x \in U$, there exists $\rho > 0$ such that $B_\rho(x) \subseteq U$.
\end{defn}

\begin{defn}{Limit Point}{limit_point}
$A \subseteq \mathbb{R}^n$. We say $b \in \mathbb{R}^n$ is a \textbf{limit point} of $A$ if there exists $\{a^{(k)}\}_{k=1}^\infty$ with $a^{(k)} \in A$ such that $b = \lim_{k\to\infty} a^{(k)}$.
\end{defn}

\begin{defn}{Closed Set}{closed_def}
$K \subseteq \mathbb{R}^n$ is \textbf{closed} if it contains all its limit points.
\end{defn}

\begin{remark}
A set can be both open and closed (e.g., $\mathbb{R}^n$ and $\varnothing$), or neither (e.g., $(a,b]$).
\end{remark}

\begin{thm}{Union of Open Sets}{union_open}
If $\{U_\alpha\}_{\alpha \in \Gamma}$ is a collection of open sets in $\mathbb{R}^n$, then $\bigcup_{\alpha \in \Gamma} U_\alpha$ is open.

If $\{K_\alpha\}_{\alpha \in \Gamma}$ is a collection of closed sets, then $\bigcap_{\alpha \in \Gamma} K_\alpha$ is closed.
\end{thm}

\begin{proof}
Let $x \in \bigcup_{\alpha \in \Gamma} U_\alpha$. Then there exists $\alpha_0 \in \Gamma$ such that $x \in U_{\alpha_0}$.

Since $U_{\alpha_0}$ is open, there exists $\rho > 0$ such that $B_\rho(x) \subseteq U_{\alpha_0} \subseteq \bigcup_{\alpha \in \Gamma} U_\alpha$.
\end{proof}

\begin{thm}{Open-Closed Duality}{open_closed_dual}
$U \subseteq \mathbb{R}^n$ is open if and only if $\mathbb{R}^n \setminus U$ is closed.
\end{thm}

\begin{proof}
$(\Rightarrow)$: Let $U$ be open and let $z \in \mathbb{R}^n$ be a limit point of $\mathbb{R}^n \setminus U$. We want to show $z \notin U$.

If $z \in U$, there exists $\rho > 0$ such that $B_\rho(z) \subseteq U$. But since $z$ is a limit point of $\mathbb{R}^n \setminus U$, there exists $y \in \mathbb{R}^n \setminus U$ with $\|y - z\| < \rho$. Hence $y \in B_\rho(z) \subseteq U$ and $y \in \mathbb{R}^n \setminus U$, contradiction.

$(\Leftarrow)$: Prove the contrapositive. If $U$ is not open, then $\exists x \in U$ such that for all $\rho > 0$, there exists $y_\rho \in B_\rho(x)$ with $y_\rho \in \mathbb{R}^n \setminus U$. Taking $\rho = 1/j$, we get a sequence converging to $x$, so $x$ is a limit point of $\mathbb{R}^n \setminus U$ but $x \in U$, hence $\mathbb{R}^n \setminus U$ is not closed.
\end{proof}

\begin{defn}{Compact Set}{compact_def}
$K \subseteq \mathbb{R}^n$ is \textbf{compact} if for all sequences $\{a^{(k)}\}_{k=1}^\infty \subseteq K$, there exists a subsequence $\{a^{(k_j)}\}$ and $a \in K$ such that $\lim_{j\to\infty} a^{(k_j)} = a$.
\end{defn}

\begin{remark}
In $\mathbb{R}^n$: closed + bounded $\Leftrightarrow$ compact.
\end{remark}

\section{Linear Algebra: Thu/Fri}

\subsection{Linear Maps}

\begin{defn}{Linear Map}{linear_map}
A function $T: V \to W$ (over field $\mathbb{F}$) is a \textbf{linear map} if
\[T(\alpha x + \beta y) = \alpha T(x) + \beta T(y)\]
for all $\alpha, \beta \in \mathbb{F}$ and $x, y \in V$.
\end{defn}

\begin{remark}
$T(0) = 0$ always!
\end{remark}

\begin{example}
\begin{enumerate}
    \item $T: \mathbb{R}^2 \to \mathbb{R}^2$ given by $T(x_1, x_2) = (x_1, x_1 x_2)$ is NOT linear.
    \item $T(x_1, x_2) = (x_1, \lambda x_2)$ for fixed $\lambda \in \mathbb{R}$ IS linear.
    \item Rotation: $(x_1, x_2) \mapsto (x_1 \cos\theta - x_2 \sin\theta, x_1 \sin\theta + x_2 \cos\theta)$ is linear.
\end{enumerate}
\end{example}

\subsection{Matrix Representation}

Fix a basis $\{v_1, \ldots, v_n\}$ of $V$ and $\{w_1, \ldots, w_m\}$ of $W$.

How does $T$ act on the basis vectors? Write
\[T(v_j) = a_{1j} w_1 + a_{2j} w_2 + \cdots + a_{mj} w_m\]

For any vector $x \in V$ with $x = \xi_1 v_1 + \cdots + \xi_n v_n$:
\begin{align*}
T(x) &= \xi_1 T(v_1) + \cdots + \xi_n T(v_n) \\
&= (a_{11}\xi_1 + \cdots + a_{1n}\xi_n) w_1 + \cdots + (a_{m1}\xi_1 + \cdots + a_{mn}\xi_n) w_m
\end{align*}

\begin{defn}{Matrix of a Linear Map}{matrix_of_map}
Given a linear map $T: V \to W$ with $T(v_j) = \sum_i a_{ij} w_i$, the \textbf{matrix of $T$} under these bases is
\[A = (a_{ij}) \in \mathbb{F}^{m \times n}\]
with $n$ columns ($= \dim V$) and $m$ rows ($= \dim W$).
\end{defn}

\begin{defn}{Linear Operator}{linear_operator}
A linear map from $V$ to $V$ is called a \textbf{linear operator} on $V$. Its matrix is square ($n \times n$).
\end{defn}

\subsection{Composition and Matrix Multiplication}

Given $T_1: V_1 \to V_2$ and $T_2: V_2 \to V_3$, define $T_2 \circ T_1: V_1 \to V_3$ by $(T_2 \circ T_1)(x) = T_2(T_1(x))$.

If $A$ is the matrix for $T_1$ and $B$ is the matrix for $T_2$, then the matrix for $T_2 \circ T_1$ is $C = B \cdot A$ where
\[c_{ik} = \sum_j b_{ij} a_{jk} = (\text{$i$th row of $B$}) \cdot (\text{$k$th column of $A$})\]

\begin{remark}
Matrix multiplication is associative: $A(BC) = (AB)C$, but NOT commutative: $AB \neq BA$ in general.
\end{remark}

\begin{exercises}
\begin{enumerate}
    \item \textbf{(3.1)} Prove that a continuous function $f: [a,b] \to \mathbb{R}$ is uniformly continuous.
    
    \item \textbf{(3.2)} Show that $f(x) = 1/x$ on $(0,1)$ is \textit{not} uniformly continuous.
    
    \item \textbf{(3.3)} (\textit{Cauchy criterion}) Prove that $\{a_n\}$ converges iff $\forall \varepsilon > 0$, $\exists N$ s.t. $|a_m - a_n| < \varepsilon$ for $m, n \geq N$.
    
    \item \textbf{(3.4)} Prove: $E \subseteq A$ is relatively open in $A$ iff $\exists$ open $U \subseteq \mathbb{R}^n$ with $E = A \cap U$.
    
    \item \textbf{(3.5)} Prove the identity $\frac{1}{2}\sum_{i,j}(a_i b_j - a_j b_i)^2 = \|a\|^2\|b\|^2 - (a \cdot b)^2$ and use it to prove Cauchy-Schwarz.
    
    \item \textbf{(3.6)} Prove that equality holds in the triangle inequality iff $y = \lambda x$ for some $\lambda \geq 0$ (or one is zero).
    
    \item \textbf{(3.7)} Prove $(AB)^T = B^T A^T$ directly from the definition of matrix multiplication.
    
    \item \textbf{(3.8)} Find the matrix of rotation by angle $\theta$ in $\mathbb{R}^2$. What is the matrix for reflection across $y = x$?
\end{enumerate}
\end{exercises}

% ==================== WEEK 4 ====================
\chapter{Continuity and Metric Spaces}

\section{Linear Algebra: Thu/Fri (continued from Week 3)}

\subsection{Kernel and Range}

\begin{defn}{Kernel and Range}{ker_range}
For a linear map $T: V \to W$:
\begin{enumerate}
    \item $\ker(T) = \{v \in V : T(v) = 0\}$
    \item $\text{Range}(T) = \{y \in W : \exists v \in V, T(v) = y\}$
\end{enumerate}
\end{defn}

\begin{lem}{Kernel and Range are Subspaces}{ker_range_subspace}
$\ker(T)$ and $\text{Range}(T)$ are subspaces of $V$ and $W$, respectively.
\end{lem}

\begin{proof}
Take $u, v \in \ker(T)$ and $\alpha, \beta \in \mathbb{F}$:
\[T(\alpha u + \beta v) = \alpha T(u) + \beta T(v) = 0 + 0 = 0\]

Take $x, y \in \text{Range}(T)$ with $x = T(u)$, $y = T(v)$:
\[\alpha x + \beta y = \alpha T(u) + \beta T(v) = T(\alpha u + \beta v) \in \text{Range}(T)\]
\end{proof}

\begin{remark}
For a matrix $A \in \mathbb{F}^{m \times n}$:
\begin{itemize}
    \item $\text{Range}(A) = \text{span}(\text{columns of } A) = $ column space of $A$
    \item $\ker(A) = \{x \in \mathbb{F}^n : Ax = 0\}$
    \item Row space: $R(A) = \text{span}(\text{rows of } A)$
\end{itemize}
\end{remark}

\begin{defn}{Injective, Surjective, Bijective}{inj_surj_bij}
\begin{itemize}
    \item $f$ is \textbf{injective}: if $u \neq v$, then $f(u) \neq f(v)$
    \item $f$ is \textbf{surjective}: $\forall y \in W$, $\exists u \in V$ s.t. $f(u) = y$
    \item $f$ is \textbf{bijective}: both injective and surjective
\end{itemize}
\end{defn}

\begin{remark}
For linear maps:
\begin{itemize}
    \item $T$ is surjective $\Leftrightarrow$ $\text{Range}(T) = W$
    \item $T$ is injective $\Leftrightarrow$ $\ker(T) = \{0\}$
\end{itemize}
\end{remark}

\subsection{Rank-Nullity Theorem}

\begin{thm}{Rank-Nullity Theorem}{rank_nullity}
If $T: V \to W$ is linear and $V$ is finite dimensional, then
\[\dim(\ker(T)) + \dim(\text{Range}(T)) = \dim V\]
\end{thm}

\begin{proof}
Construct a basis $\{u_1, \ldots, u_k\}$ of $\ker(T)$ (assuming $\dim(\ker(T)) = k$; if $k = 0$, this is empty).

Extend to a basis of $V$: $\{u_1, \ldots, u_k, v_{k+1}, \ldots, v_n\}$.

\textbf{Claim:} $\text{Range}(T) = \text{span}(T(v_{k+1}), \ldots, T(v_n))$.

If $y \in \text{Range}(T)$, then $y = T(x)$ for some $x = \alpha_1 u_1 + \cdots + \alpha_k u_k + \alpha_{k+1} v_{k+1} + \cdots + \alpha_n v_n$, so $y = \alpha_{k+1} T(v_{k+1}) + \cdots + \alpha_n T(v_n)$.

\textbf{Claim:} $T(v_{k+1}), \ldots, T(v_n)$ are linearly independent.

Assume $\beta_{k+1} T(v_{k+1}) + \cdots + \beta_n T(v_n) = 0$. Then $T(\beta_{k+1} v_{k+1} + \cdots + \beta_n v_n) = 0$, so $\beta_{k+1} v_{k+1} + \cdots + \beta_n v_n \in \ker(T) = \text{span}(u_1, \ldots, u_k)$.

By linear independence of the full basis, $\beta_{k+1} = \cdots = \beta_n = 0$.

Hence $\dim(\text{Range}(T)) = n - k$.
\end{proof}

\begin{defn}{Rank}{rank_def}
Define the \textbf{rank} of $A$ as $\text{rank}(A) = \dim(\text{Range}(A)) = \dim(\text{column space})$.
\end{defn}

\begin{lem}{Rank from RREF}{rank_rref}
$\dim(\ker(A)) = n - \#\text{pivots}$ and $\text{rank}(A) = \#\text{pivots}$.
\end{lem}

\section{Analysis: Mon/Tue}

\subsection{Continuity in Multiple Dimensions}

\begin{defn}{Continuity in $\mathbb{R}^n$}{cont_Rn}
Given $A \subseteq \mathbb{R}^n$, $f: A \to \mathbb{R}^m$, $a \in A$, we say $f$ is \textbf{continuous at $a$} if for all $\varepsilon > 0$, there exists $\delta > 0$ such that $\|f(x) - f(a)\| < \varepsilon$ whenever $\|x - a\| < \delta$ and $x \in A$.

$f$ is \textbf{continuous} if it is continuous at all $a \in A$.
\end{defn}

\begin{remark}
Some useful facts:
\begin{enumerate}
    \item $a \in A$ is an \textbf{isolated point} if there exists $\rho > 0$ such that $B_\rho(a) \cap A = \{a\}$.
    \item $f$ is continuous at a non-isolated point $a$ iff $\lim_{x \to a} f(x) = f(a)$.
    \item Equivalently: $f$ is continuous at $a$ iff for any sequence $\{a^{(k)}\} \subseteq A$ with $\lim a^{(k)} = a$, we have $\lim f(a^{(k)}) = f(a)$.
\end{enumerate}
\end{remark}

\begin{thm}{Extreme Value Theorem (Compact Version)}{evt_compact}
If $A \subseteq \mathbb{R}^n$ is compact and $f: A \to \mathbb{R}$ is continuous, then $f$ is bounded and achieves its maximum and minimum.
\end{thm}

\begin{proof}
Suppose not. Then for all $k \in \mathbb{N}$, there exists $x^{(k)} \in A$ such that $|f(x^{(k)})| \geq k$.

Since $A$ is compact, there exists a subsequence $\{x^{(k_j)}\}$ and $a \in A$ such that $\lim_{j\to\infty} x^{(k_j)} = a$.

By continuity, $\lim_{j\to\infty} f(x^{(k_j)}) = f(a)$, which is convergent and thus bounded. But $|f(x^{(k_j)})| \geq k_j \to \infty$, contradiction.
\end{proof}

\subsection{Metric Spaces}

\begin{defn}{Metric Space}{metric_space}
$(X, d)$ is a \textbf{metric space} if $X$ is a set and $d: X \times X \to \mathbb{R}$ satisfies:
\begin{enumerate}
    \item[(M1)] $d(x, y) \geq 0$ for all $x, y \in X$, with equality iff $x = y$.
    \item[(M2)] $d(x, y) = d(y, x)$ for all $x, y \in X$.
    \item[(M3)] $d(x, z) \leq d(x, y) + d(y, z)$ for all $x, y, z \in X$ (triangle inequality).
\end{enumerate}
\end{defn}

\begin{example}
\begin{enumerate}
    \item $X = \mathbb{R}^n$, $d(x, y) = \|x - y\| = \sqrt{\sum_{i=1}^n (x_i - y_i)^2}$
    \item $X \subseteq \mathbb{R}^n$, $d(x, y) = \|x - y\|$
    \item Discrete metric: $d(x, y) = \begin{cases} 1 & x \neq y \\ 0 & x = y \end{cases}$
    \item Any normed vector space $(V, \|\cdot\|)$ with $d(x, y) = \|x - y\|$
\end{enumerate}
\end{example}

\begin{remark}
In a metric space $(X, d)$:
\begin{enumerate}
    \item $\{x_n\} \subseteq X$ converges to $x \in X$ if $\forall \varepsilon > 0$, $\exists N$ such that $d(x_n, x) < \varepsilon$ for $n \geq N$.
    \item Ball: $B_\rho(y) = \{x \in X : d(x, y) < \rho\}$
    \item $x$ is a limit point of $A$ if $\exists \{a_n\} \subseteq A$ with $\lim a_n = x$.
    \item $A$ is open if $\forall x \in A$, $\exists \rho > 0$ with $B_\rho(x) \subseteq A$.
    \item $A$ is closed if it contains all its limit points.
    \item $A$ is bounded if $\exists y \in X$, $\rho > 0$ with $A \subseteq B_\rho(y)$.
    \item $A$ is compact if every sequence in $A$ has a convergent subsequence with limit in $A$.
\end{enumerate}
\end{remark}

\begin{thm}{Compact implies Closed and Bounded}{compact_closed_bdd}
Let $(X, d)$ be a metric space and $A \subseteq X$. If $A$ is compact, then $A$ is closed and bounded.
\end{thm}

\begin{proof}
\textbf{Boundedness:} Suppose not. Then for all $y \in X$ and $\rho > 0$, $A \not\subseteq B_\rho(y)$. So there exists $a_n \in A$ with $d(a_n, y) \geq n$. By compactness, some subsequence $\{a_{n_k}\}$ converges to $a \in A$, so $\{a_{n_k}\}$ is bounded. But $d(a_{n_k}, y) \geq n_k \geq k \to \infty$, contradiction.

\textbf{Closedness:} Suppose $x \in X$ is a limit point of $A$. We WTS $x \in A$. By definition, there exists $\{a_n\} \subseteq A$ with $\lim a_n = x$. By compactness, there exists $\{a_{n_k}\}$ with $\lim a_{n_k} = a \in A$. Since limits are unique, $x = a \in A$.
\end{proof}

\begin{defn}{Continuity in Metric Spaces}{cont_metric}
Given metric spaces $(X, d_X)$, $(Y, d_Y)$ and $f: X \to Y$:
\begin{enumerate}
    \item $f$ is \textbf{continuous at $c \in X$} if $\forall \varepsilon > 0$, $\exists \delta > 0$ such that $d_Y(f(x), f(c)) < \varepsilon$ whenever $d_X(x, c) < \delta$.
    \item $f$ is \textbf{continuous} if it is continuous at all points.
\end{enumerate}
\end{defn}

\begin{thm}{Image of Compact is Compact}{compact_image}
Let $(X, d_X)$, $(Y, d_Y)$ be metric spaces with $X$ compact and $f: X \to Y$ continuous. Then $f(X) = \{f(x) : x \in X\}$ is compact.
\end{thm}

\begin{proof}
Let $\{y_n\} \subseteq f(X)$. Then there exists $x_n \in X$ with $y_n = f(x_n)$.

By compactness of $X$, there exists $\{x_{n_k}\}$ and $x \in X$ with $\lim x_{n_k} = x$.

By continuity of $f$, $\lim y_{n_k} = \lim f(x_{n_k}) = f(x) \in f(X)$.
\end{proof}

\section{Linear Algebra: Thu/Fri (continued)}

See Week 3 for kernel, range, and rank-nullity theorem content.

\begin{exercises}
\begin{enumerate}
    \item \textbf{(4.1)} Let $f: [a,b] \to \mathbb{R}^n$ be differentiable. Prove $\exists c \in (a,b)$ with $\|f(b) - f(a)\| \leq (b-a)\|f'(c)\|$.
    
    \item \textbf{(4.2)} Prove that $\bar{A}$ (closure of $A$) is closed in any metric space.
    
    \item \textbf{(4.3)} Prove that $f: X \to Y$ continuous and $U \subseteq Y$ open implies $f^{-1}(U)$ is open in $X$.
    
    \item \textbf{(4.4)} If $A$ is dense in $X$ and $f, g: X \to Y$ are continuous with $f|_A = g|_A$, prove $f = g$.
    
    \item \textbf{(4.5)} Prove that $\text{rank}(AB) \leq \min(\text{rank}(A), \text{rank}(B))$.
    
    \item \textbf{(4.6)} Prove that $\dim \ker(T \circ S) \leq \dim \ker T + \dim \ker S$.
    
    \item \textbf{(4.7)} Find the matrix of orthogonal projection onto $W = \text{span}\{v\}$ where $v$ is a unit vector in $\mathbb{R}^n$.
    
    \item \textbf{(4.8)} Solve the system with augmented matrix $\begin{pmatrix} 1 & 0 & 1 & 1 & 1 & | & 1 \\ 2 & 1 & 1 & 3 & 1 & | & 1 \\ 1 & 1 & 2 & 0 & 2 & | & 2 \\ 0 & 0 & 1 & -1 & 1 & | & 1 \end{pmatrix}$.
\end{enumerate}
\end{exercises}

% ==================== WEEK 5 ====================
\chapter{Series and Differentiation in Higher Dimensions}

\section{Linear Algebra: Thu/Fri (Gaussian Elimination)}

\subsection{Solving Linear Systems}

\begin{lem}{Solutions to $Ax = 0$}{ax_zero_solutions}
\begin{enumerate}
    \item Solutions to $Ax = 0$ = $\ker(A)$
    \item $Ax = b$ has a solution $\Leftrightarrow$ $b \in \text{Range}(A)$
\end{enumerate}
Moreover, if $\bar{x}$ is a particular solution, then the set of all solutions is $\{\bar{x} + w : w \in \ker(A)\} = \bar{x} + \ker(A)$.
\end{lem}

\begin{proof}
Suppose $x = \bar{x} + w$ with $w \in \ker(A)$. Then $Ax = A(\bar{x} + w) = A\bar{x} + Aw = b + 0 = b$.

Conversely, if $x$ is a solution, then $A(x - \bar{x}) = b - b = 0$, so $x - \bar{x} \in \ker(A)$.
\end{proof}

\subsection{Row Operations}

\begin{defn}{Row Operations}{row_ops}
\begin{enumerate}
    \item[(R1)] Switch two rows
    \item[(R2)] Multiply row $i$ by nonzero constant $\lambda \in \mathbb{F}$
    \item[(R3)] Replace row $i$ by (row $i$ + $\lambda \cdot$ row $j$)
\end{enumerate}
\end{defn}

\begin{lem}{Row Operations Preserve Kernel}{row_ops_kernel}
\begin{enumerate}
    \item Performing row operations on $A$ doesn't change $\ker(A)$ (it may change $\text{Range}(A)$ though!)
    \item Performing row operations on the augmented matrix $(A \mid b)$ does not change the solution set.
\end{enumerate}
\end{lem}

\begin{lem}{RREF Existence}{rref_exists}
Every (augmented) matrix can be reduced to RREF by finitely many row operations.
\end{lem}

\begin{example}
To find $\ker(A)$: reduce to RREF. Non-pivot variables are ``free.'' E.g., if RREF gives
\[\begin{pmatrix} 1 & 0 & 3 \\ 0 & 1 & -3 \end{pmatrix}\]
then $\ker(A) = \{(-3\beta, 3\beta, \beta) : \beta \in \mathbb{F}\}$.
\end{example}

\subsection{Orthogonal Complements}

\begin{defn}{Orthogonal Complement}{orth_comp}
For $V \subseteq \mathbb{F}^n$ a subspace,
\[V^\perp = \{x \in \mathbb{F}^n : x \cdot v = 0 \text{ for all } v \in V\}\]
\end{defn}

\begin{thm}{Dimension of Orthogonal Complement}{dim_orth_comp}
\begin{enumerate}
    \item $\dim V + \dim V^\perp = n$
    \item $(V^\perp)^\perp = V$
    \item If $\mathbb{F} = \mathbb{R}$, then $V \cap V^\perp = \{0\}$
    \item If $\mathbb{F} = \mathbb{R}$, then $\forall x \in \mathbb{R}^n$, $\exists!$ decomposition $x = v + v^\perp$ with $v \in V$, $v^\perp \in V^\perp$
\end{enumerate}
\end{thm}

\begin{proof}[Proof of (1)]
Let $\{v_1, \ldots, v_k\}$ be a basis of $V$. Form matrix $A$ with rows $v_i$. Note $\ker(A) = V^\perp$ since $Ax = 0$ iff $v_i \cdot x = 0$ for all $i$.

Hence $\dim(\text{Range}(A)) + \dim(V^\perp) = n$, and $\dim(\text{Range}(A)) = k = \dim V$.
\end{proof}

\section{Analysis: Mon/Tue}

\subsection{Series}

\begin{defn}{Series}{series_def}
Let $\{a_k\}_{k=1}^\infty \subseteq \mathbb{R}^n$. We say $\sum_{k=1}^\infty a_k = s \in \mathbb{R}^n$ if the sequence of partial sums $S_j = \sum_{k=1}^j a_k$ satisfies $\lim_{j\to\infty} S_j = s$.

We say $\sum a_k$ \textbf{converges} if such $s$ exists, and \textbf{diverges} otherwise.
\end{defn}

\begin{thm}{Terms Must Vanish}{terms_vanish}
Suppose $\{a_k\}_{k=1}^\infty \subseteq \mathbb{R}^n$ and $\sum a_k$ converges. Then $\lim_{k\to\infty} a_k = 0$.
\end{thm}

\begin{proof}
Let $\varepsilon > 0$. We WTS $\|a_k\| < \varepsilon$ for $k \geq \mathcal{K}$.

By assumption, there exists $S \in \mathbb{R}^n$ with $\sum a_k = S$. So $\|S_j - S\| < \frac{\varepsilon}{2}$ whenever $j \geq \mathcal{J}$.

Let $\mathcal{K} = \mathcal{J} + 1$. Then for $k \geq \mathcal{K}$:
\[\|a_k\| = \|S_k - S_{k-1}\| \leq \|S - S_k\| + \|S - S_{k-1}\| < \frac{\varepsilon}{2} + \frac{\varepsilon}{2} = \varepsilon\]
\end{proof}

\begin{thm}{Non-negative Series}{nonneg_series}
Suppose $\{a_k\}_{k=1}^\infty \subseteq \mathbb{R}$ with $a_k \geq 0$ for all $k$. Then $\sum_{k=1}^\infty a_k$ converges iff $\{S_j\}$ is bounded above.
\end{thm}

\begin{proof}
Observe that $\{S_j\}$ is increasing since $a_k \geq 0$. Recall that an increasing sequence converges iff it is bounded above.
\end{proof}

\begin{example}
\textbf{Geometric Series:} $\sum_{k=0}^\infty r^k$ converges if $r \in (0,1)$ and diverges if $r \geq 1$.

\begin{proof}
For $r \neq 1$: $S_j = 1 + r + \cdots + r^j$ and $rS_j = r + \cdots + r^{j+1}$.

So $(r-1)S_j = r^{j+1} - 1$, giving $S_j = \frac{r^{j+1} - 1}{r - 1}$.

This is bounded (converges to $\frac{1}{1-r}$) if $r \in (0,1)$ and unbounded if $r > 1$.
\end{proof}
\end{example}

\begin{example}
\textbf{Harmonic Series:} $\sum_{k=1}^\infty \frac{1}{k}$ diverges.

\begin{proof}
\[S_{2^j} = 1 + \frac{1}{2} + \left(\frac{1}{3} + \frac{1}{4}\right) + \left(\frac{1}{5} + \cdots + \frac{1}{8}\right) + \cdots\]
\[\geq 1 + \frac{1}{2} + \frac{2}{4} + \frac{4}{8} + \cdots = 1 + \frac{j}{2} \to \infty\]
\end{proof}
\end{example}

\begin{example}
\textbf{$p$-series:} $\sum_{k=1}^\infty \frac{1}{k^p}$ converges if $p \geq 2$ (in fact, if $p > 1$).

\begin{proof}
\[\sum_{k=1}^{2^{j-1}} \frac{1}{k^p} \leq 1 + \left(\frac{1}{2^p} + \frac{1}{2^p}\right) + \left(\frac{4}{4^p}\right) + \cdots \leq 1 + \left(\frac{1}{2}\right)^{p-1} + \left(\frac{1}{2}\right)^{2(p-1)} + \cdots\]
Since $p > 1$, this is a convergent geometric series.
\end{proof}
\end{example}

\begin{defn}{Absolute Convergence}{abs_conv}
$\sum_{k=1}^\infty a_k$ is \textbf{absolutely convergent} if $\sum_{k=1}^\infty |a_k|$ converges.
\end{defn}

\begin{thm}{Absolute Convergence Implies Convergence}{abs_implies_conv}
If $\sum a_k$ is absolutely convergent, then it is convergent.
\end{thm}

\begin{proof}
Let $p_k = \max\{a_k, 0\}$ and $q_k = \max\{-a_k, 0\}$. Then $|a_k| = p_k + q_k$ and $a_k = p_k - q_k$.

Since $\sum |a_k|$ converges, $\sum_{k=1}^j |a_k|$ is bounded. Then $\sum_{k=1}^j p_k \leq \sum_{k=1}^j |a_k|$ is bounded, so $\sum p_k$ converges. Similarly $\sum q_k$ converges.

Hence $\sum a_k = \sum(p_k - q_k)$ converges.
\end{proof}

\begin{defn}{Rearrangement}{rearrangement}
Given $\{a_k\}_{k=1}^\infty \subseteq \mathbb{R}$, we say $\sum_{k=1}^\infty a_{j(k)}$ is a \textbf{rearrangement} if $j: \mathbb{N} \to \mathbb{N}$ is a bijection.
\end{defn}

\begin{thm}{Rearrangements of Absolutely Convergent Series}{rearrange_abs}
If $\sum a_k$ converges absolutely, then any rearrangement $\sum a_{j(k)}$ also converges to the same sum.
\end{thm}

\begin{proof}
First consider the case $a_k \geq 0$. Suppose $\sum a_k = S$.

For every $\ell \in \mathbb{N}$, define $p(\ell) = \max\{j(1), \ldots, j(\ell)\}$.

Then $\sum_{k=1}^\ell a_{j(k)} \leq \sum_{k=1}^{p(\ell)} a_k \leq S$.

So $\sum a_{j(k)}$ converges. Since the partial sums are increasing,
\[\sum_{k=1}^\infty a_{j(k)} = \sup_\ell \sum_{k=1}^\ell a_{j(k)} \leq S\]

By symmetry (the original is a rearrangement of the rearrangement), we get equality.

For the general case, use $a_k = p_k - q_k$ where both $\sum p_k$ and $\sum q_k$ converge by the previous theorem.
\end{proof}

\subsection{Differentiation in Higher Dimensions}

\begin{defn}{Derivative Matrix}{deriv_matrix}
Given $f: U \to \mathbb{R}^m$ where $U \subseteq \mathbb{R}^n$ is open, and $c \in U$. We say $f$ is \textbf{differentiable at $c$} if there exists $A \in \mathbb{R}^{m \times n}$ such that
\[\lim_{h \to 0} \frac{\|f(c+h) - f(c) - Ah\|}{\|h\|} = 0\]

If such $A$ exists, it is called the \textbf{derivative matrix} of $f$ at $c$, denoted $Df(c)$.
\end{defn}

\begin{lem}{Matrix Norm Inequality}{matrix_norm}
If $\|A\| = \left(\sum_{i,j} a_{ij}^2\right)^{1/2}$, then $\|Ax\| \leq \|A\| \|x\|$.
\end{lem}

\begin{proof}
Write $A = (a_{ij})$. Then
\[\|Ax\|^2 = \sum_{i=1}^m (a_i \cdot x)^2 \leq \sum_{i=1}^m \|a_i\|^2 \|x\|^2 = \|A\|^2 \|x\|^2\]
by Cauchy-Schwarz. Take square roots.
\end{proof}

\begin{thm}{Differentiable Implies Continuous}{diff_implies_cont}
If $f: U \to \mathbb{R}^m$ is differentiable at $c$, then $f$ is continuous at $c$.
\end{thm}

\begin{proof}
Since $U$ is open, there exists $\rho > 0$ with $B_\rho(c) \subseteq U$.

Using the definition with $\varepsilon = 1$, there exists $\delta_0 \in (0, \rho)$ such that
\[\frac{\|f(c+h) - f(c) - Ah\|}{\|h\|} < 1 \text{ whenever } 0 < \|h\| < \delta_0\]

Let $\varepsilon > 0$. We WTS there exists $\delta$ such that $\|f(c+h) - f(c)\| < \varepsilon$ whenever $\|h\| < \delta$.

Let $\delta = \min\{\delta_0, \frac{\varepsilon}{2}, \frac{\varepsilon}{2\|A\|}\}$ (assuming $A \neq 0$).

Then:
\[\|f(c+h) - f(c)\| \leq \|f(c+h) - f(c) - Ah\| + \|Ah\| < \|h\| + \|A\|\|h\| \leq \frac{\varepsilon}{2} + \frac{\varepsilon}{2} = \varepsilon\]
\end{proof}

\section{Linear Algebra: Thu/Fri}
\begin{center}
    \textit{Content from Linear Algebra PDF needed}
\end{center}


% ==================== WEEK 6 ====================
\chapter{Directional Derivatives and the Chain Rule}

\section{Linear Algebra: Thu/Fri (Determinants)}

\subsection{Permutations and Signs}

\begin{defn}{Permutation}{permutation}
A \textbf{permutation} $\sigma: \{1, \ldots, n\} \to \{1, \ldots, n\}$ is a bijection. The set of all permutations on $n$ elements is denoted $S_n$.
\end{defn}

\begin{defn}{Transposition}{transposition}
A \textbf{transposition} $\tau_{kl}$ is a permutation that swaps $k$ and $l$ and fixes everything else.
\end{defn}

\begin{lem}{Permutations as Transpositions}{perm_as_trans}
Every permutation can be written as a composition of adjacent transpositions $\tau_{k,k+1}$.
\end{lem}

\begin{defn}{Inversions and Sign}{inversions_sign}
For $\sigma \in S_n$, the \textbf{number of inversions} is
\[N(\sigma) = |\{(i,j) : i < j \text{ but } \sigma(i) > \sigma(j)\}|\]
Define $\text{sgn}(\sigma) = (-1)^{N(\sigma)}$.
\end{defn}

\begin{remark}
$\text{sgn}(\tau_{kl}) = -1$ for any transposition. Each adjacent swap changes the sign.
\end{remark}

\subsection{The Determinant}

\begin{defn}{Determinant}{determinant}
For $A = (a_{ij}) \in \mathbb{F}^{n \times n}$:
\[\det(A) = \sum_{\sigma \in S_n} \text{sgn}(\sigma) \cdot a_{1,\sigma(1)} a_{2,\sigma(2)} \cdots a_{n,\sigma(n)}\]
\end{defn}

\begin{remark}
The meaning of the determinant is a ``signed volume'' --- the sign captures the orientation of the vectors.
\end{remark}

\begin{thm}{Properties of Determinant}{det_properties}
\begin{enumerate}
    \item $\det(A^T) = \det(A)$
    \item Switching two rows/columns multiplies $\det$ by $-1$
    \item $\det$ is linear in each row/column
    \item If two rows/columns are equal, $\det = 0$
    \item $\det(AB) = \det(A) \cdot \det(B)$
\end{enumerate}
\end{thm}

\subsection{Inverses}

\begin{defn}{Inverse Matrix}{inverse_def}
$B \in \mathbb{F}^{n \times n}$ is an \textbf{inverse} of $A$ if $AB = I = BA$.
\end{defn}

\begin{lem}{Uniqueness of Inverse}{inverse_unique}
An inverse, if it exists, is unique.
\end{lem}

\begin{proof}
Suppose $B, C = A^{-1}$. Then $B = BI = B(AC) = (BA)C = IC = C$.
\end{proof}

\begin{lem}{One-Sided Implies Two-Sided}{one_sided_inverse}
If $A, B \in \mathbb{F}^{n \times n}$ and $AB = I$, then $BA = I$.
\end{lem}

\begin{proof}
$AB = I$ implies $\text{Range}(AB) = \mathbb{F}^n$, so $B$ is surjective. By rank-nullity, $\ker(B) = \{0\}$.

Now $AB = I$ implies $ABy = y$ for all $y$. Hence $BABy = By$ for all $y$. Since $B$ is surjective, $BAx = x$ for all $x$, so $BA = I$.
\end{proof}

\begin{thm}{Invertibility Criterion}{invert_criterion}
$A$ is invertible iff $\det(A) \neq 0$.
\end{thm}

\begin{proof}
$(\Rightarrow)$: $AA^{-1} = I$ implies $\det(A)\det(A^{-1}) = 1$, so $\det(A) \neq 0$.

$(\Leftarrow)$: See adjugate formula: $A \cdot \text{adj}(A) = \det(A) \cdot I$.
\end{proof}

\begin{thm}{Equivalent Statements for Invertibility}{equiv_invert}
The following are equivalent:
\begin{enumerate}
    \item $A$ is invertible
    \item $\exists B$ s.t. $AB = I$
    \item $\exists B$ s.t. $BA = I$
    \item $\det(A) \neq 0$
    \item $\text{rref}(A) = I$
    \item $\ker(A) = \{0\}$
    \item $\text{Range}(A) = \mathbb{F}^n$
    \item $\text{rank}(A) = n$
\end{enumerate}
\end{thm}

\section{Analysis: Mon/Tue}

\subsection{Directional and Partial Derivatives}

\begin{defn}{Directional Derivative}{dir_deriv}
Let $f: U \to \mathbb{R}^m$ where $U \subseteq \mathbb{R}^n$ is open, $a \in U$, $v \in \mathbb{R}^n$. The \textbf{directional derivative} of $f$ at $a$ in direction $v$ is
\[(D_v f)(a) = \lim_{t \to 0} \frac{f(a + tv) - f(a)}{t}\]
if this limit exists.
\end{defn}

\begin{remark}
This is the derivative of the single-variable function $g(t) = f(a + tv)$ at $t = 0$.

When $v = e_j = (0, \ldots, 1, \ldots, 0)$, the directional derivative is called a \textbf{partial derivative}, denoted $(D_j f)(a)$ or $\frac{\partial f}{\partial x_j}(a)$.
\end{remark}

\begin{thm}{Directional Derivative via Derivative Matrix}{dir_deriv_matrix}
Suppose $f$ is differentiable at $a \in U$. Then for any $v \in \mathbb{R}^n$, the directional derivative $(D_v f)(a)$ exists and
\[(D_v f)(a) = [Df(a)] v\]
\end{thm}

\begin{proof}
Let $\varepsilon > 0$. We WTS $\left\|\frac{f(a+tv) - f(a)}{t} - Df(a) \cdot v\right\| < \varepsilon$ if $0 < |t| < \delta$.

If $v = 0$, this is obvious. Assume $v \neq 0$.

By definition of $Df(a)$, there exists $\delta' > 0$ such that
\[\frac{\|f(a+h) - f(a) - Df(a) \cdot h\|}{\|h\|} < \frac{\varepsilon}{\|v\|} \text{ whenever } \|h\| < \delta'\]

Let $\delta = \frac{\delta'}{\|v\|}$. If $0 < |t| < \delta$, then $h = tv$ satisfies $0 < \|h\| = |t|\|v\| < \delta'$.

Hence:
\[\frac{\|f(a+tv) - f(a) - Df(a) \cdot tv\|}{|t| \|v\|} < \frac{\varepsilon}{\|v\|}\]
which gives
\[\left\|\frac{f(a+tv) - f(a)}{t} - Df(a) \cdot v\right\| < \varepsilon\]
\end{proof}

\begin{thm}{Corollary: Partial Derivatives}{partial_deriv_cor}
If $f$ is differentiable at $a$:
\begin{enumerate}
    \item $(D_j f)(a) = Df(a) \cdot e_j =$ the $j$-th column of $Df(a)$.
    \item For $v = \sum_{j=1}^n v_j e_j$: $(D_v f)(a) = \sum_{j=1}^n v_j (D_j f)(a)$.
\end{enumerate}
\end{thm}

\begin{thm}{Continuous Partials Imply Differentiability}{cont_partials}
Given $U \subseteq \mathbb{R}^n$ open, $f: U \to \mathbb{R}^m$. Suppose there exists $a \in U$ and $\rho > 0$ such that $D_j f$ exists everywhere in $B_\rho(a)$ for all $j$, and $D_j f$ is continuous at $a$. Then $f$ is differentiable at $a$.
\end{thm}

\begin{proof}
Let $\varepsilon > 0$. We WTS 
\[\frac{\left\|f(a+h) - f(a) - \sum_{j=1}^n (D_j f)(a) h_j\right\|}{\|h\|} < \varepsilon\]
whenever $0 < \|h\| < \delta$, where $h = \sum_{j=1}^n h_j e_j$.

Define $h^{(k)} = \sum_{j=1}^k h_j e_j$ for $k = 1, \ldots, n$.

Then:
\[f(a+h) - f(a) = \sum_{k=1}^n \left[f(a + h^{(k)}) - f(a + h^{(k-1)})\right]\]

By MVT applied to each term:
\[f(a + h^{(k)}) - f(a + h^{(k-1)}) = h_k (D_k f)(a + h^{(k-1)} + \theta_k h_k e_k)\]
for some $\theta_k \in (0, 1)$.

Hence:
\[\frac{\left\|f(a+h) - f(a) - \sum_{j=1}^n (D_j f)(a) h_j\right\|}{\|h\|}\]
\[= \frac{\left\|\sum_{j=1}^n h_j \left[(D_j f)(a + h^{(j-1)} + \theta_j h_j e_j) - (D_j f)(a)\right]\right\|}{\|h\|}\]
\[\leq \sum_{j=1}^n \left\|(D_j f)(a + h^{(j-1)} + \theta_j h_j e_j) - (D_j f)(a)\right\|\]

Since $D_j f$ is continuous at $a$, there exists $\delta_j$ such that each term is $< \frac{\varepsilon}{n}$ when $\|h\| < \delta_j$.

Take $\delta = \min\{\delta_1, \ldots, \delta_n\}$.
\end{proof}

\subsection{The Chain Rule}

\begin{thm}{Chain Rule}{chain_rule}
Let $U \subseteq \mathbb{R}^n$, $V \subseteq \mathbb{R}^m$ be open. Let $f: U \to V$ and $g: V \to \mathbb{R}^\ell$.

Suppose $f$ is differentiable at $a \in U$ and $g$ is differentiable at $f(a) \in V$. Then $g \circ f: U \to \mathbb{R}^\ell$ is differentiable at $a$, and
\[\underbrace{D(g \circ f)(a)}_{\mathbb{R}^{\ell \times n}} = \underbrace{Dg(f(a))}_{\mathbb{R}^{\ell \times m}} \cdot \underbrace{Df(a)}_{\mathbb{R}^{m \times n}}\]
\end{thm}

\begin{proof}
Let $\varepsilon > 0$. We WTS there exists $\delta > 0$ such that $B_\delta(a) \subseteq U$ and
\[\|(g \circ f)(a+h) - (g \circ f)(a) - Dg(f(a)) Df(a) h\| < \varepsilon \|h\|\]
whenever $\|h\| < \delta$.

Let $\varepsilon_1, \varepsilon_2 > 0$ to be chosen. By differentiability of $f$ at $a$, there exists $\delta_1 > 0$ such that
\[\|f(a+h) - f(a) - Df(a) h\| < \varepsilon_1 \|h\| \text{ for } \|h\| < \delta_1\]

By differentiability of $g$ at $f(a)$, there exists $\delta_2 > 0$ such that
\[\|g(y) - g(f(a)) - Dg(f(a))(y - f(a))\| < \varepsilon_2 \|y - f(a)\| \text{ for } \|y - f(a)\| < \delta_2\]

Since $f$ is continuous at $a$, there exists $\delta_3 > 0$ such that $\|f(a+h) - f(a)\| < \delta_2$ for $\|h\| < \delta_3$.

For $\|h\| < \min\{\delta_1, \delta_3\}$:
\begin{align*}
&\|(g \circ f)(a+h) - (g \circ f)(a) - Dg(f(a)) Df(a) h\| \\
&\leq \|g(f(a+h)) - g(f(a)) - Dg(f(a))(f(a+h) - f(a))\| \\
&\quad + \|Dg(f(a))\| \|f(a+h) - f(a) - Df(a) h\| \\
&< \varepsilon_2 \|f(a+h) - f(a)\| + \|Dg(f(a))\| \varepsilon_1 \|h\| \\
&\leq (\varepsilon_1 \varepsilon_2 + \varepsilon_2 \|Df(a)\| + \varepsilon_1 \|Dg(f(a))\|) \|h\|
\end{align*}

Choose $\varepsilon_1, \varepsilon_2$ small enough that this is $< \varepsilon \|h\|$.
\end{proof}

\begin{thm}{Chain Rule: Component Form}{chain_component}
With $f, g$ as above:
\[D_i(g \circ f)(a) = \sum_{k=1}^m (D_k g_i)(f(a)) \cdot (D_i f_k)(a)\]
\end{thm}

\begin{thm}{One-dimensional Chain Rule}{chain_1d}
If $n = m = 1$: $(g \circ f)'(a) = g'(f(a)) \cdot f'(a)$.
\end{thm}

\section{Linear Algebra: Thu/Fri}
\begin{center}
    \textit{Content from Linear Algebra PDF needed}
\end{center}


% ==================== WEEK 7 ====================
\chapter{Higher Derivatives and Taylor's Theorem}

\section{Analysis: Mon/Tue}

\subsection{Higher Derivatives}

Given $f: U \to \mathbb{R}^m$ where $U \subseteq \mathbb{R}^n$ is open, we can take partial derivatives $D_i f$.

\begin{defn}{Higher Partial Derivatives}{higher_partials}
Suppose $D_i f$ is defined everywhere on $U$. Then we can think of $D_i f: U \to \mathbb{R}^m$. Define
\[D_j D_i f(a) = D_j(D_i f)(a)\]
if it exists. Define higher derivatives inductively by
\[D_{i_1} D_{i_2} \cdots D_{i_k} f = D_{i_1}(D_{i_2} \cdots D_{i_k} f)\]
if $D_{i_2} \cdots D_{i_k} f: U \to \mathbb{R}^m$ is a well-defined function and this exists.
\end{defn}

\begin{defn}{$C^k$ Functions}{ck_functions}
We say $f: U \to \mathbb{R}^m$ is $C^k$ if $f$ and all higher derivatives up to $k$-th order are well-defined and continuous.
\end{defn}

\begin{thm}{Symmetry of Mixed Partials}{mixed_partials}
Let $f: U \to \mathbb{R}^m$, $U \subseteq \mathbb{R}^n$ open. Suppose $a \in U$, $\rho > 0$ s.t. $B_\rho(a) \subseteq U$ and $D_i f$, $D_j f$, $D_i D_j f$, $D_j D_i f$ are defined in $B_\rho(a)$.

Assume, moreover, that $D_i D_j f$ and $D_j D_i f$ are continuous at $a$. Then
\[D_i D_j f(a) = D_j D_i f(a)\]
\end{thm}

\begin{proof}[Sketch]
Consider $\Delta = f(a + he_i + ke_j) - f(a + he_i) - f(a + ke_j) + f(a)$.

Define $\psi_k(x) = f(x + ke_j) - f(x)$. Then $\Delta = \psi_k(a + he_i) - \psi_k(a)$.

By MVT: $\Delta = D_i \psi_k(a + \theta h e_i) \cdot h = [D_i f(a + he_i + \theta' k e_j) - D_i f(a + \theta' k e_j)] \cdot h$

$= D_j D_i f(a + \theta'' h e_i + \theta' k e_j) \cdot hk$ for some $\theta, \theta', \theta'' \in (0,1)$.

Similarly, defining $\phi_h(x) = f(x + he_i) - f(x)$, we get $\Delta = D_i D_j f(a + \theta''' h e_i + \theta'''' k e_j) \cdot hk$.

Taking $h, k \to 0$ and using continuity gives the result.
\end{proof}

\begin{cor}{Order Independence for $C^k$}{order_indep}
Suppose $f: U \to \mathbb{R}^m$ is $C^k$, $k \geq 2$. Then
\[D_{i_1} \cdots D_{i_k} f(x) = D_{j_1} \cdots D_{j_k} f(x)\]
whenever $(j_1, \ldots, j_k)$ is a permutation of $(i_1, \ldots, i_k)$.
\end{cor}

\subsection{Taylor's Theorem}

\begin{defn}{Gradient and Hessian}{grad_hessian}
\begin{enumerate}
    \item For $f: U \to \mathbb{R}$ differentiable, the \textbf{gradient} is $\nabla f = (D_1 f, \ldots, D_n f)$.
    \item A function $Q: \mathbb{R}^n \to \mathbb{R}$ is called a \textbf{quadratic form} if $Q(x) = \sum_{i,j} h_{ij} x_i x_j$ for some symmetric matrix $(h_{ij})$.
    \item The quadratic form $\sum_{i,j} D_i D_j f(a) \cdot x_i x_j$ is called the \textbf{Hessian} of $f$ at $a$.
    \item $Q$ is \textbf{positive definite} if $Q(x) > 0$ for all $x \neq 0$.
\end{enumerate}
\end{defn}

\begin{thm}{Second-Order Taylor Expansion}{taylor_2nd}
Let $U \subseteq \mathbb{R}^n$ open, $f: U \to \mathbb{R}$ be $C^2$. Then for any $a \in U$:
\[f(x) = f(a) + \nabla f(a) \cdot (x-a) + \frac{1}{2} \sum_{i,j} D_i D_j f(a)(x-a)_i(x-a)_j + o(\|x-a\|^2)\]
where $\lim_{x \to a} \frac{o(\|x-a\|^2)}{\|x-a\|^2} = 0$.
\end{thm}

\subsection{Second Derivative Test}

\begin{lem}{Critical Points}{critical_pts}
Let $U$ open, $f: U \to \mathbb{R}$ differentiable. Suppose $f$ has a local min/max at $a$. Then $\nabla f(a) = 0$.
\end{lem}

\begin{thm}{Second Derivative Test}{2nd_deriv_test}
Let $U$ open, $f: U \to \mathbb{R}$ be $C^2$. Suppose $\nabla f(a) = 0$ and the Hessian of $f$ at $a$ is positive definite. Then $f$ has a strict local minimum at $a$.
\end{thm}

\begin{proof}[Sketch]
\textbf{Lemma:} If $Q$ is a positive definite quadratic form, then $\exists M > 0$ s.t. $Q(x) \geq M\|x\|^2$ for all $x \in \mathbb{R}^n$.

Assuming the lemma, then $\exists \rho$ s.t. for $x \in B_\rho(a)$:
\[f(x) \approx f(a) + \frac{1}{2} \text{Hess}_f(x-a) + o(\|x-a\|^2) \geq f(a) + \frac{M}{2}\|x-a\|^2 - o(\|x-a\|^2)\]

For small enough $\|x-a\|$, this is $> f(a)$.
\end{proof}

\section{Linear Algebra: Thu/Fri (Orthonormal Bases)}

\subsection{Orthonormal Bases}

\begin{defn}{Orthonormal Set}{orthonormal}
Vectors $\{u_1, \ldots, u_k\} \subseteq V$ are \textbf{orthonormal} if
\[\langle u_i, u_j \rangle = \begin{cases} 0 & \text{if } i \neq j \\ 1 & \text{if } i = j \end{cases}\]
An \textbf{orthonormal basis} is a basis that is orthonormal.
\end{defn}

\begin{lem}{Orthonormal Implies Independent}{orth_indep}
If $\{u_1, \ldots, u_k\}$ is orthonormal, then it is linearly independent.
\end{lem}

\begin{proof}
Suppose $\sum_{j=1}^k \alpha_j u_j = 0$. Take dot product with $u_i$:
\[0 = \sum_{j=1}^k \alpha_j \langle u_i, u_j \rangle = \alpha_i\]
Hence $\alpha_i = 0$ for all $i$.
\end{proof}

\subsection{Orthogonal Projection}

\begin{defn}{Orthogonal Projection}{orth_proj}
For $W \subseteq \mathbb{F}^n$ a subspace and $x \in \mathbb{F}^n$, there exists a unique decomposition $x = w + w^\perp$ where $w \in W$ and $w^\perp \in W^\perp$.

Define $P_W: \mathbb{F}^n \to W$ by $P_W(x) = w$. This is the \textbf{orthogonal projection} onto $W$.
\end{defn}

\begin{thm}{Projection Formula}{proj_formula}
Suppose $\{u_1, \ldots, u_k\}$ is an orthonormal basis of $W$. Then
\[P_W(x) = \sum_{j=1}^k \langle x, u_j \rangle u_j\]
\end{thm}

\begin{proof}
By definition, it suffices to show $v = x - \sum_j \langle x, u_j \rangle u_j \in W^\perp$.

Let $\langle v, u_i \rangle$ for any $i$. Since $\{u_1, \ldots, u_k\}$ is a basis, we need $\langle v, u_i \rangle = 0$ for all $i$.

$\langle v, u_i \rangle = \langle x, u_i \rangle - \sum_j \langle x, u_j \rangle \langle u_j, u_i \rangle = \langle x, u_i \rangle - \langle x, u_i \rangle = 0$.
\end{proof}

\subsection{Gram-Schmidt Process}

\begin{thm}{Gram-Schmidt}{gram_schmidt}
Suppose $\{v_1, \ldots, v_k\} \subseteq V$ is linearly independent. Then there exists an orthonormal set $\{u_1, \ldots, u_k\}$ s.t.
\[\text{span}(u_1, \ldots, u_j) = \text{span}(v_1, \ldots, v_j)\]
for any $j = 1, \ldots, k$.
\end{thm}

\begin{proof}
Define $u_1 = \frac{v_1}{\|v_1\|}$ (note $v_1 \neq 0$ since linearly independent). Moreover, $\{u_1\}$ is orthonormal and $\text{span}(u_1) = \text{span}(v_1)$.

Inductively, suppose we have $\{u_1, \ldots, u_j\}$ orthonormal with the right span.

Define $\tilde{u}_{j+1} = v_{j+1} - P_{\text{span}(u_1,\ldots,u_j)}(v_{j+1}) = v_{j+1} - \sum_{i=1}^j \langle v_{j+1}, u_i \rangle u_i$.

Then $\tilde{u}_{j+1} \neq 0$ (otherwise $v_{j+1} \in \text{span}(v_1, \ldots, v_j)$, contradicting independence).

Set $u_{j+1} = \frac{\tilde{u}_{j+1}}{\|\tilde{u}_{j+1}\|}$.
\end{proof}

\subsection{Eigenvalues and the Spectral Theorem}

Now we come to one of the most important results in linear algebra: the Spectral Theorem for symmetric matrices.

\begin{defn}{Characteristic Polynomial}{char_poly}
Let $A$ be an $n \times n$ matrix. The \textbf{characteristic polynomial} of $A$ is
\[p_A(\lambda) = \det(A - \lambda I)\]
This is a degree $n$ polynomial in $\lambda$, with leading coefficient $(-1)^n$.
\end{defn}

\begin{defn}{Eigenvalue and Eigenvector}{eigenvalue}
Let $A$ be an $n \times n$ matrix.
\begin{enumerate}
    \item The \textbf{eigenvalues} of $A$ are the roots of the characteristic polynomial, i.e., values $\lambda$ such that $\det(A - \lambda I) = 0$.
    \item An \textbf{eigenvector} corresponding to eigenvalue $\lambda$ is a nonzero vector $v \in \mathbb{R}^n$ such that $Av = \lambda v$.
\end{enumerate}
\end{defn}

\begin{remark}
By the Fundamental Theorem of Algebra, every $n \times n$ matrix has exactly $n$ eigenvalues (counting multiplicity), though some may be complex. We're mainly interested in the case where eigenvalues are real.
\end{remark}

\begin{lem}{Properties of Eigenvalues}{eigen_properties}
Let $A$ be an $n \times n$ matrix with eigenvalues $\lambda_1, \ldots, \lambda_n$. Then:
\begin{enumerate}
    \item $\det A = \lambda_1 \cdots \lambda_n$ (product of eigenvalues)
    \item For $\lambda \in \mathbb{R}$: $\lambda$ is an eigenvalue $\iff$ $\exists v \neq 0$ with $Av = \lambda v$
\end{enumerate}
\end{lem}

\begin{proof}
For (1): Set $\lambda = 0$ in $\det(A - \lambda I) = (-1)^n(\lambda - \lambda_1) \cdots (\lambda - \lambda_n)$.

For (2): $\det(A - \lambda I) = 0 \iff N(A - \lambda I) \neq \{0\} \iff \exists v \neq 0$ with $(A - \lambda I)v = 0 \iff Av = \lambda v$.
\end{proof}

\begin{lem}{Eigenvalues of Symmetric Matrices are Real}{symmetric_real_eigen}
If $A$ is a real symmetric matrix, then all eigenvalues of $A$ are real.
\end{lem}

\begin{proof}
Suppose $Av = \lambda v$ for some $v \neq 0$ (possibly complex). Taking the conjugate transpose:
\[\bar{v}^T A = \bar{\lambda} \bar{v}^T\]
since $A = A^T$ and $A$ is real. Multiplying the original equation by $\bar{v}^T$:
\[\bar{v}^T A v = \lambda \bar{v}^T v = \lambda \|v\|^2\]
But also $\bar{v}^T A v = \bar{\lambda} \bar{v}^T v = \bar{\lambda} \|v\|^2$.

Since $\|v\|^2 \neq 0$, we get $\lambda = \bar{\lambda}$, so $\lambda$ is real.
\end{proof}

\begin{thm}{Spectral Theorem}{spectral_thm}
Let $A$ be a real symmetric $n \times n$ matrix. Then there exists an orthonormal basis $\{v_1, \ldots, v_n\}$ of $\mathbb{R}^n$ consisting of eigenvectors of $A$.
\end{thm}

\begin{proof}
We proceed by induction on $n$.

\textbf{Base case:} $n = 1$ is trivial --- $A$ is just a scalar and $\mathbb{R}^1$ has the standard basis.

\textbf{Inductive step:} Assume the theorem holds for $(n-1) \times (n-1)$ symmetric matrices.

\textit{Step 1: Find a first eigenvector.}

Consider the quadratic form $Q(x) = x^T A x = \sum_{i,j} a_{ij} x_i x_j$.

Restricted to the unit sphere $S^{n-1}$, $Q$ is continuous on a compact set, so it attains a minimum at some point $v_1 \in S^{n-1}$.

Since $Q(x) = \|x\|^2 Q(\hat{x})$ where $\hat{x} = x/\|x\|$, the function $\|x\|^{-2} Q(x)$ also attains its minimum at $v_1$.

Taking the gradient and setting it to zero (this is a constrained optimization on the sphere, using Lagrange multipliers), we find that
\[Av_1 = \lambda_1 v_1\]
where $\lambda_1 = Q(v_1)$ is the minimum value. So $v_1$ is an eigenvector!

\textit{Step 2: Restrict to the orthogonal complement.}

Let $V = \text{span}\{v_1\}$ and $V^\perp$ be its orthogonal complement. Since $\dim V = 1$, we have $\dim V^\perp = n - 1$.

\textbf{Claim:} $A$ maps $V^\perp$ to itself, i.e., if $w \perp v_1$, then $Aw \perp v_1$.

\textit{Proof of claim:} For $w \in V^\perp$:
\[\langle Aw, v_1 \rangle = \langle w, A^T v_1 \rangle = \langle w, Av_1 \rangle = \langle w, \lambda_1 v_1 \rangle = \lambda_1 \langle w, v_1 \rangle = 0\]
using that $A = A^T$ (symmetric).

\textit{Step 3: Apply induction.}

Choose an orthonormal basis $\{w_1, \ldots, w_{n-1}\}$ for $V^\perp$. The restriction $A|_{V^\perp}: V^\perp \to V^\perp$ is represented by some $(n-1) \times (n-1)$ matrix $A_0$ relative to this basis.

One can check that $A_0$ is also symmetric (same argument using $A = A^T$).

By the inductive hypothesis, there exists an orthonormal basis $\{u_2, \ldots, u_n\}$ of $\mathbb{R}^{n-1}$ consisting of eigenvectors of $A_0$:
\[A_0 u_j = \lambda_j u_j \text{ for } j = 2, \ldots, n\]

Define $v_j = \sum_{i=1}^{n-1} (u_j)_i w_i \in V^\perp$ for $j = 2, \ldots, n$.

Then $\{v_1, v_2, \ldots, v_n\}$ is an orthonormal basis for $\mathbb{R}^n$ (since $v_1 \perp V^\perp$ and $\{v_2, \ldots, v_n\}$ is orthonormal in $V^\perp$), and each $v_j$ is an eigenvector of $A$.
\end{proof}

\begin{cor}{Diagonalization of Symmetric Matrices}{diag_symmetric}
If $A$ is a real symmetric matrix, then there exists an orthogonal matrix $Q$ such that $Q^T A Q$ is diagonal. The diagonal entries are the eigenvalues of $A$.
\end{cor}

\begin{proof}
Let $\{v_1, \ldots, v_n\}$ be the orthonormal eigenvector basis from the Spectral Theorem. Let $Q$ be the matrix with columns $v_1, \ldots, v_n$. Then $Q$ is orthogonal and
\[Q^T A Q = \begin{pmatrix} \lambda_1 & & \\ & \ddots & \\ & & \lambda_n \end{pmatrix}\]
\end{proof}


% ==================== WEEK 8 ====================
\chapter{Integration and the Fundamental Theorem}

\section{Analysis: Mon/Tue}

\subsection{The Riemann Integral}

\begin{defn}{Riemann Integral}{riemann_integral}
For $f: [a,b] \to \mathbb{R}$:
\begin{enumerate}
    \item A \textbf{partition} $P$ is a choice $a = t_0 < t_1 < \cdots < t_n = b$.
    \item Define
    \[L(f, P) = \sum_{i=1}^n \inf_{t \in [t_{i-1}, t_i]} f(t) \cdot (t_i - t_{i-1})\]
    \[U(f, P) = \sum_{i=1}^n \sup_{t \in [t_{i-1}, t_i]} f(t) \cdot (t_i - t_{i-1})\]
    \item Define $\underline{\int} f = \sup_P L(f,P)$ and $\overline{\int} f = \inf_P U(f,P)$.
    \item $f$ is \textbf{integrable} if $\underline{\int} f = \overline{\int} f$, and we write this value as $\int_a^b f(t) \, dt$.
\end{enumerate}
\end{defn}

\begin{thm}{Properties of the Integral}{int_properties}
Let $f, g: [a,b] \to \mathbb{R}$ be integrable and $c_1, c_2 \in \mathbb{R}$. Then:
\begin{enumerate}
    \item $c_1 f + c_2 g$ is integrable and $\int (c_1 f + c_2 g) = c_1 \int f + c_2 \int g$
    \item (\textbf{Positivity}) If $f \geq 0$, then $\int f \geq 0$
    \item (\textbf{Boundedness}) $|\int f| \leq \sup|f| \cdot (b-a)$
    \item (\textbf{Additivity}) If $a < c < b$, then $\int_a^b f = \int_a^c f + \int_c^b f$
\end{enumerate}
\end{thm}

\subsection{Fundamental Theorem of Calculus}

\begin{thm}{FTC Version 1}{ftc1}
Suppose $f: [a,b] \to \mathbb{R}$ is continuous. Define $F: [a,b] \to \mathbb{R}$ by
\[F(t) = \int_a^t f(s) \, ds\]
Then $F$ is $C^1$ and $F'(t) = f(t)$ for all $t \in (a,b)$.
\end{thm}

\begin{proof}
Suffices to show $F'(t) = f(t)$ (then $f$ cts implies $F'$ cts).

Let $\varepsilon > 0$. Consider $|h| > 0$.

$\left|\frac{F(t+h) - F(t)}{h} - f(t)\right| = \left|\frac{1}{h}\int_t^{t+h} f(s) \, ds - f(t)\right| = \left|\frac{1}{h}\int_t^{t+h} (f(s) - f(t)) \, ds\right|$

By continuity of $f$, $\exists \delta > 0$ s.t. $|f(s) - f(t)| < \varepsilon$ whenever $|s - t| < \delta$.

If $|h| < \delta$:
\[\left|\frac{1}{h}\int_t^{t+h} (f(s) - f(t)) \, ds\right| \leq \frac{1}{|h|} \cdot \varepsilon \cdot |h| = \varepsilon\]
\end{proof}

\begin{thm}{FTC Version 2}{ftc2}
Suppose $f: [a,b] \to \mathbb{R}$ is $C^1$. Then
\[\int_a^b f'(t) \, dt = f(b) - f(a)\]
\end{thm}

\begin{proof}
Define $F(t) = \int_a^t f'(s) \, ds$. By FTC v1, $F$ is $C^1$ and $F'(t) = f'(t)$.

Let $g(t) = F(t) - f(t)$. Then $g$ is $C^1$ and $g'(t) = F'(t) - f'(t) = 0$ for all $t$.

\textbf{Lemma:} If $f$ is continuous on $[a,b]$, differentiable on $(a,b)$, and $f'(t) = 0$ for all $t$, then $f$ is constant.

Hence $g$ is constant, so $F(a) - f(a) = g(a) = g(b) = F(b) - f(b)$.

Rearranging: $f(b) - f(a) = F(b) - F(a) = F(b) = \int_a^b f'(s) \, ds$.
\end{proof}

\subsection{Integration by Parts}

\begin{thm}{Integration by Parts}{int_by_parts}
Let $f, g: [a,b] \to \mathbb{R}$ be $C^1$. Then
\[\int_a^b f'(t) g(t) \, dt = f(b)g(b) - f(a)g(a) - \int_a^b f(t) g'(t) \, dt\]
\end{thm}

\begin{proof}
Apply FTC v2 to $fg: [a,b] \to \mathbb{R}$:
\[f(b)g(b) - f(a)g(a) = \int_a^b (fg)'(t) \, dt = \int_a^b f'g \, dt + \int_a^b fg' \, dt\]
Rearranging gives the result.
\end{proof}

\section{Linear Algebra: Thu/Fri (Computing Inverses)}

\subsection{Finding Inverses via Row Reduction}

\begin{thm}{Inverse via Augmented Matrix}{inverse_augmented}
Augment $(A \mid I)$ and row reduce. If $\text{rref}(A) = I$, then $\text{rref}(A \mid I) = (I \mid A^{-1})$.
\end{thm}

\begin{proof}
Recall that row operations correspond to multiplying by matrices on the left.

Hence if $R_1, \ldots, R_k$ are the row operations:
\[\text{rref}(A \mid I) = R_k \cdots R_1 (A \mid I) = (R_k \cdots R_1 A \mid R_k \cdots R_1 I)\]

Here $B = R_k \cdots R_1 = A^{-1}$.
\end{proof}

\begin{remark}
The row operations correspond to matrices:
\begin{itemize}
    \item Switching rows $i$ and $j$: permutation matrix
    \item Multiply row $i$ by $\lambda \neq 0$: diagonal matrix with $\lambda$ in position $(i,i)$
    \item Add $\lambda \times$ row $j$ to row $i$: elementary matrix with $\lambda$ in position $(i,j)$
\end{itemize}
All of these are invertible!
\end{remark}


% ==================== WEEK 9 ====================
\chapter{Curves and Submanifolds}

\section{Analysis: Mon/Tue}

\subsection{Curves}

\begin{defn}{Curve}{curve_def}
A \textbf{curve} in $\mathbb{R}^m$ is a map $\gamma: [a,b] \to \mathbb{R}^m$.

We say $\gamma$ is a $C^k$ curve for $k \in \mathbb{N} \cup \{\infty\}$ if each component function is $C^k$.
\end{defn}

\begin{defn}{Velocity and Unit Tangent}{velocity_tangent}
\begin{enumerate}
    \item If $\gamma: [a,b] \to \mathbb{R}^m$ is $C^1$, define the \textbf{velocity vector} to be $\gamma'(t) \in \mathbb{R}^m$.
    \item If $\gamma'(t) \neq 0$, define the \textbf{unit tangent vector} to be $T(t) = \frac{\gamma'(t)}{\|\gamma'(t)\|}$.
\end{enumerate}
\end{defn}

\subsection{Arc Length}

\begin{defn}{Length of a Curve}{arc_length_def}
For a partition $P: a = t_0 < t_1 < \cdots < t_n = b$, define
\[L(\gamma, P) = \sum_{i=1}^n \|\gamma(t_i) - \gamma(t_{i-1})\|\]
We say $\gamma$ has \textbf{finite length} if $\{L(\gamma, P) : P \text{ a partition}\}$ is bounded above.

If $\gamma$ has finite length, define $L(\gamma) = \sup_P L(\gamma, P)$.
\end{defn}

\begin{thm}{Arc Length Formula}{arc_length_formula}
Suppose $\gamma: [a,b] \to \mathbb{R}^m$ is $C^1$. Then $\gamma$ has finite length and
\[L(\gamma) = \int_a^b \|\gamma'(t)\| \, dt\]
\end{thm}

\begin{proof}[Sketch]
\textbf{Step 1:} $\gamma$ is of finite length.

For any partition $P$:
\[L(\gamma, P) = \sum \|\gamma(t_i) - \gamma(t_{i-1})\| = \sum \left\|\int_{t_{i-1}}^{t_i} \gamma'(s) \, ds\right\| \leq \sum \int_{t_{i-1}}^{t_i} \|\gamma'(s)\| \, ds = \int_a^b \|\gamma'(s)\| \, ds\]

Hence $\{L(\gamma, P)\}$ is bounded above by $\int_a^b \|\gamma'\| \, dt$.

\textbf{Step 2:} Define $S(t) = L(\gamma|_{[a,t]})$.

\textbf{Claim:} $S'(t)$ exists and $S'(t) = \|\gamma'(t)\|$.

Using FTC: $S(b) = S(a) + \int_a^b S'(t) \, dt = 0 + \int_a^b \|\gamma'(t)\| \, dt$.
\end{proof}

\begin{lem}{Additivity of Length}{length_additive}
$\gamma: [a,b] \to \mathbb{R}^m$ has finite length. Suppose $c \in (a,b)$. Then $\gamma|_{[a,c]}$ and $\gamma|_{[c,b]}$ have finite lengths and
\[L(\gamma) = L(\gamma|_{[a,c]}) + L(\gamma|_{[c,b]})\]
\end{lem}

\subsection{Submanifolds}

\begin{defn}{Graph Map}{graph_map}
Let $U \subseteq \mathbb{R}^k$ be open, $g: U \to \mathbb{R}^{n-k}$ be $C^1$.

Define the \textbf{graph map} $\Gamma: U \to \mathbb{R}^n$ by
\[\Gamma(x) = (x, g(x)) = (x_1, \ldots, x_k, g_1(x_1, \ldots, x_k), \ldots, g_{n-k}(x_1, \ldots, x_k))\]
\end{defn}

\begin{defn}{$k$-Dimensional Submanifold}{submanifold}
We say $M \subseteq \mathbb{R}^n$ is a \textbf{$k$-dimensional $C^1$ submanifold} if for all $a \in M$, $\exists \rho > 0$ s.t.
\[M \cap B_\rho(a) = P_{j_1, \ldots, j_k}^{-1}(\Gamma(U))\]
where:
\begin{itemize}
    \item $\Gamma$ is the graph map for some $C^1$ function $g: U \to \mathbb{R}^{n-k}$, $U \subseteq \mathbb{R}^k$ open
    \item $P_{j_1, \ldots, j_k}$ is a permutation of coordinates
\end{itemize}
I.e., $M \cap B_\rho(a)$ can be represented as a graph after permuting coordinates.
\end{defn}

\begin{thm}{Level Set Theorem}{level_set}
Suppose $\phi_1, \ldots, \phi_{n-k}: \mathbb{R}^n \to \mathbb{R}$ are $C^1$.

Let $M = \{x \in \mathbb{R}^n : \phi_1(x) = \cdots = \phi_{n-k}(x) = 0\}$.

If $\nabla\phi_1(x), \ldots, \nabla\phi_{n-k}(x)$ are linearly independent for all $x \in M$, then (if $M \neq \emptyset$) $M$ is a $k$-dimensional $C^1$ submanifold.
\end{thm}

\begin{example}
$S^{n-1} = \{x \in \mathbb{R}^n : \|x\| = 1\}$.

Here $\phi(x) = \|x\|^2 - 1$. We have $\nabla\phi(x) = 2x$, which is linearly independent when $\phi(x) = 0$ (i.e., $x \neq 0$).

So $S^{n-1}$ is an $(n-1)$-dimensional submanifold.
\end{example}

\subsection{Tangent Spaces}

\begin{defn}{Tangent Space}{tangent_space}
Let $M \subseteq \mathbb{R}^n$ be a $k$-dimensional submanifold, $a \in M$.

Define the \textbf{tangent space} of $M$ at $a$ by
\[T_a M = \{\gamma'(t_0) : \gamma: (-\varepsilon, \varepsilon) \to M \text{ is a $C^1$ curve s.t. } \gamma(t_0) = a\}\]
\end{defn}

\begin{thm}{Tangent Space is a Subspace}{tangent_subspace}
$T_a M$ is a $k$-dimensional linear subspace of $\mathbb{R}^n$.
\end{thm}

\begin{proof}[Sketch]
Suffices to consider the case of a graph. If $M \cap B_\rho(a) = \Gamma(U)$ where $\Gamma(x) = (x, g(x))$, then:

\textbf{Claim:} $T_a M = \text{span}(D_1 \Gamma(a), \ldots, D_k \Gamma(a))$.

Here $D_i \Gamma(a) = (e_i, D_i g(a))$.

These are linearly independent (first $k$ components are $e_1, \ldots, e_k$), so $\dim T_a M = k$.
\end{proof}

\section{Linear Algebra: Thu/Fri (Abstract Vector Spaces)}

\begin{remark}
Many results from $\mathbb{R}^n$ extend to abstract vector spaces over any field $\mathbb{F}$. The key tools:
\begin{itemize}
    \item Bases and dimension
    \item Linear maps and their matrices
    \item Kernel, range, rank-nullity theorem
    \item Change of basis
\end{itemize}
For inner product spaces specifically (over $\mathbb{R}$ or $\mathbb{C}$):
\begin{itemize}
    \item Orthogonality and orthonormal bases
    \item Gram-Schmidt process
    \item Orthogonal projections
\end{itemize}
\end{remark}


% ==================== WEEK 10 ====================
\chapter{Power Series and Integration in Higher Dimensions}

\section{Analysis: Mon/Tue}

\subsection{Power Series}

Consider a power series $f(x) = \sum_{n=0}^\infty a_n x^n$ where the coefficients $a_n \in \mathbb{R}$.

\begin{lem}{Convergence Propagates}{conv_propagates}
Suppose $\exists x_0 \neq 0$ s.t. $\sum a_n x_0^n$ converges. Then:
\begin{enumerate}
    \item For all $x$ with $|x| < |x_0|$, the power series converges absolutely.
    \item For all $R' < |x_0|$ and any polynomial $p(n)$, $\sum p(n) a_n x^n$ converges absolutely and uniformly on $[-R', R']$.
\end{enumerate}
\end{lem}

\begin{proof}[Sketch]
Let $r = |x|/|x_0| < 1$.

Since $\sum a_n x_0^n$ converges, $\lim_{n \to \infty} a_n x_0^n = 0$. So $\exists B > 0$ s.t. $|a_n x_0^n| \leq B$ for all $n$.

Hence $|a_n x^n| = |a_n x_0^n| \cdot r^n \leq B r^n$.

Since $\sum B r^n$ converges (geometric series with $r < 1$), we get absolute convergence.

For uniform convergence, note that $|p(n)| \cdot |a_n| \cdot |x|^n \leq |p(n)| B r^n$ on $[-R', R']$, and $\sum |p(n)| r^n$ converges.
\end{proof}

\begin{thm}{Radius of Convergence}{radius_conv}
For any power series $\sum a_n x^n$, exactly one of the following holds:
\begin{enumerate}
    \item $\forall x \in \mathbb{R}$, $f(x)$ is convergent (radius $R = \infty$)
    \item $f(x)$ is divergent $\forall x \neq 0$ (radius $R = 0$)
    \item $\exists R > 0$ s.t. $f(x)$ is convergent if $|x| < R$ and divergent if $|x| > R$
\end{enumerate}
In case (3), nothing is claimed for $x = \pm R$.
\end{thm}

\begin{defn}{Radius of Convergence}{radius_def}
The \textbf{radius of convergence} is $\infty$ in case (1), $0$ in case (2), and $R$ in case (3).
\end{defn}

\subsection{Uniform Convergence}

\begin{defn}{Uniform Convergence}{unif_conv}
Let $f_j: S \to \mathbb{R}$. We say $f_j$ \textbf{converges to $f$ uniformly} if
\[\lim_{j \to \infty} \sup_{x \in S} |f_j(x) - f(x)| = 0\]
I.e., $\forall \varepsilon > 0$, $\exists N$ s.t. $|f_j(x) - f(x)| < \varepsilon$ for all $j \geq N$ and all $x \in S$.
\end{defn}

\begin{thm}{Uniform Limit of Continuous Functions}{unif_cts}
Suppose $f_j: S \to \mathbb{R}$ are continuous and $f_j \to f$ uniformly. Then $f$ is continuous.
\end{thm}

\begin{thm}{Weierstrass M-Test}{m_test}
Suppose $\sum g_j: S \to \mathbb{R}$ satisfies $|g_j(x)| \leq M_j$ for some $M_j$ with $\sum M_j$ convergent.

Then $\sum g_j$ converges uniformly to some $g$.

In particular, if each $g_j$ is continuous, then $g$ is also continuous.
\end{thm}

\begin{thm}{Continuity of Power Series}{power_cts}
Let $f(x) = \sum a_n x^n$ have radius of convergence $R > 0$. Then $f$ is continuous on $(-R, R)$.
\end{thm}

\begin{thm}{Differentiability of Power Series}{power_diff}
Let $f(x) = \sum a_n x^n$ have radius of convergence $R > 0$. Then:
\begin{enumerate}
    \item $f$ is $C^k$ on $(-R, R)$ for all $k$
    \item $f^{(k)}(x) = \sum_{n \geq k} n(n-1)\cdots(n-k+1) a_n x^{n-k}$
\end{enumerate}
This means we can differentiate term by term!
\end{thm}

\subsection{Taylor Series}

\begin{thm}{Taylor's Theorem (General)}{taylor_general}
Let $f: [x-r, x+r] \to \mathbb{R}$ be $C^{k+1}$ for some $L \in \mathbb{R}$, $r > 0$. Then
\[f(x) = \sum_{j=0}^k \frac{f^{(j)}(L)}{j!}(x-L)^j + R_k(x)\]
where the remainder satisfies $R_k(x) = \frac{f^{(k+1)}(\xi)}{(k+1)!}(x-L)^{k+1}$ for some $\xi$ between $L$ and $x$.
\end{thm}

\begin{cor}{Power Series Representation}{power_rep}
If $f: (L-r, L+r) \to \mathbb{R}$ is $C^\infty$ and $|f^{(k)}(x)| \to 0$ sufficiently fast, then
\[f(x) = \sum_{n=0}^\infty \frac{f^{(n)}(L)}{n!}(x-L)^n\]
In particular, $a_n = \frac{f^{(n)}(L)}{n!}$.
\end{cor}

\subsection{Integration in $\mathbb{R}^n$}

\begin{defn}{Rectangle and Partition}{rectangle}
\begin{enumerate}
    \item A \textbf{rectangle} $R = R[a_1, b_1; \ldots; a_n, b_n] = \{(x_1, \ldots, x_n) : a_i \leq x_i \leq b_i\}$
    \item $\text{Vol}(R) = (b_1 - a_1) \cdots (b_n - a_n)$
    \item A \textbf{partition} $P$ of $R$ is a choice of partitions $P_i$ of $[a_i, b_i]$ for each $i$
\end{enumerate}
\end{defn}

\begin{defn}{Riemann Integral in $\mathbb{R}^n$}{riemann_rn}
For $f: R \to \mathbb{R}$ bounded:
\[L(f, P) = \sum_{R' \text{ subrectangle}} \inf_{x \in R'} f(x) \cdot \text{Vol}(R')\]
\[U(f, P) = \sum_{R' \text{ subrectangle}} \sup_{x \in R'} f(x) \cdot \text{Vol}(R')\]

Define $\underline{\int} f = \sup_P L(f,P)$, $\overline{\int} f = \inf_P U(f,P)$.

$f$ is \textbf{integrable} if $\underline{\int} f = \overline{\int} f$, written $\int_R f$.
\end{defn}

\begin{thm}{Fubini's Theorem}{fubini}
Given $f: R \to \mathbb{R}$ integrable, $R \subseteq \mathbb{R}^n$ a rectangle, $R = R[a_1, b_1; \ldots; a_n, b_n]$.

For $x \in \mathbb{R}^n$, write $x = (y, z)$ where $y \in \mathbb{R}^k$ and $z \in \mathbb{R}^{n-k}$.

Then
\[\int_R f \, dx_1 \cdots dx_n = \int_{R'} \left(\int_{R''} f(y, z) \, dz\right) dy\]
where $R'$ and $R''$ are the projections of $R$ onto the $y$ and $z$ coordinates.

I.e., we can compute higher-dimensional integrals as iterated 1-D integrals in any order!
\end{thm}

\section{Linear Algebra: Thu/Fri (Review and Connections)}

\begin{remark}
The main theorems connecting linear algebra and analysis in this course:
\begin{enumerate}
    \item \textbf{Derivative as a linear map:} $(Df)(a): \mathbb{R}^n \to \mathbb{R}^m$ is a linear transformation, represented by the Jacobian matrix.
    \item \textbf{Chain rule as matrix multiplication:} $D(g \circ f)(a) = Dg(f(a)) \cdot Df(a)$.
    \item \textbf{Tangent spaces as subspaces:} $T_a M$ is a linear subspace of $\mathbb{R}^n$.
    \item \textbf{Hessian as a quadratic form:} The second derivative test uses positive definiteness of the Hessian matrix.
    \item \textbf{Inverse Function Theorem:} $f$ is locally invertible near $a$ iff $Df(a)$ is invertible (i.e., $\det Df(a) \neq 0$).
\end{enumerate}
\end{remark}

\section{Advanced Analysis: Contraction Mapping and Function Theorems}

\subsection{Contraction Mapping Principle}

\begin{defn}{Contraction}{contraction}
Let $E \subseteq \mathbb{R}^n$. A map $f: E \to \mathbb{R}^n$ is a \textbf{contraction} if there exists $\theta \in (0, 1)$ such that
\[\|f(x) - f(y)\| \leq \theta \|x - y\| \quad \forall x, y \in E\]
\end{defn}

\begin{thm}{Contraction Mapping Theorem}{contraction_mapping}
Let $E \subseteq \mathbb{R}^n$ be closed and let $f: E \to E$ be a contraction (so $f(E) \subseteq E$). Then $f$ has a unique fixed point, i.e., there exists a unique $z \in E$ such that $f(z) = z$.
\end{thm}

\begin{proof}
The proof is constructive --- we'll actually find the fixed point!

Pick any $x_0 \in E$ and define iteratively:
\[x_k = f(x_{k-1}) \text{ for } k \geq 1\]

\textbf{Step 1: The sequence is Cauchy.}

By the contraction property:
\[\|x_{k+1} - x_k\| = \|f(x_k) - f(x_{k-1})\| \leq \theta \|x_k - x_{k-1}\|\]

By induction: $\|x_{k+1} - x_k\| \leq \theta^k \|x_1 - x_0\|$.

For $\ell > k$:
\begin{align*}
\|x_\ell - x_k\| &\leq \sum_{j=k}^{\ell-1} \|x_{j+1} - x_j\| \\
&\leq \sum_{j=k}^{\ell-1} \theta^j \|x_1 - x_0\| \\
&\leq \frac{\theta^k}{1 - \theta} \|x_1 - x_0\|
\end{align*}
Since $\theta^k \to 0$, the sequence is Cauchy.

\textbf{Step 2: The limit is a fixed point.}

Since $E$ is closed and $\{x_k\}$ is Cauchy in $E$, it converges to some $z \in E$.

Since $f$ is continuous (contractions are Lipschitz, hence continuous):
\[f(z) = f(\lim_{k \to \infty} x_k) = \lim_{k \to \infty} f(x_k) = \lim_{k \to \infty} x_{k+1} = z\]

\textbf{Step 3: Uniqueness.}

If $f(z) = z$ and $f(w) = w$, then $\|z - w\| = \|f(z) - f(w)\| \leq \theta \|z - w\|$.

Since $\theta < 1$, this forces $\|z - w\| = 0$, so $z = w$.
\end{proof}

\begin{remark}
The contraction mapping theorem is incredibly useful! It gives both existence and uniqueness, plus a practical algorithm (iteration) to find the fixed point.
\end{remark}

\subsection{Inverse Function Theorem}

The Inverse Function Theorem tells us when a $C^1$ map is locally invertible.

\begin{thm}{Inverse Function Theorem}{IFT}
Let $U \subseteq \mathbb{R}^n$ be open, $f: U \to \mathbb{R}^n$ be $C^1$, and $x_0 \in U$ with $Df(x_0)$ invertible. Then there exist open sets $V, W$ with $x_0 \in V$ and $f(x_0) \in W$ such that:
\begin{enumerate}
    \item $f|_V: V \to W$ is a bijection
    \item The inverse $g: W \to V$ is also $C^1$
    \item $Dg(y) = (Df(g(y)))^{-1}$ for all $y \in W$
\end{enumerate}
\end{thm}

\begin{proof}[Sketch]
\textbf{Step 1: Reduce to the case $Df(x_0) = I$.}

If $Df(x_0) = A$ is invertible, consider $\tilde{f}(x) = A^{-1} f(x)$. Then $D\tilde{f}(x_0) = A^{-1} A = I$.

If the theorem holds for $\tilde{f}$, it holds for $f$.

\textbf{Step 2: Setup.}

Since $Df$ is continuous and $Df(x_0) = I$, there exists $\rho > 0$ such that $\|Df(x) - I\| < \frac{1}{2}$ for $x \in B_\rho(x_0)$.

\textbf{Step 3: Injectivity.}

For $x, y \in B_\rho(x_0)$:
\[\|f(x) - f(y) - (x - y)\| = \left\|\int_0^1 (Df(x + t(y-x)) - I)(y - x) \, dt\right\| \leq \frac{1}{2}\|x - y\|\]

Hence $\|f(x) - f(y)\| \geq \frac{1}{2}\|x - y\|$, so $f$ is injective on $B_\rho(x_0)$.

\textbf{Step 4: Surjectivity (using contraction mapping).}

For any $v$ near $f(x_0)$, we want to find $x$ with $f(x) = v$.

Define $F_v(x) = x - f(x) + v$. We show $F_v$ is a contraction on a closed ball centered at $x_0$:

$\|F_v(x) - F_v(y)\| = \|(x - f(x)) - (y - f(y))\| \leq \frac{1}{2}\|x - y\|$.

By the Contraction Mapping Theorem, $F_v$ has a fixed point $x$, meaning $x = x - f(x) + v$, i.e., $f(x) = v$.

\textbf{Step 5: Differentiability of inverse.}

The formula $Dg(y) = (Df(g(y)))^{-1}$ follows from differentiating $f(g(y)) = y$ using the chain rule.
\end{proof}

\subsection{Implicit Function Theorem}

The Implicit Function Theorem is closely related to the IFT and tells us when we can solve equations locally.

\begin{thm}{Implicit Function Theorem}{ImplicitFT}
Let $U \subseteq \mathbb{R}^n \times \mathbb{R}^m$ be open and $F: U \to \mathbb{R}^m$ be $C^1$. Suppose $(a, b) \in U$ satisfies $F(a, b) = 0$ and the partial derivative $D_y F(a, b): \mathbb{R}^m \to \mathbb{R}^m$ is invertible.

Then there exist open neighborhoods $V$ of $a$ in $\mathbb{R}^n$ and $W$ of $b$ in $\mathbb{R}^m$, and a unique $C^1$ function $g: V \to W$ such that:
\begin{enumerate}
    \item $g(a) = b$
    \item $F(x, g(x)) = 0$ for all $x \in V$
    \item $Dg(x) = -(D_y F(x, g(x)))^{-1} D_x F(x, g(x))$
\end{enumerate}
\end{thm}

\begin{proof}[Sketch]
Define $\Phi: U \to \mathbb{R}^n \times \mathbb{R}^m$ by $\Phi(x, y) = (x, F(x, y))$.

Then $D\Phi(a, b) = \begin{pmatrix} I & 0 \\ D_x F & D_y F \end{pmatrix}$, which is invertible since $D_y F$ is invertible.

By the Inverse Function Theorem, $\Phi$ has a local inverse $\Psi = (\Psi_1, \Psi_2)$.

For $(x, 0)$ near $(a, 0)$: $\Psi(x, 0) = (x, g(x))$ for some function $g$.

Then $\Phi(x, g(x)) = (x, 0)$ implies $F(x, g(x)) = 0$.
\end{proof}

\begin{example}
\textbf{(Inverse Function Theorem)}
Consider $f: \mathbb{R}^2 \to \mathbb{R}^2$ given by $f(x, y) = (e^x \cos y, e^x \sin y)$ (polar to Cartesian in disguise).

At $(0, 0)$: $Df(0,0) = \begin{pmatrix} 1 & 0 \\ 0 & 1 \end{pmatrix}$ is invertible.

So $f$ is locally invertible near $(0,0)$. The local inverse is essentially $(r, \theta) \mapsto (\ln r, \theta)$ for $r > 0$ near 1.
\end{example}

\begin{example}
\textbf{(Implicit Function Theorem)}
Consider $F(x, y, z) = x^2 + y^2 + z^2 - 1$ (the unit sphere). At $(0, 0, 1)$:
\[D_z F = 2z = 2 \neq 0\]

So near $(0, 0)$, we can solve for $z = g(x, y) = \sqrt{1 - x^2 - y^2}$ (the upper hemisphere).

The derivative is: $Dg = -\frac{1}{2z}(2x, 2y) = \left(-\frac{x}{\sqrt{1-x^2-y^2}}, -\frac{y}{\sqrt{1-x^2-y^2}}\right)$
\end{example}

\begin{example}
\textbf{(Implicit Function Theorem: Level Curves)}
$F(x, y) = x^3 + y^3 - 3xy$ (the folium of Descartes). At $(1, 1)$: $F(1,1) = 1 + 1 - 3 = -1 \neq 0$.

But at $(3/2, 3/2)$: $F(3/2, 3/2) = 27/8 + 27/8 - 27/4 = 0$ and $F_y = 3y^2 - 3x = 3(9/4) - 9/2 = 27/4 - 9/2 = 9/4 \neq 0$.

So locally near $(3/2, 3/2)$, we can write $y = g(x)$ for some smooth function $g$.

The formula for $Dg$ follows from differentiating $F(x, g(x)) = 0$.
\end{example}

\begin{remark}
The IFT and Implicit FT are fundamental tools in differential geometry --- they let us show that submanifolds defined by equations are ``nice'' smooth objects.
\end{remark}


% =======================================================
% BACK MATTER: SOLUTIONS
% =======================================================
\backmatter 

\chapter{Solutions to Selected Exercises}
\markboth{Solutions}{Solutions} 

% --- Solutions for Part I (Math 61) ---
\section*{Part I: Math 61}

\subsection*{Week 1: The Real Numbers}

\begin{solution}[1.1]
(i) Assume for contradiction that $-a > 0$. Then $a + (-a) > 0$ (O2), so $0 > 0$ (F5), contradiction. Also $-a \neq 0$ since $a \neq 0$. Hence $-a < 0$.

(ii) Assume $\frac{1}{a} < 0$. Then $a \cdot (-\frac{1}{a}) > 0$ (O2). But $a \cdot (-\frac{1}{a}) = -1$, so $(-1)(-1) = 1 > 0$ implies $1 + (-1) > 0$, contradicting F5. Also $\frac{1}{a} \neq 0$ since $a \cdot \frac{1}{a} = 1 \neq 0$. Hence $\frac{1}{a} > 0$.

(iii) Assume $\frac{1}{b} < \frac{1}{a}$. Multiplying by $ab > 0$ (using O2 and part (ii)): $a < b$, contradiction. Hence $\frac{1}{a} < \frac{1}{b}$.
\end{solution}

\begin{solution}[1.3]
Let $S = \{x : x^2 < a, x > 0\}$. First, $S$ is nonempty (take $x = \min(1, a/2)$) and bounded above (by $\max(1, a)$). Let $\alpha = \sup S$.

\textbf{Claim:} $\alpha^2 = a$.

If $\alpha^2 > a$: Choose $n$ large enough that $\frac{1}{n} < \frac{\alpha^2 - a}{2\alpha}$. Then $(\alpha - \frac{1}{n})^2 > a$, so $\alpha - \frac{1}{n}$ is an upper bound for $S$, contradicting $\alpha = \sup S$.

If $\alpha^2 < a$: Choose $n$ large enough that $(\alpha + \frac{1}{n})^2 < a$. Then $\alpha + \frac{1}{n} \in S$, contradicting $\alpha = \sup S$.

Hence $\alpha^2 = a$, and $\alpha = \sqrt{a}$.
\end{solution}

\begin{solution}[1.5]
Let $\varepsilon > 0$. By Archimedean property, $\exists N \in \mathbb{N}$ with $N > 1/\varepsilon$, i.e., $\frac{1}{N} < \varepsilon$. For $n \geq N$: $\frac{1}{n} \leq \frac{1}{N} < \varepsilon$, so $|\frac{1}{n} - 0| < \varepsilon$. Hence $\lim \frac{1}{n} = 0$.
\end{solution}

\begin{solution}[1.7]
No! Note that $(v_1 - v_2) + (v_2 - v_3) + (v_3 - v_4) + (v_4 - v_1) = 0$. This is a nontrivial linear combination equaling zero, so they are linearly dependent.
\end{solution}

\subsection*{Week 2: Subsequences and Continuity}

\begin{solution}[2.1]
Let $L = \lim a_n = \lim b_n$. Given $\varepsilon > 0$, choose $N$ such that $|a_n - L| < \varepsilon$ and $|b_n - L| < \varepsilon$ for $n \geq N$.

Since $a_n \leq c_n \leq b_n$, we have $a_n - L \leq c_n - L \leq b_n - L$.

If $c_n - L \geq 0$: $|c_n - L| = c_n - L \leq b_n - L \leq |b_n - L| < \varepsilon$.

If $c_n - L < 0$: $|c_n - L| = L - c_n \leq L - a_n = |a_n - L| < \varepsilon$.

Either way, $|c_n - L| < \varepsilon$ for $n \geq N$. Hence $\lim c_n = L$.
\end{solution}

\begin{solution}[2.3]
Suppose $f$ is not continuous at $c$. Then $\exists \varepsilon_0 > 0$ such that $\forall \delta > 0$, $\exists x$ with $|x - c| < \delta$ but $|f(x) - f(c)| \geq \varepsilon_0$.

Choose $\delta = \frac{1}{n}$. Then $\exists x_n$ with $|x_n - c| < \frac{1}{n}$ but $|f(x_n) - f(c)| \geq \varepsilon_0$.

This sequence $\{x_n\}$ satisfies $x_n \to c$, but $f(x_n) \not\to f(c)$, contradicting the hypothesis. Hence $f$ is continuous at $c$.
\end{solution}

\begin{solution}[2.5]
($\Rightarrow$) If $U = V$, they have the same basis, so $\dim U = \dim V$.

($\Leftarrow$) Let $\{u_1, \ldots, u_n\}$ be a basis for $U$. Since $U \subseteq V$ and these are linearly independent vectors in $V$ with $|$set$| = \dim V$, they form a basis for $V$. Hence $V = \text{span}\{u_1, \ldots, u_n\} = U$.
\end{solution}

\begin{solution}[2.7]
$V = \{p \in \mathcal{P}_4 : p(1) = 0\}$ consists of polynomials with root at $x = 1$, so $p(x) = (x-1)q(x)$ where $\deg q \leq 3$.

\textbf{Basis for $V$:} $\{(x-1), x(x-1), x^2(x-1), x^3(x-1)\}$.

These span $V$ since any $(x-1)(a + bx + cx^2 + dx^3)$ is a linear combination. They're independent since $(x-1)(c_1 + c_2 x + c_3 x^2 + c_4 x^3) = 0$ implies all $c_i = 0$.

\textbf{Extend to $\mathcal{P}_4$:} Add $\{1\}$. Check: $1 \notin V$ (since $1(1) = 1 \neq 0$). So basis for $\mathcal{P}_4$ is $\{1, (x-1), x(x-1), x^2(x-1), x^3(x-1)\}$.
\end{solution}

\subsection*{Week 3: Differentiation and Topology}

\begin{solution}[3.1]
Suppose $f: [a,b] \to \mathbb{R}$ is continuous but not uniformly continuous. Then $\exists \varepsilon_0 > 0$ such that $\forall \delta > 0$, $\exists x, y$ with $|x - y| < \delta$ but $|f(x) - f(y)| \geq \varepsilon_0$.

Take $\delta = 1/n$. Get sequences $\{x_n\}, \{y_n\}$ with $|x_n - y_n| < 1/n$ and $|f(x_n) - f(y_n)| \geq \varepsilon_0$.

By Bolzano-Weierstrass, $\{x_n\}$ has a convergent subsequence $x_{n_k} \to c \in [a,b]$. Since $|x_{n_k} - y_{n_k}| < 1/n_k \to 0$, we have $y_{n_k} \to c$ also.

By continuity: $f(x_{n_k}) \to f(c)$ and $f(y_{n_k}) \to f(c)$, so $|f(x_{n_k}) - f(y_{n_k})| \to 0$. But this contradicts $|f(x_{n_k}) - f(y_{n_k})| \geq \varepsilon_0$. Hence $f$ is uniformly continuous.
\end{solution}

\begin{solution}[3.3]
($\Rightarrow$) If $a_n \to L$, then $\forall \varepsilon > 0$, $\exists N$ with $|a_n - L| < \varepsilon/2$ for $n \geq N$. For $m, n \geq N$:
\[|a_m - a_n| \leq |a_m - L| + |a_n - L| < \varepsilon/2 + \varepsilon/2 = \varepsilon\]

($\Leftarrow$) First show $\{a_n\}$ is bounded: Take $\varepsilon = 1$. Then $\exists N$ with $|a_m - a_n| < 1$ for $m, n \geq N$. So $|a_m| \leq |a_N| + 1$ for $m \geq N$. Hence $\{a_n\}$ is bounded.

By BW, $\exists$ subsequence $a_{n_k} \to A$. Given $\varepsilon > 0$, choose $K$ so $|a_{n_k} - A| < \varepsilon/2$ for $k \geq K$, and choose $N$ so $|a_m - a_n| < \varepsilon/2$ for $m, n \geq N$. For $n \geq \max(N, n_K)$:
\[|a_n - A| \leq |a_n - a_{n_K}| + |a_{n_K} - A| < \varepsilon\]
Hence $a_n \to A$.
\end{solution}

\begin{solution}[3.5]
Expanding:
\begin{align*}
\sum_{i,j}(a_i b_j - a_j b_i)^2 &= \sum_{i,j}(a_i^2 b_j^2 - 2a_i b_j a_j b_i + a_j^2 b_i^2) \\
&= \sum_i a_i^2 \sum_j b_j^2 + \sum_j a_j^2 \sum_i b_i^2 - 2\sum_i a_i b_i \sum_j a_j b_j \\
&= 2\|a\|^2\|b\|^2 - 2(a \cdot b)^2
\end{align*}
Hence $\frac{1}{2}\sum_{i,j}(a_i b_j - a_j b_i)^2 = \|a\|^2\|b\|^2 - (a \cdot b)^2 \geq 0$.

This gives $(a \cdot b)^2 \leq \|a\|^2\|b\|^2$, so $|a \cdot b| \leq \|a\|\|b\|$ (Cauchy-Schwarz).
\end{solution}

\begin{solution}[3.7]
$(AB)^T_{ij} = (AB)_{ji} = \sum_k a_{jk} b_{ki}$.

$(B^T A^T)_{ij} = \sum_k (B^T)_{ik}(A^T)_{kj} = \sum_k b_{ki} a_{jk} = \sum_k a_{jk} b_{ki}$.

These are equal, so $(AB)^T = B^T A^T$.
\end{solution}

\subsection*{Week 4: Continuity and Metric Spaces}

\begin{solution}[4.1]
Define $h: [a,b] \to \mathbb{R}$ by $h(t) = (f(b) - f(a)) \cdot f(t)$. Then $h$ is continuous on $[a,b]$ and differentiable on $(a,b)$.

By MVT, $\exists c \in (a,b)$ with $h'(c) = \frac{h(b) - h(a)}{b-a}$.

Computing: $h(b) - h(a) = (f(b) - f(a)) \cdot f(b) - (f(b) - f(a)) \cdot f(a) = \|f(b) - f(a)\|^2$.

And $h'(c) = (f(b) - f(a)) \cdot f'(c)$.

So $\frac{\|f(b) - f(a)\|^2}{b-a} = (f(b) - f(a)) \cdot f'(c) \leq \|f(b) - f(a)\| \|f'(c)\|$ (Cauchy-Schwarz).

Hence $\|f(b) - f(a)\| \leq (b-a)\|f'(c)\|$.
\end{solution}

\begin{solution}[4.3]
Let $c \in f^{-1}(U)$, so $f(c) \in U$. Since $U$ is open, $\exists \varepsilon > 0$ with $B_\varepsilon(f(c)) \subseteq U$.

By continuity of $f$, $\exists \delta > 0$ with $d_X(x, c) < \delta \Rightarrow d_Y(f(x), f(c)) < \varepsilon$.

So $x \in B_\delta(c) \Rightarrow f(x) \in B_\varepsilon(f(c)) \subseteq U \Rightarrow x \in f^{-1}(U)$.

Hence $B_\delta(c) \subseteq f^{-1}(U)$, so $f^{-1}(U)$ is open.
\end{solution}

\begin{solution}[4.5]
For $\text{rank}(AB) \leq \text{rank}(A)$: Each column of $AB$ is a linear combination of columns of $A$ (with coefficients from $B$). So $\text{col}(AB) \subseteq \text{col}(A)$, hence $\text{rank}(AB) \leq \text{rank}(A)$.

For $\text{rank}(AB) \leq \text{rank}(B)$: Each row of $AB$ is a linear combination of rows of $B$. So $\text{row}(AB) \subseteq \text{row}(B)$, hence $\text{rank}(AB) \leq \text{rank}(B)$.

Therefore $\text{rank}(AB) \leq \min(\text{rank}(A), \text{rank}(B))$.
\end{solution}

\begin{solution}[4.7]
The projection of $x$ onto $W = \text{span}\{v\}$ (where $\|v\| = 1$) is $(v \cdot x)v$.

In matrix form: $P_W(x) = v(v^T x) = (vv^T)x$.

So the matrix is $P_W = vv^T$ (an $n \times n$ matrix).

The projection onto $W^\perp$ is $P_{W^\perp} = I - P_W = I - vv^T$.
\end{solution}

\subsection*{Week 5: Series and Higher-Dimensional Differentiation}

\begin{solution}[5.1]
By the ratio test hypothesis: $|a_n| \leq |a_N| \lambda^{n-N}$ for $n \geq N$.

So $\sum_{n=N}^\infty |a_n| \leq |a_N| \sum_{n=N}^\infty \lambda^{n-N} = |a_N| \cdot \frac{1}{1-\lambda}$ (geometric series with $\lambda < 1$).

Since $\sum_{n=1}^{N-1} |a_n|$ is finite, $\sum_{n=1}^\infty |a_n|$ converges. Hence $\sum a_n$ converges absolutely.
\end{solution}

\begin{solution}[5.3]
Let $S_{2n} = \sum_{k=1}^{2n} (-1)^{k+1} a_k = (a_1 - a_2) + (a_3 - a_4) + \cdots + (a_{2n-1} - a_{2n})$.

Each parenthesis is $\geq 0$ (since $\{a_n\}$ is decreasing), so $S_{2n}$ is increasing.

Also $S_{2n} = a_1 - (a_2 - a_3) - (a_4 - a_5) - \cdots \leq a_1$, so $S_{2n}$ is bounded.

Hence $S_{2n} \to L$ for some $L$.

Now $S_{2n+1} = S_{2n} + a_{2n+1} \to L + 0 = L$ (since $a_n \to 0$).

So $S_n \to L$.
\end{solution}

\begin{solution}[5.5]
Since $\sum \|a_k\|$ converges, each component $\sum |a_{k,i}|$ converges (as $|a_{k,i}| \leq \|a_k\|$).

For real series: if $\sum |b_k|$ converges, then $\sum b_k$ converges (define $b_k^+ = \max(b_k, 0)$, $b_k^- = \max(-b_k, 0)$; both are bounded by $|b_k|$).

So each component $\sum a_{k,i}$ converges. Hence $\sum a_k$ converges in $\mathbb{R}^n$.
\end{solution}

\begin{solution}[5.7]
We compute $\det$ using row/column operations, tracking sign changes.

After reducing to upper triangular form (work omitted), we find $\det = 15$.

Alternatively: The determinant equals $\prod_{1 \leq i < j \leq n}(x_j - x_i)$ for the Vandermonde matrix.
\end{solution}

%%%%%%%%%%%%%%%%%%%%%%%%%%%%%%%%%%%%%%%%%%%%%%%%%%%%%%%%%%%%%%%%%%%%%%%%%%%%%%%
% PART II: MATH 62CM - DIFFERENTIAL FORMS AND MANIFOLDS
%%%%%%%%%%%%%%%%%%%%%%%%%%%%%%%%%%%%%%%%%%%%%%%%%%%%%%%%%%%%%%%%%%%%%%%%%%%%%%%

\mainmatter % Restore normal chapter numbering after backmatter

% Change chapter format for Part II (from "Week X" to "Chapter X")
\titleformat{\chapter}[display]
  {\sffamily\bfseries\color{MagicTeal}}
  {\flushright\fontsize{40}{40}\selectfont\color{MagicPink!50}Chapter \thechapter} 
  {-10pt}
  {\Huge}
  [\vspace{1ex}{\color{MagicMint}\titlerule[3pt]}]

\part{Math 62CM: Differential Forms and Manifolds}

% Reset chapter counter for Part II
\setcounter{chapter}{0}

\chapter{Dual Spaces and Topology}

\section{Dual Spaces}

The highlight of this course is \textbf{Stokes' theorem}, which unifies Green's theorem, the divergence theorem, and the classical Stokes theorem into a single elegant statement. To develop this, we need multilinear algebra and the theory of differential forms.

\begin{defn}{Dual Space}{dual_space}
Let $V$ be a real vector space. A function $f: V \to \mathbb{R}$ is called \textbf{linear} if $f(\lambda v) = \lambda f(v)$ and $f(v_1 + v_2) = f(v_1) + f(v_2)$ for all $\lambda \in \mathbb{R}$ and $v, v_1, v_2 \in V$.

The set of all linear functions on $V$ forms a vector space called the \textbf{dual space}, denoted $V^*$.
\end{defn}

\begin{thm}{Dimension of Dual Space}{dim_dual}
If $\dim(V) = n$ is finite, then $\dim(V^*) = n$.
\end{thm}

\begin{proof}
Let $\{v_1, \ldots, v_n\}$ be a basis for $V$. Define $v_i^*: V \to \mathbb{R}$ by $v_i^*(v_j) = \delta_{ij}$ (the Kronecker delta). We claim $\{v_1^*, \ldots, v_n^*\}$ is a basis for $V^*$.

\textbf{Linear independence:} If $\sum a_i v_i^* = 0$, then applying to $v_j$ gives $a_j = 0$ for all $j$.

\textbf{Spanning:} Given $f \in V^*$, let $a_i = f(v_i)$. Then $(f - \sum a_i v_i^*)(v_j) = f(v_j) - a_j = 0$, so $f = \sum a_i v_i^*$.
\end{proof}

\begin{defn}{Dual Basis}{dual_basis}
The basis $\{v_1^*, \ldots, v_n^*\}$ constructed above is called the \textbf{dual basis} to $\{v_1, \ldots, v_n\}$.
\end{defn}

\begin{remark}
While $V$ and $V^*$ are isomorphic (they have the same dimension), there is no \emph{canonical} isomorphism---it depends on a choice of basis. However, if $V$ has an inner product, we get a canonical isomorphism $D: V \to V^*$ via $D(v) = \ell_v$ where $\ell_v(w) = \langle v, w \rangle$.
\end{remark}

\begin{thm}{Double Dual}{double_dual}
$(V^*)^*$ is canonically isomorphic to $V$ (assuming $\dim V$ is finite).
\end{thm}

\begin{proof}
Define $\iota: V \to (V^*)^*$ by $\iota(v)(f) = f(v)$ for $f \in V^*$. This is linear and injective (if $f(v) = 0$ for all $f$, then $v = 0$). Since $\dim V = \dim V^* = \dim (V^*)^*$, it's an isomorphism.
\end{proof}

\begin{defn}{Dual Map}{dual_map}
If $A: V \to W$ is linear, the \textbf{dual map} $A^*: W^* \to V^*$ is defined by $(A^*f)(v) = f(Av)$.
\end{defn}

In matrix form, if $A$ is represented by matrix $M$, then $A^*$ is represented by $M^T$.

\begin{thm}{Properties of Dual Map}{dual_map_props}
Let $A: V \to W$ and $B: W \to U$ be linear maps. Then:
\begin{enumerate}
    \item $(B \circ A)^* = A^* \circ B^*$
    \item $(A^*)^* = A$ under the canonical identification $V \cong (V^*)^*$
    \item If $V$ has an inner product and $A: V \to V$, the \textbf{adjoint} $\tilde{A}$ satisfies $\langle Av, w \rangle = \langle v, \tilde{A}w \rangle$. In an orthonormal basis, $\tilde{A} = A^T$.
\end{enumerate}
\end{thm}

\begin{proof}
(1) For $f \in U^*$ and $v \in V$:
\[((B \circ A)^* f)(v) = f(BAv) = (B^* f)(Av) = (A^*(B^* f))(v)\]

(2) The canonical map $\iota: V \to (V^*)^*$ satisfies $\iota(v)(f) = f(v)$. We need $(A^*)^* \circ \iota = \iota \circ A$:
\[((A^*)^* \iota(v))(g) = \iota(v)(A^* g) = (A^* g)(v) = g(Av) = \iota(Av)(g)\]

(3) Using the isomorphism $D: V \to V^*$ with $D(v)(w) = \langle v, w \rangle$:
\[\langle Av, w \rangle = D(Av)(w) = (D \circ A)(v)(w)\]
We want this to equal $\langle v, \tilde{A}w \rangle = D(v)(\tilde{A}w) = (A^* \circ D)(w)(v)$.
\end{proof}

\section{Algebraic Constructions: Quotient Spaces and Direct Sums}

These constructions are fundamental to everything that follows. They give us canonical ways to build new vector spaces from old ones.

\subsection{Quotient Spaces}

\begin{defn}{Quotient Space}{quotient_space}
Let $V$ be a vector space and $L \subseteq V$ a linear subspace. For any $a \in V$, define the \textbf{affine subspace} (or \textbf{coset}):
\[L_a = a + L = \{a + x : x \in L\}\]

Two cosets $L_a$ and $L_b$ are either identical (if $a - b \in L$) or disjoint (if $a - b \notin L$).

The \textbf{quotient space} $V/L$ is the set of all cosets, with operations:
\[L_a + L_b := L_{a+b}, \quad \lambda L_a := L_{\lambda a}\]
\end{defn}

% TikZ diagram for quotient space
\begin{center}
\begin{tikzpicture}[scale=0.8]
    % Draw the subspace L
    \draw[thick, blue] (-2,-1) -- (2,1) node[right] {$L$};
    % Draw parallel cosets
    \draw[dashed, blue!60] (-2,0) -- (2,2) node[right] {$L_a$};
    \draw[dashed, blue!60] (-2,1) -- (2,3) node[right] {$L_b$};
    \draw[dashed, blue!60] (-2,-2) -- (2,0);
    % Mark points
    \fill (0,0) circle (2pt) node[below right] {$0$};
    \fill (0,1) circle (2pt) node[below right] {$a$};
    \fill (0,2) circle (2pt) node[below right] {$b$};
    % Label
    \node at (3.5,1) {$V/L = $ cosets};
\end{tikzpicture}
\end{center}

\begin{thm}{Well-Definedness of Quotient Operations}{quotient_welldef}
The operations on $V/L$ are well-defined, i.e., they don't depend on the choice of representative.
\end{thm}

\begin{proof}
Suppose $L_a = L_{a'}$ and $L_b = L_{b'}$, meaning $a - a' \in L$ and $b - b' \in L$.

For addition: $(a + b) - (a' + b') = (a - a') + (b - b') \in L$, so $L_{a+b} = L_{a'+b'}$.

For scalar multiplication: $\lambda a - \lambda a' = \lambda(a - a') \in L$, so $L_{\lambda a} = L_{\lambda a'}$.
\end{proof}

\begin{thm}{Dimension of Quotient Space}{dim_quotient}
$\dim(V/L) = \dim(V) - \dim(L) = \text{codim}_V(L)$.
\end{thm}

\begin{proof}
Let $N \subseteq V$ be any subspace complementary to $L$, meaning $L \cap N = \{0\}$ and $\dim L + \dim N = \dim V$.

Define $\phi: N \to V/L$ by $\phi(x) = L_x$.

\textbf{Injective:} If $L_x = L_y$ for $x, y \in N$, then $x - y \in L \cap N = \{0\}$, so $x = y$.

\textbf{Surjective:} Any coset $L_a$ contains a unique element of $N$. Write $a = \ell + n$ where $\ell \in L$ and $n \in N$. Then $L_a = L_n = \phi(n)$.

So $\phi$ is an isomorphism, and $\dim(V/L) = \dim N = \dim V - \dim L$.
\end{proof}

\begin{defn}{Canonical Projection}{canonical_proj}
The \textbf{canonical projection} $\pi_L: V \to V/L$ is defined by $\pi_L(v) = L_v$.

This is a surjective linear map with $\ker(\pi_L) = L$.
\end{defn}

\begin{thm}{Universal Property of Quotient}{quotient_universal}
Let $f: V \to W$ be a linear map with $L \subseteq \ker(f)$. Then there exists a unique linear map $\bar{f}: V/L \to W$ such that $f = \bar{f} \circ \pi_L$.

% Commutative diagram
\begin{center}
\begin{tikzpicture}
    \node (V) at (0,1.5) {$V$};
    \node (VL) at (0,0) {$V/L$};
    \node (W) at (2,0) {$W$};
    \draw[->] (V) -- node[left] {$\pi_L$} (VL);
    \draw[->] (V) -- node[above right] {$f$} (W);
    \draw[->, dashed] (VL) -- node[below] {$\exists!\bar{f}$} (W);
\end{tikzpicture}
\end{center}
\end{thm}

\begin{proof}
Define $\bar{f}(L_a) = f(a)$. This is well-defined: if $L_a = L_b$, then $a - b \in L \subseteq \ker(f)$, so $f(a) = f(b)$.

Linearity: $\bar{f}(L_a + L_b) = \bar{f}(L_{a+b}) = f(a+b) = f(a) + f(b) = \bar{f}(L_a) + \bar{f}(L_b)$.

Uniqueness: If $g: V/L \to W$ satisfies $f = g \circ \pi_L$, then $g(L_a) = g(\pi_L(a)) = f(a) = \bar{f}(L_a)$.
\end{proof}

\begin{cor}{First Isomorphism Theorem}{first_iso_thm}
For any surjective linear map $f: V \to W$:
\[V/\ker(f) \cong W\]
The isomorphism is $\bar{f}: L_a \mapsto f(a)$.
\end{cor}

\subsection{Direct Sums}

\begin{defn}{External Direct Sum}{direct_sum_ext}
Given vector spaces $V$ and $W$, their \textbf{(external) direct sum} $V \oplus W$ is the set of pairs $(v, w)$ with $v \in V$, $w \in W$, with operations:
\[(v_1, w_1) + (v_2, w_2) = (v_1 + v_2, w_1 + w_2), \quad \lambda(v, w) = (\lambda v, \lambda w)\]
\end{defn}

\begin{thm}{Dimension of Direct Sum}{dim_direct_sum}
$\dim(V \oplus W) = \dim(V) + \dim(W)$.

If $\{v_1, \ldots, v_n\}$ is a basis for $V$ and $\{w_1, \ldots, w_m\}$ is a basis for $W$, then
\[\{(v_1, 0), \ldots, (v_n, 0), (0, w_1), \ldots, (0, w_m)\}\]
is a basis for $V \oplus W$.
\end{thm}

\begin{defn}{Internal Direct Sum}{direct_sum_int}
Subspaces $V, W \subseteq H$ form an \textbf{internal direct sum} $H = V \oplus W$ if:
\begin{enumerate}
    \item $V \cap W = \{0\}$
    \item $V + W = H$ (i.e., every $h \in H$ can be written as $h = v + w$)
\end{enumerate}
Equivalently, every $h \in H$ can be written \emph{uniquely} as $h = v + w$ with $v \in V$, $w \in W$.
\end{defn}

\begin{thm}{Projection Maps}{projection_maps}
If $H = V \oplus W$ (internal), there are canonical \textbf{projection maps}:
\[\pi_V: H \to V, \quad \pi_W: H \to W\]
defined by $\pi_V(v + w) = v$ and $\pi_W(v + w) = w$.

These satisfy:
\begin{enumerate}
    \item $\pi_V + \pi_W = \text{id}_H$
    \item $\pi_V^2 = \pi_V$, $\pi_W^2 = \pi_W$ (idempotent)
    \item $\pi_V \circ \pi_W = \pi_W \circ \pi_V = 0$
    \item $\ker(\pi_V) = W$, $\ker(\pi_W) = V$
\end{enumerate}
\end{thm}

\begin{proof}
Well-definedness follows from uniqueness of the decomposition.

(1) $(\pi_V + \pi_W)(v + w) = v + w$.

(2) $\pi_V(v + w) = v = v + 0$, so $\pi_V(\pi_V(v+w)) = \pi_V(v) = v$.

(3) $\pi_V(\pi_W(v+w)) = \pi_V(w) = \pi_V(0 + w) = 0$.

(4) $\pi_V(v+w) = 0 \iff v = 0 \iff v + w \in W$.
\end{proof}

\begin{remark}
The external and internal direct sums are related: if $H = V \oplus W$ internally, the map $(v, w) \mapsto v + w$ is an isomorphism from the external $V \oplus W$ to $H$.
\end{remark}

\subsection{Free Vector Spaces}

\begin{defn}{Free Vector Space}{free_vs}
Given any set $S$, the \textbf{free vector space} on $S$, denoted $\mathbb{R}[S]$ or $\text{Free}(S)$, is the vector space of all formal finite linear combinations:
\[\sum_{s \in S} a_s \cdot s, \quad a_s \in \mathbb{R}, \text{ only finitely many } a_s \neq 0\]
\end{defn}

\begin{thm}{Universal Property of Free Space}{free_universal}
For any vector space $V$ and any function $f: S \to V$, there exists a unique linear map $\tilde{f}: \mathbb{R}[S] \to V$ extending $f$.

\begin{center}
\begin{tikzpicture}
    \node (S) at (0,1.5) {$S$};
    \node (RS) at (0,0) {$\mathbb{R}[S]$};
    \node (V) at (2,0) {$V$};
    \draw[right hook->] (S) -- node[left] {$\iota$} (RS);
    \draw[->] (S) -- node[above right] {$f$} (V);
    \draw[->, dashed] (RS) -- node[below] {$\exists!\tilde{f}$} (V);
\end{tikzpicture}
\end{center}
\end{thm}

\begin{proof}
Define $\tilde{f}(\sum a_s \cdot s) = \sum a_s f(s)$. This is clearly linear and extends $f$ (i.e., $\tilde{f}(s) = f(s)$).

Uniqueness: Any linear extension must satisfy $\tilde{f}(\sum a_s \cdot s) = \sum a_s \tilde{f}(s) = \sum a_s f(s)$.
\end{proof}

\begin{remark}
$\dim(\mathbb{R}[S]) = |S|$ when $S$ is finite. The elements of $S$ form a basis for $\mathbb{R}[S]$.
\end{remark}

\section{Topology in Vector Spaces}

Let $V$ be a finite-dimensional Euclidean space (vector space with inner product).

\begin{defn}{Open and Closed Sets}{open_closed}
\begin{itemize}
    \item The \textbf{open ball} of radius $r$ centered at $p$ is $B_r(p) = \{x \in V : \|x - p\| < r\}$.
    \item The \textbf{closed ball} is $\overline{B}_r(p) = \{x \in V : \|x - p\| \leq r\}$.
    \item The \textbf{sphere} is $S_r(p) = \{x \in V : \|x - p\| = r\}$.
    \item A set $U \subseteq V$ is \textbf{open} if for every $x \in U$, there exists $\delta > 0$ such that $B_\delta(x) \subseteq U$.
    \item A set $A \subseteq V$ is \textbf{closed} if $V \setminus A$ is open.
\end{itemize}
\end{defn}

\begin{thm}{Properties of Open and Closed Sets}{open_closed_props}
\begin{enumerate}
    \item Arbitrary unions of open sets are open.
    \item Finite intersections of open sets are open.
    \item Arbitrary intersections of closed sets are closed.
    \item Finite unions of closed sets are closed.
\end{enumerate}
\end{thm}

\begin{defn}{Interior, Closure, Boundary}{int_clos_bdy}
For $X \subseteq V$:
\begin{itemize}
    \item The \textbf{interior} $\text{Int}(X)$ is the union of all open sets contained in $X$.
    \item The \textbf{closure} $\overline{X}$ is the intersection of all closed sets containing $X$.
    \item The \textbf{boundary} $\partial X = \overline{X} \setminus \text{Int}(X)$.
\end{itemize}
\end{defn}

\begin{defn}{Compact Sets}{compact}
A set $A \subseteq V$ is \textbf{compact} if:
\begin{enumerate}
    \item (Definition 1) $A$ is closed and bounded.
    \item (Definition 2) Every sequence $\{x_n\} \subseteq A$ has a subsequence converging to a point in $A$.
    \item (Definition 3) Every open cover of $A$ has a finite subcover.
\end{enumerate}
These definitions are equivalent in $\mathbb{R}^n$ (Heine-Borel theorem).
\end{defn}

\section{Connected and Path-Connected Sets}

\begin{defn}{Connected}{connected}
A set $A \subseteq V$ is \textbf{connected} if we cannot write $A = A_0 \cup A_1$ where $A_0, A_1$ are both non-empty, disjoint, and relatively open in $A$.
\end{defn}

\begin{defn}{Path-Connected}{path_connected}
A set $A$ is \textbf{path-connected} if for any $a, b \in A$, there exists a continuous path $\gamma: [0,1] \to A$ with $\gamma(0) = a$ and $\gamma(1) = b$.
\end{defn}

\begin{thm}{Path-Connected Implies Connected}{path_conn_implies}
\end{thm}

\begin{proof}
Suppose $A$ is path-connected but not connected. Write $A = A_0 \cup A_1$ with $A_0, A_1$ open, disjoint, non-empty. Pick $a_0 \in A_0, a_1 \in A_1$ and a path $\gamma: [0,1] \to A$ connecting them. Then $f = \mathbf{1}_{A_1} \circ \gamma: [0,1] \to \{0,1\}$ is continuous with $f(0) = 0, f(1) = 1$. But this contradicts the Intermediate Value Theorem.
\end{proof}

\begin{remark}
The converse is false. The ``topologist's sine curve'' $\{(x, \sin(1/x)) : x > 0\} \cup \{(0,0)\}$ is connected but not path-connected.
\end{remark}

\begin{thm}{$\mathbb{R}^n$ is Connected}{rn_connected}
$\mathbb{R}^n$ is connected (in fact, path-connected via linear paths $\gamma(t) = (1-t)x + ty$).
\end{thm}

\begin{defn}{Homeomorphism}{homeomorphism}
$f: A \to B$ is a \textbf{homeomorphism} if $f$ is continuous, bijective, and $f^{-1}$ is continuous. We say $A$ and $B$ are \textbf{homeomorphic} (``topologically the same'').
\end{defn}

\begin{thm}{Continuous Image of Connected}{cont_connected}
Continuous images of connected sets are connected.
\end{thm}

\begin{exercises}

\begin{exercise}
Prove the converse: given any basis $\ell_1, \ldots, \ell_n$ of $V^*$, construct a dual basis $w_1, \ldots, w_n$ of $V$ so that $\ell_i(w_j) = \delta_{ij}$.
\end{exercise}

\begin{exercise}
Let $V$ be a vector space with inner product $\langle \cdot, \cdot \rangle$. Prove that $\mathcal{D}: V \to V^*$, $\mathcal{D}(v) = \ell_v$ where $\ell_v(x) = \langle v, x \rangle$, is an isomorphism. Show that $\mathcal{D} = i_S$ for any orthonormal basis $S$.
\end{exercise}

\begin{exercise}
Show that both maps $i_S$ and $\mathcal{D}$ are injective in the infinite-dimensional case as well.
\end{exercise}

\begin{exercise}
Prove that the definition of open and closed sets is independent of the choice of inner product on a finite-dimensional vector space.
\end{exercise}

\begin{exercise}
Show that the closure of a connected set is connected.
\end{exercise}

\begin{exercise}
Prove that the closure of the graph of $\sin(1/x)$, $x > 0$, is connected but not path-connected.
\end{exercise}

\begin{exercise}
\begin{enumerate}
    \item Prove that $[0,1]$, $[0,1)$, and $(0,1)$ are pairwise not homeomorphic.
    \item Prove that $\mathbb{R}^n$ is not homeomorphic to $\mathbb{R}$ if $n > 1$.
\end{enumerate}
\end{exercise}

\end{exercises}

\chapter{Differential Forms}

\section{Vector Fields and 1-Forms}

\begin{defn}{Vector Field}{vector_field}
A \textbf{vector field} on an open set $U \subseteq V$ is an assignment $x \mapsto v_x \in V_x$ (where $V_x$ is a copy of $V$ at $x$).
\end{defn}

\begin{defn}{Differential 1-Form}{diff_1form}
A \textbf{differential 1-form} $\alpha$ on $U$ is an assignment $x \mapsto \alpha_x \in V_x^*$ varying smoothly.
\end{defn}

In coordinates $(x_1, \ldots, x_n)$, a 1-form is $\alpha = \sum a_i(x) dx_i$ where $dx_i$ are the coordinate differentials (dual to the standard basis).

\begin{defn}{Differential of a Function}{diff_func}
For $f: U \to \mathbb{R}$ differentiable, the \textbf{differential} $df$ is a 1-form defined by:
\[
df_x(v) = \lim_{t \to 0} \frac{f(x + tv) - f(x)}{t} = D_v f(x)
\]
In coordinates: $df = \sum \frac{\partial f}{\partial x_i} dx_i$.
\end{defn}

\begin{defn}{Exact and Closed Forms}{exact_closed}
A 1-form $\alpha$ is:
\begin{itemize}
    \item \textbf{Exact} if $\alpha = df$ for some function $f$ (called a \textbf{primitive}).
    \item \textbf{Closed} if $\frac{\partial a_i}{\partial x_j} = \frac{\partial a_j}{\partial x_i}$ for all $i, j$.
\end{itemize}
\end{defn}

\begin{thm}{Exact Implies Closed}{exact_implies_closed}
Exact $\Rightarrow$ Closed.
\end{thm}

\section{Pullback}

\begin{defn}{Pullback}{pullback}
Let $f: U \to U'$ be differentiable and $\alpha$ a 1-form on $U'$. The \textbf{pullback} $f^*\alpha$ is a 1-form on $U$ defined by:
\[
(f^*\alpha)_x(v) = \alpha_{f(x)}((d_x f)(v))
\]
\end{defn}

\begin{thm}{Properties of Pullback}{pullback_props}
\begin{enumerate}
    \item $f^*(\alpha + \beta) = f^*\alpha + f^*\beta$
    \item $f^*(\phi \alpha) = (\phi \circ f) f^*\alpha$ for functions $\phi$
    \item $f^*(dg) = d(g \circ f)$
    \item $(g \circ f)^* = f^* \circ g^*$
\end{enumerate}
\end{thm}

\begin{example}
Let $f(r, \theta) = (r\cos\theta, r\sin\theta)$ (polar coordinates) and $\alpha = x dy - y dx$ on $\mathbb{R}^2$. Then:
\[
f^*\alpha = r\cos\theta \cdot d(r\sin\theta) - r\sin\theta \cdot d(r\cos\theta) = r^2 d\theta
\]
\end{example}

\section{Integration on Curves}

\begin{defn}{Line Integral}{line_integral}
Let $\gamma: [a,b] \to U$ be a smooth path and $\alpha$ a 1-form on $U$. The \textbf{line integral} is:
\[
\int_\gamma \alpha = \int_a^b \gamma^*\alpha = \int_a^b \alpha_{\gamma(t)}(\gamma'(t)) dt
\]
\end{defn}

\begin{thm}{Independence of Parametrization}{param_indep}
The line integral is independent of parametrization (up to sign for orientation reversal).
\end{thm}

\begin{thm}{Fundamental Theorem for Exact Forms}{ftc_exact}
If $\alpha = df$ is exact and $\gamma$ is a path from $A$ to $B$:
\[
\int_\gamma df = f(B) - f(A)
\]
In particular, $\oint_\gamma df = 0$ for any closed loop.
\end{thm}

\section{Change of Variables (Jacobian)}

This is one of the most important applications of differential forms to multivariable calculus.

\begin{thm}{Pullback of Volume Form}{pullback_vol}
Let $\phi: U \to V$ be a smooth map between open sets in $\mathbb{R}^n$. For the volume form $\omega = dy_1 \wedge \cdots \wedge dy_n$ on $V$:
\[\phi^* \omega = \det(D\phi(x)) \, dx_1 \wedge \cdots \wedge dx_n\]
The factor $\det(D\phi)$ is the \textbf{Jacobian determinant}.
\end{thm}

\begin{proof}
Write $\phi = (\phi_1, \ldots, \phi_n)$. Then:
\[\phi^*(dy_j) = d(\phi_j) = \sum_i \frac{\partial \phi_j}{\partial x_i} dx_i\]
So:
\begin{align*}
\phi^*(dy_1 \wedge \cdots \wedge dy_n) &= d\phi_1 \wedge \cdots \wedge d\phi_n \\
&= \left(\sum_i \frac{\partial \phi_1}{\partial x_i} dx_i\right) \wedge \cdots \wedge \left(\sum_i \frac{\partial \phi_n}{\partial x_i} dx_i\right)
\end{align*}
The only surviving term is when all indices are distinct, giving $\det(D\phi) \, dx_1 \wedge \cdots \wedge dx_n$.
\end{proof}

\begin{thm}{Change of Variables Formula}{change_var}
Let $\phi: U \to V$ be a diffeomorphism. For any integrable function $f$ on $V$:
\[\int_V f(y)\, dy = \int_U f(\phi(x)) \, |\det D\phi(x)| \, dx\]
\end{thm}

\begin{remark}
The absolute value is needed because orientation-reversing maps give negative Jacobians, but volumes are always positive. If we work with \emph{signed} integrals (oriented integration), we can drop the absolute value.
\end{remark}

\begin{example}
\textbf{Polar coordinates:} $\phi(r, \theta) = (r\cos\theta, r\sin\theta)$. The Jacobian is:
\[\det \begin{pmatrix} \cos\theta & -r\sin\theta \\ \sin\theta & r\cos\theta \end{pmatrix} = r\]
So $dx\,dy = r\,dr\,d\theta$.
\end{example}

\section{Poincaré Lemma}

\begin{defn}{Star-Shaped}{star_shaped}
A set $U$ is \textbf{star-shaped} with respect to $a \in U$ if for every $x \in U$ and $t \in [0,1]$, $(1-t)a + tx \in U$.
\end{defn}

\begin{thm}{Poincaré Lemma}{poincare_1form}
On a star-shaped open set $U$ (star-shaped with respect to origin), every closed 1-form is exact.
\end{thm}

\begin{proof}
Let $\alpha = \sum_{i=1}^n a_i(x) dx_i$ be closed, meaning $\frac{\partial a_i}{\partial x_j} = \frac{\partial a_j}{\partial x_i}$ for all $i, j$.

Define $f: U \to \mathbb{R}$ by:
\[f(x) = \int_0^1 \sum_{i=1}^n a_i(tx) x_i \, dt\]

\textbf{Claim:} $df = \alpha$, i.e., $\frac{\partial f}{\partial x_j} = a_j(x)$ for all $j$.

Differentiating under the integral:
\begin{align*}
\frac{\partial f}{\partial x_j} &= \int_0^1 \left( \sum_{i=1}^n t \frac{\partial a_i}{\partial x_j}(tx) x_i + a_j(tx) \right) dt
\end{align*}

Using the closedness condition $\frac{\partial a_i}{\partial x_j} = \frac{\partial a_j}{\partial x_i}$:
\begin{align*}
\frac{\partial f}{\partial x_j} &= \int_0^1 \left( \sum_{i=1}^n t \frac{\partial a_j}{\partial x_i}(tx) x_i + a_j(tx) \right) dt \\
&= \int_0^1 \frac{d}{dt}\left( t \cdot a_j(tx) \right) dt \\
&= [t \cdot a_j(tx)]_0^1 = a_j(x)
\end{align*}
where we used the chain rule: $\frac{d}{dt}(a_j(tx)) = \sum_i \frac{\partial a_j}{\partial x_i}(tx) \cdot x_i$.
\end{proof}

\begin{thm}{Poincaré Lemma (General)}{poincare_general}
On a star-shaped open set $U$ (star-shaped with respect to the origin), every closed $k$-form ($k \geq 1$) is exact.
\end{thm}

\begin{proof}
We define the \textbf{homotopy operator} $H: \Omega^k(U) \to \Omega^{k-1}(U)$ as follows.

For a $k$-form $\omega = \sum_{I} a_I(x) dx_{i_1} \wedge \cdots \wedge dx_{i_k}$ (sum over multi-indices $I = (i_1, \ldots, i_k)$ with $i_1 < \cdots < i_k$), define:
\[H\omega = \sum_I \sum_{j=1}^k (-1)^{j-1} \left( \int_0^1 t^{k-1} a_I(tx) \, dt \right) x_{i_j} \, dx_{i_1} \wedge \cdots \wedge \widehat{dx_{i_j}} \wedge \cdots \wedge dx_{i_k}\]
where $\widehat{dx_{i_j}}$ means that term is omitted.

\textbf{Key identity:} For any $k$-form $\omega$:
\[dH\omega + Hd\omega = \omega\]

To verify this, let's check on a monomial $\omega = a(x) dx_{i_1} \wedge \cdots \wedge dx_{i_k}$. A computation shows that when we expand $dH\omega + Hd\omega$ and collect terms, everything cancels except $\omega$ itself. The key step uses:
\[\frac{d}{dt}[t^k a(tx)] = k t^{k-1} a(tx) + t^k \sum_i \frac{\partial a}{\partial x_i}(tx) x_i\]
and integrating from 0 to 1 gives $a(x)$.

\textbf{Application:} If $d\omega = 0$, then:
\[\omega = dH\omega + Hd\omega = dH\omega + H(0) = d(H\omega)\]
So $\omega$ is exact with primitive $H\omega$.
\end{proof}

\begin{remark}
The operator $H$ is sometimes called the \textbf{cone operator} because it ``pulls back'' forms along the rays from the origin. The identity $dH + Hd = \text{id}$ is a \textbf{chain homotopy} showing that $d$ is ``invertible up to exact forms'' on star-shaped domains.
\end{remark}

\begin{defn}{Simply Connected}{simply_connected}
$U$ is \textbf{simply connected} if it is path-connected and every two paths connecting the same endpoints are homotopic.
\end{defn}

\begin{thm}{Simply Connected Poincaré}{poincare_sc}
On a simply connected open set, every closed form is exact.
\end{thm}

\begin{example}
The form $\alpha = \frac{-y dx + x dy}{x^2 + y^2}$ on $\mathbb{R}^2 \setminus \{0\}$ is closed but not exact (since $\oint_{S^1} \alpha = 2\pi \neq 0$). This shows $\mathbb{R}^2 \setminus \{0\}$ is not simply connected.
\end{example}

\begin{example}
\textbf{(Pullback Computation)}
Let $f: \mathbb{R}^2 \to \mathbb{R}^3$, $f(u,v) = (u^2, uv, v^2)$ and $\omega = x\,dy \wedge dz + y\,dz \wedge dx$. Then:
\begin{align*}
f^*\omega &= u^2 \cdot d(uv) \wedge d(v^2) + uv \cdot d(v^2) \wedge d(u^2) \\
&= u^2 (v\,du + u\,dv) \wedge (2v\,dv) + uv(2v\,dv) \wedge (2u\,du) \\
&= u^2 \cdot 2v^2\,du \wedge dv + uv \cdot (-4uv)\,du \wedge dv \\
&= (2u^2v^2 - 4u^2v^2)\,du \wedge dv = -2u^2v^2\,du \wedge dv
\end{align*}
\end{example}

\begin{example}
\textbf{(Finding a Primitive)}
Is $\omega = (2xy + z^2)\,dx + x^2\,dy + 2xz\,dz$ exact on $\mathbb{R}^3$?

Check: $\frac{\partial}{\partial y}(2xy + z^2) = 2x = \frac{\partial}{\partial x}(x^2)$ \checkmark

$\frac{\partial}{\partial z}(2xy + z^2) = 2z = \frac{\partial}{\partial x}(2xz)$ \checkmark

$\frac{\partial}{\partial z}(x^2) = 0 = \frac{\partial}{\partial y}(2xz)$ \checkmark

So $d\omega = 0$. Since $\mathbb{R}^3$ is star-shaped, $\omega = df$ for some $f$.

Finding $f$: We need $\frac{\partial f}{\partial x} = 2xy + z^2$, $\frac{\partial f}{\partial y} = x^2$, $\frac{\partial f}{\partial z} = 2xz$.

From the second: $f = x^2 y + g(x,z)$.

Substituting into the first: $2xy + \frac{\partial g}{\partial x} = 2xy + z^2 \Rightarrow g = xz^2 + h(z)$.

Check with the third: $\frac{\partial}{\partial z}(x^2 y + xz^2 + h(z)) = 2xz + h'(z) = 2xz \Rightarrow h = C$.

So $f(x,y,z) = x^2 y + xz^2 + C$.
\end{example}

\begin{exercises}

\begin{exercise}
Let $f: \mathbb{R}^2 \to \mathbb{R}^3$, $f(x,y) = (e^x - 2y, x + y^2, \sin(y) + 1)$ and $\alpha = z\,dx - 3\,dy + e^y\,dz$. Compute $f^*\alpha$.
\end{exercise}

\begin{exercise}
In $\mathbb{R}^4$ with coordinates $(x,y,z,u)$, consider the plane field defined by $dz - y\,dx = 0$ and $dy - u\,dx = 0$. For a curve $\Gamma$ given by $y = f(x)$, $z = g(x)$, $u = h(x)$, find necessary and sufficient conditions for $\Gamma$ to be tangent to this plane field.
\end{exercise}

\begin{exercise}
In $\mathbb{R}^3$, find a curve $\Gamma$ whose projection to the $(x,y)$-plane is $x = y^2$, tangent to $dz - y\,dx = 0$, and containing the origin.
\end{exercise}

\begin{exercise}
Find a primitive in $\mathbb{R}^2$ of the exact 1-form:
\[
\alpha = e^x(e^y(x-y+2) + y)dx + e^x(e^y(x-y) + 1)dy
\]
\end{exercise}

\begin{exercise}
Prove that $\mathbb{R}^3 \setminus \{y = x^2, z = 0\}$ is not simply connected.
\end{exercise}

\begin{exercise}
Show that any star-shaped open set is simply connected.
\end{exercise}

\end{exercises}

\chapter{Multilinear Algebra}

\section{Bilinear and Multilinear Forms}

\begin{defn}{Multilinear Form}{multilinear}
A function $f: V \times \cdots \times V \to \mathbb{R}$ ($k$ copies) is \textbf{$k$-linear} (or multilinear) if it is linear in each argument when the others are fixed.
\end{defn}

\begin{defn}{Bilinear Form}{bilinear}
A \textbf{bilinear form} $f: V \times V \to \mathbb{R}$ is \textbf{symmetric} if $f(x,y) = f(y,x)$ and \textbf{skew-symmetric} if $f(x,y) = -f(y,x)$.
\end{defn}

\begin{defn}{Quadratic Form}{quadratic}
A \textbf{quadratic form} is $Q(x) = f(x,x)$ for some symmetric bilinear $f$.
\end{defn}

\begin{thm}{Polarization}{polarization}
Given a quadratic form $Q$, there is a unique symmetric bilinear form $f$ with $Q(x) = f(x,x)$:
\[
f(x,y) = \frac{1}{2}(Q(x+y) - Q(x) - Q(y))
\]
\end{thm}

\begin{example}
\textbf{(Quadratic Forms on $\mathbb{R}^2$)}
$Q(x,y) = x^2 + 4xy + 3y^2$ has matrix $\begin{pmatrix} 1 & 2 \\ 2 & 3 \end{pmatrix}$.

By polarization: $f((x_1, y_1), (x_2, y_2)) = x_1 x_2 + 2(x_1 y_2 + x_2 y_1) + 3y_1 y_2$.

The eigenvalues of the matrix determine the signature: $\lambda = 2 \pm \sqrt{5}$, both positive, so this is positive definite.
\end{example}

\begin{example}
\textbf{(Skew-Symmetric Bilinear Forms)}
On $\mathbb{R}^3$, the cross product defines a skew-symmetric ``bilinear form'' to $\mathbb{R}^3$:
\[f(u, v) = u \times v\]

A truly scalar-valued example: $\omega(u, v) = \det \begin{pmatrix} u_1 & v_1 \\ u_2 & v_2 \end{pmatrix}$ on $\mathbb{R}^2$ is skew-symmetric.
\end{example}

\begin{defn}{Index and Rank}{index_rank}
For a quadratic form $Q$ written as $Q = \sum \lambda_j x_j^2$:
\begin{itemize}
    \item The number of positive $\lambda_j$ is the \textbf{positive index}.
    \item The number of negative $\lambda_j$ is the \textbf{negative index}.
    \item The rank is the total number of nonzero $\lambda_j$.
\end{itemize}
\end{defn}

\begin{thm}{Sylvester's Law of Inertia}{sylvester}
The positive index, negative index, and nullity of a quadratic form are independent of the choice of basis (i.e., the presentation as a sum of squares).
\end{thm}

\begin{proof}
Let $Q$ be a quadratic form on $V$ with associated symmetric bilinear form $f$. Suppose we have two presentations:
\[Q = \sum_{i=1}^{p} x_i^2 - \sum_{i=p+1}^{p+q} x_i^2 \quad \text{and} \quad Q = \sum_{j=1}^{p'} y_j^2 - \sum_{j=p'+1}^{p'+q'} y_j^2\]
in bases $\{e_1, \ldots, e_n\}$ and $\{f_1, \ldots, f_n\}$ respectively.

Define $V^+ = \text{span}(e_1, \ldots, e_p)$ and $V'^- = \text{span}(f_{p'+1}, \ldots, f_{p'+q'})$.

\textbf{Claim:} $p \leq p'$.

Suppose $p > p'$. Then $\dim V^+ + \dim V'^- = p + q' > p' + q' = \dim V - \text{nullity}$.

So $V^+ \cap V'^- \neq \{0\}$. Pick $v \neq 0$ in the intersection.

Since $v \in V^+$, write $v = \sum_{i=1}^p a_i e_i$, so $Q(v) = \sum_{i=1}^p a_i^2 > 0$ (since $v \neq 0$).

Since $v \in V'^-$, write $v = \sum_{j=p'+1}^{p'+q'} b_j f_j$, so $Q(v) = -\sum_{j=p'+1}^{p'+q'} b_j^2 \leq 0$.

Contradiction. Hence $p \leq p'$. By symmetry, $p' \leq p$, so $p = p'$.

Similarly, $q = q'$.
\end{proof}

\begin{remark}
The proof uses a dimension-counting argument: if we could find more positive directions in one basis than another, we'd get a vector that is both positive and non-positive, which is impossible.
\end{remark}

\section{Exterior (Wedge) Product}

\begin{defn}{Skew-Symmetric $k$-Linear Form}{skew_sym}
$f: V^k \to \mathbb{R}$ is \textbf{skew-symmetric} if swapping any two arguments negates the value:
\[
f(v_1, \ldots, v_i, \ldots, v_j, \ldots, v_k) = -f(v_1, \ldots, v_j, \ldots, v_i, \ldots, v_k)
\]
\end{defn}

The space of skew-symmetric $k$-linear forms on $V$ is denoted $\Lambda^k(V^*)$ and has dimension $\binom{n}{k}$ where $n = \dim V$.

\begin{defn}{Wedge Product}{wedge}
For $\varphi \in \Lambda^k(V^*)$ and $\psi \in \Lambda^\ell(V^*)$, the \textbf{wedge product} $\varphi \wedge \psi \in \Lambda^{k+\ell}(V^*)$ is:
\[
(\varphi \wedge \psi)(v_1, \ldots, v_{k+\ell}) = \frac{1}{k!\ell!} \sum_{\sigma \in S_{k+\ell}} \text{sign}(\sigma) \varphi(v_{\sigma(1)}, \ldots, v_{\sigma(k)}) \psi(v_{\sigma(k+1)}, \ldots, v_{\sigma(k+\ell)})
\]
\end{defn}

\begin{thm}{Properties of Wedge Product}{wedge_props}
\begin{enumerate}
    \item \textbf{Anticommutativity:} $\varphi \wedge \psi = (-1)^{k\ell} \psi \wedge \varphi$
    \item \textbf{Associativity:} $(\varphi \wedge \psi) \wedge \chi = \varphi \wedge (\psi \wedge \chi)$
    \item \textbf{Dependence criterion:} $\ell_1 \wedge \cdots \wedge \ell_k = 0 \iff \ell_1, \ldots, \ell_k$ are linearly dependent
\end{enumerate}
\end{thm}

\begin{proof}
(1) The transposition that swaps the first $k$ and last $\ell$ indices has sign $(-1)^{k\ell}$ (it's the product of $k\ell$ transpositions).

(2) Both sides equal $\frac{1}{k!\ell!m!}(\varphi \otimes \psi \otimes \chi)^{\text{asym}}$.

(3) $(\Rightarrow)$: Suppose $\ell_1 \wedge \cdots \wedge \ell_k = 0$ but the $\ell_i$ are independent. Extend to a basis $\ell_1, \ldots, \ell_n$ of $V^*$ and let $w_1, \ldots, w_n$ be the dual basis. Then:
\[(\ell_1 \wedge \cdots \wedge \ell_k)(w_1, \ldots, w_k) = \det(\ell_i(w_j))_{i,j \leq k} = \det(I_k) = 1 \neq 0\]
Contradiction.

$(\Leftarrow)$: If the $\ell_i$ are dependent, say $\ell_k = \sum_{i<k} c_i \ell_i$, then:
\[\ell_1 \wedge \cdots \wedge \ell_k = \sum_{i<k} c_i (\ell_1 \wedge \cdots \wedge \ell_{k-1} \wedge \ell_i) = 0\]
since each term has $\ell_i$ appearing twice.
\end{proof}

\begin{thm}{Wedge of 1-Forms as Determinant}{wedge_det}
For 1-forms $\ell_1, \ldots, \ell_k$ and vectors $v_1, \ldots, v_k$:
\[(\ell_1 \wedge \cdots \wedge \ell_k)(v_1, \ldots, v_k) = \det(\ell_i(v_j))\]
\end{thm}

\begin{proof}
By definition, $\ell_1 \wedge \cdots \wedge \ell_k = (\ell_1 \otimes \cdots \otimes \ell_k)^{\text{asym}}$. Evaluating:
\[(\ell_1 \wedge \cdots \wedge \ell_k)(v_1, \ldots, v_k) = \sum_{\sigma} \text{sign}(\sigma) \ell_1(v_{\sigma(1)}) \cdots \ell_k(v_{\sigma(k)})\]
which is exactly the Leibniz formula for $\det(\ell_i(v_j))$.
\end{proof}

\begin{thm}{Basis for $\Lambda^k(V^*)$}{lambda_basis}
If $x_1, \ldots, x_n$ is a basis for $V^*$, then $\{x_{i_1} \wedge \cdots \wedge x_{i_k} : i_1 < \cdots < i_k\}$ is a basis for $\Lambda^k(V^*)$.
\end{thm}

\section{Determinant as a Volume Form}

\begin{thm}{Determinant via Exterior Power}{det_ext}
For a linear map $A: V \to V$, the induced map $A^*: \Lambda^n(V^*) \to \Lambda^n(V^*)$ is multiplication by $\det(A)$.
\end{thm}

This gives an intrinsic definition of the determinant: it's the scalar by which $A$ scales the ``top'' exterior power.

\textbf{Key idea:} $\det(A)$ is a \emph{signed volume}. The sign captures orientation!

\section{Orientation of Vector Spaces}

\begin{defn}{Orientation}{orientation_vs}
Two ordered bases $(v_1, \ldots, v_n)$ and $(w_1, \ldots, w_n)$ of $V$ are \textbf{equivalent} if the change-of-basis matrix $C$ (where $w_j = \sum_i c_{ij} v_i$) has $\det(C) > 0$.

This is an equivalence relation on ordered bases, and it partitions them into exactly \textbf{two} equivalence classes. Each class is called an \textbf{orientation} of $V$.
\end{defn}

\begin{example}
In $\mathbb{R}^2$:
\begin{itemize}
    \item $(e_1, e_2)$ and $(e_1, e_1 + e_2)$ have the same orientation (counterclockwise)
    \item $(e_1, e_2)$ and $(e_2, e_1)$ have opposite orientations ($\det = -1$)
\end{itemize}

% Diagram for orientation in R^2
\begin{center}
\begin{tikzpicture}[scale=1.2]
    % Left: counterclockwise
    \draw[->] (-0.3,0) -- (1.5,0);
    \draw[->] (0,-0.3) -- (0,1.5);
    \draw[->, thick, MagicBlue] (0,0) -- (1,0) node[below] {$e_1$};
    \draw[->, thick, MagicPink] (0,0) -- (0,1) node[left] {$e_2$};
    \draw[->, thick, gray] (0.6,0.15) arc (15:75:0.5);
    \node at (0.5,-0.7) {counterclockwise};
    \node at (0.5,-1.1) {$\det > 0$};
    
    % Right: clockwise
    \begin{scope}[xshift=4cm]
    \draw[->] (-0.3,0) -- (1.5,0);
    \draw[->] (0,-0.3) -- (0,1.5);
    \draw[->, thick, MagicPink] (0,0) -- (1,0) node[below] {$e_2$};
    \draw[->, thick, MagicBlue] (0,0) -- (0,1) node[left] {$e_1$};
    \draw[->, thick, gray] (0.15,0.6) arc (105:15:0.5);
    \node at (0.5,-0.7) {clockwise};
    \node at (0.5,-1.1) {$\det < 0$};
    \end{scope}
\end{tikzpicture}
\end{center}

In $\mathbb{R}^3$: the ``standard'' orientation is the right-hand rule---thumb along $e_1$, fingers curl from $e_2$ toward $e_3$.
\end{example}

\begin{defn}{Oriented Vector Space}{oriented_vs}
An \textbf{oriented vector space} is a pair $(V, \mathcal{O})$ where $\mathcal{O}$ is a choice of orientation on $V$. Bases in the chosen class are called \textbf{positively oriented} or \textbf{preferred}.
\end{defn}

\begin{remark}
For $\dim V = n$, we have $\dim \Lambda^n(V^*) = 1$. An orientation on $V$ is the same as a choice of ``positive'' direction in $\Lambda^n(V^*)$.
\end{remark}

\subsection{Orthogonal and Special Orthogonal Groups}

\begin{defn}{Orthogonal Transformation}{orthogonal_def}
Let $V$ have an inner product $\langle \cdot, \cdot \rangle$. A linear map $U: V \to V$ is \textbf{orthogonal} if it preserves the inner product:
\[\langle U(x), U(y) \rangle = \langle x, y \rangle \quad \forall x, y \in V\]

Equivalently, in an orthonormal basis: $U^T U = I$.
\end{defn}

\begin{thm}{Properties of Orthogonal Maps}{orth_props}
If $U$ is orthogonal:
\begin{enumerate}
    \item $U$ takes orthonormal bases to orthonormal bases
    \item $\det(U^T U) = (\det U)^2 = 1$, so $\det U = \pm 1$
    \item $U$ preserves lengths and angles
\end{enumerate}
\end{thm}

\begin{defn}{O(n) and SO(n)}{on_son}
\begin{itemize}
    \item $O(n) = \{A \in \mathcal{M}_{n \times n} : A^T A = I\}$ is the \textbf{orthogonal group}
    \item $SO(n) = \{A \in O(n) : \det A = +1\}$ is the \textbf{special orthogonal group}
\end{itemize}
These are groups under matrix multiplication.
\end{defn}

\begin{thm}{Orientation and Orthogonal Maps}{orth_orientation}
\begin{itemize}
    \item $A \in O(n)$ with $\det A = +1$: \textbf{orientation-preserving} (rotation)
    \item $A \in O(n)$ with $\det A = -1$: \textbf{orientation-reversing} (reflection)
\end{itemize}
So $SO(n)$ consists exactly of rotations.
\end{thm}

\begin{example}
\begin{itemize}
    \item $SO(2)$: rotations in the plane. Every element looks like $\begin{pmatrix} \cos\theta & -\sin\theta \\ \sin\theta & \cos\theta \end{pmatrix}$.
    \item $O(2) \setminus SO(2)$: reflections, like $\begin{pmatrix} 1 & 0 \\ 0 & -1 \end{pmatrix}$.
    \item $SO(3)$: rotations in 3D (3-dimensional manifold!)
\end{itemize}
\end{example}

\begin{thm}{Dimension of O(n) and SO(n)}{dim_on}
$O(n)$ and $SO(n)$ are both $\frac{n(n-1)}{2}$-dimensional manifolds.
\end{thm}

\begin{proof}
Define $F: \mathcal{M}_{n \times n} \to \text{Sym}(n)$ by $F(A) = A^T A$. Then $O(n) = F^{-1}(I)$.

We check that $I$ is a regular value. The derivative is $dF_A(H) = H^T A + A^T H$. At $A \in O(n)$: $dF_A(H) = H^T A + A^T H$.

For any symmetric $S$, we need to solve $H^T A + A^T H = S$. Taking $H = \frac{1}{2}AS$ works since:
\[H^T A + A^T H = \frac{1}{2}S^T A^T A + \frac{1}{2}A^T A S = \frac{1}{2}S + \frac{1}{2}S = S\]

So $dF_A$ is surjective onto $\text{Sym}(n)$, which has dimension $\frac{n(n+1)}{2}$. Thus:
\[\dim O(n) = n^2 - \frac{n(n+1)}{2} = \frac{n(n-1)}{2}\]

$SO(n) = O(n) \cap \{\det = 1\}$ is an open subset of $O(n)$ (same dimension).
\end{proof}

\section{Volume of Parallelepipeds}

\begin{defn}{Parallelepiped}{parallelepiped}
The $k$-dimensional \textbf{parallelepiped} spanned by $v_1, \ldots, v_k \in V$ is:
\[P(v_1, \ldots, v_k) = \left\{\sum_{i=1}^k t_i v_i : 0 \leq t_i \leq 1\right\}\]
\end{defn}

\begin{thm}{Volume Formula}{vol_formula}
If $V$ has an inner product and $v_1, \ldots, v_k$ are linearly independent:
\[\text{Vol}(P(v_1, \ldots, v_k)) = \sqrt{\det G}\]
where $G = (g_{ij})$ with $g_{ij} = \langle v_i, v_j \rangle$ is the \textbf{Gram matrix}.
\end{thm}

\begin{proof}
Write $v_i$ in an orthonormal basis $(e_1, \ldots, e_n)$: let $U$ be the $n \times k$ matrix with columns $v_i$.

Then $G = U^T U$, and by Gram-Schmidt induction:
\[\text{Vol}(P(v_1, \ldots, v_k)) = \text{Vol}(P(v_1, \ldots, v_{k-1})) \cdot \text{dist}(v_k, \text{span}(v_1, \ldots, v_{k-1}))\]

For orthonormal vectors, volume is 1. The general formula follows by induction.
\end{proof}

\begin{cor}{Special Case: Full Dimension}{vol_full_dim}
For $v_1, \ldots, v_n \in \mathbb{R}^n$:
\[\text{Vol}(P(v_1, \ldots, v_n)) = |\det(v_1 | \cdots | v_n)|\]
\end{cor}

\begin{remark}
If we keep track of orientation (signed volume), we get $\det$ without the absolute value. This is why determinant = signed volume!
\end{remark}

\section{Tensor Products}

The tensor product is one of the most important constructions in algebra. We give both the abstract definition via universal properties and the concrete construction.

\subsection{Construction of the Tensor Product}

\begin{defn}{Tensor Product (Construction)}{tensor_construct}
Given vector spaces $V$ and $W$, we construct $V \otimes W$ as follows:

\textbf{Step 1:} Form the free vector space $\mathbb{R}[V \times W]$ generated by all pairs $(v, w)$.

\textbf{Step 2:} Let $L$ be the subspace spanned by all vectors of the form:
\begin{align*}
&\lambda(v, w) - (\lambda v, w), \quad \lambda(v, w) - (v, \lambda w) \\
&(v_1 + v_2, w) - (v_1, w) - (v_2, w) \\
&(v, w_1 + w_2) - (v, w_1) - (v, w_2)
\end{align*}

\textbf{Step 3:} Define $V \otimes W = \mathbb{R}[V \times W] / L$.

For $v \in V$ and $w \in W$, we write $v \otimes w := \pi_L((v, w)) \in V \otimes W$.
\end{defn}

The relations in $L$ ensure that the map $(v, w) \mapsto v \otimes w$ is bilinear:
\begin{align*}
\lambda(v \otimes w) &= (\lambda v) \otimes w = v \otimes (\lambda w) \\
(v_1 + v_2) \otimes w &= v_1 \otimes w + v_2 \otimes w \\
v \otimes (w_1 + w_2) &= v \otimes w_1 + v \otimes w_2
\end{align*}

\begin{thm}{Universal Property of Tensor Product}{tensor_universal}
For any bilinear map $f: V \times W \to Z$, there exists a unique linear map $\tilde{f}: V \otimes W \to Z$ such that $f(v, w) = \tilde{f}(v \otimes w)$.

\begin{center}
\begin{tikzpicture}
    \node (VW) at (0,1.5) {$V \times W$};
    \node (T) at (0,0) {$V \otimes W$};
    \node (Z) at (2.5,0) {$Z$};
    \draw[->] (VW) -- node[left] {$\otimes$} (T);
    \draw[->] (VW) -- node[above right] {$f$ (bilinear)} (Z);
    \draw[->, dashed] (T) -- node[below] {$\exists!\tilde{f}$ (linear)} (Z);
\end{tikzpicture}
\end{center}
\end{thm}

\begin{proof}
By the universal property of free vector spaces, the function $f: V \times W \to Z$ extends to a linear map $\hat{f}: \mathbb{R}[V \times W] \to Z$.

Since $f$ is bilinear, we have $\hat{f}(L) = 0$. For example:
\[\hat{f}((v_1 + v_2, w) - (v_1, w) - (v_2, w)) = f(v_1 + v_2, w) - f(v_1, w) - f(v_2, w) = 0\]

By the universal property of quotients, $\hat{f}$ descends to $\tilde{f}: V \otimes W \to Z$.

Uniqueness: Any such $\tilde{f}$ must satisfy $\tilde{f}(v \otimes w) = f(v, w)$, and since pure tensors $v \otimes w$ span $V \otimes W$, $\tilde{f}$ is determined.
\end{proof}

\begin{thm}{Dimension of Tensor Product}{dim_tensor}
$\dim(V \otimes W) = \dim(V) \cdot \dim(W)$.

If $\{v_1, \ldots, v_n\}$ is a basis for $V$ and $\{w_1, \ldots, w_m\}$ is a basis for $W$, then $\{v_i \otimes w_j : 1 \leq i \leq n, 1 \leq j \leq m\}$ is a basis for $V \otimes W$.
\end{thm}

\begin{proof}
\textbf{Spanning:} Any $v \otimes w$ can be written as:
\[v \otimes w = \left(\sum_i a_i v_i\right) \otimes \left(\sum_j b_j w_j\right) = \sum_{i,j} a_i b_j (v_i \otimes w_j)\]
Since pure tensors span $V \otimes W$, so do the $v_i \otimes w_j$.

\textbf{Linear independence:} Define $f_{ij}: V \times W \to \mathbb{R}$ by $f_{ij}(v, w) = v_i^*(v) \cdot w_j^*(w)$ where $\{v_i^*\}$ and $\{w_j^*\}$ are dual bases. This is bilinear, so it induces $\tilde{f}_{ij}: V \otimes W \to \mathbb{R}$.

Note that $\tilde{f}_{ij}(v_k \otimes w_\ell) = \delta_{ik} \delta_{j\ell}$.

If $\sum_{i,j} c_{ij}(v_i \otimes w_j) = 0$, applying $\tilde{f}_{ij}$ gives $c_{ij} = 0$.
\end{proof}

\begin{remark}
Not every element of $V \otimes W$ is a pure tensor $v \otimes w$! A general element is a sum $\sum_k v_k \otimes w_k$. The \textbf{rank} of a tensor is the minimum number of pure tensors needed to express it.
\end{remark}

\begin{thm}{Properties of Tensor Product}{tensor_props}
\begin{enumerate}
    \item \textbf{Associativity:} $(U \otimes V) \otimes W \cong U \otimes (V \otimes W)$
    \item \textbf{Commutativity:} $V \otimes W \cong W \otimes V$ via $v \otimes w \mapsto w \otimes v$
    \item \textbf{Unit:} $V \otimes \mathbb{R} \cong V$ via $v \otimes \lambda \mapsto \lambda v$
    \item \textbf{Distribution over direct sum:} $U \otimes (V \oplus W) \cong (U \otimes V) \oplus (U \otimes W)$
\end{enumerate}
\end{thm}

\begin{proof}
Each follows from the universal property. For (1): Both $(U \otimes V) \otimes W$ and $U \otimes (V \otimes W)$ have the same universal property with respect to trilinear maps $U \times V \times W \to Z$.
\end{proof}

\subsection{Tensors and Multilinear Maps}

\begin{thm}{Tensor-Hom Adjunction}{tensor_hom}
There is a canonical isomorphism:
\[V^* \otimes W \cong \text{Hom}(V, W)\]
sending $f \otimes w$ to the linear map $v \mapsto f(v) \cdot w$.
\end{thm}

\begin{proof}
Define $\Phi: V^* \times W \to \text{Hom}(V, W)$ by $\Phi(f, w)(v) = f(v) \cdot w$. This is bilinear, so it induces $\tilde{\Phi}: V^* \otimes W \to \text{Hom}(V, W)$.

\textbf{Surjective:} If $\{v_i\}$ is a basis for $V$ with dual basis $\{v_i^*\}$, any $T \in \text{Hom}(V, W)$ satisfies:
\[T(v) = T\left(\sum_i v_i^*(v) v_i\right) = \sum_i v_i^*(v) T(v_i)\]
so $T = \tilde{\Phi}(\sum_i v_i^* \otimes T(v_i))$.

\textbf{Injective:} Both spaces have dimension $(\dim V)(\dim W)$.
\end{proof}

\begin{cor}{Tensors as Multilinear Forms}{tensor_multilinear}
$(V^*)^{\otimes k} \cong \{\text{$k$-linear forms on } V\}$.

In particular, $V^* \otimes V^* \cong \text{Bil}(V \times V, \mathbb{R})$.
\end{cor}

\begin{defn}{Symmetric and Exterior Powers}{sym_ext_powers}
The \textbf{$k$-th symmetric power} $S^k(V)$ is the quotient of $V^{\otimes k}$ by the relations $v_1 \otimes \cdots \otimes v_k = v_{\sigma(1)} \otimes \cdots \otimes v_{\sigma(k)}$ for all permutations $\sigma$.

The \textbf{$k$-th exterior power} $\Lambda^k(V)$ is the quotient of $V^{\otimes k}$ by the relations $v_1 \otimes \cdots \otimes v_k = \text{sign}(\sigma) \cdot v_{\sigma(1)} \otimes \cdots \otimes v_{\sigma(k)}$.

We write $v_1 \cdots v_k$ for elements of $S^k(V)$ and $v_1 \wedge \cdots \wedge v_k$ for elements of $\Lambda^k(V)$.
\end{defn}

\begin{thm}{Dimensions of Symmetric and Exterior Powers}{dim_sym_ext}
\begin{enumerate}
    \item $\dim S^k(V) = \binom{n + k - 1}{k}$ where $n = \dim V$
    \item $\dim \Lambda^k(V) = \binom{n}{k}$
\end{enumerate}
In particular, $\Lambda^k(V) = 0$ for $k > n$.
\end{thm}

\begin{remark}
We have $(V^*)^{\otimes k} \cong (V^{\otimes k})^*$. The symmetric and skew-symmetric $k$-tensors correspond to $S^k(V^*)$ and $\Lambda^k(V^*)$ respectively.
\end{remark}

\begin{exercises}

\begin{exercise}
Let $M$ be the space of $n \times n$ real matrices. For the bilinear form $B(X, Y) = \text{Tr}(XY)$, show $B$ is symmetric and compute the rank and index of its associated quadratic form.
\end{exercise}

\begin{exercise}
\begin{enumerate}
    \item Prove that any exterior 2-form on an odd-dimensional space is degenerate.
    \item Show that the rank of any skew-symmetric bilinear form is even.
\end{enumerate}
\end{exercise}

\begin{exercise}
Let $\omega$ be an exterior 2-form on a $2n$-dimensional space $V$.
\begin{enumerate}
    \item Prove $\omega$ is non-degenerate iff $\omega^{\wedge n} \neq 0$.
    \item Prove there exists a basis where $\omega$ has the standard symplectic matrix form.
\end{enumerate}
\end{exercise}

\begin{exercise}
Show that $S^2(V)$ (the symmetric square) is generated by tensors of the form $v \otimes v$.
\end{exercise}

\begin{exercise}
Compute the 3-dimensional volume of each face of the parallelepiped $P(v_1, v_2, v_3, v_4) \subset \mathbb{R}^4$ where $v_1 = (1,1,1,1)$, $v_2 = (1,-1,1,1)$, $v_3 = (1,1,-1,1)$, $v_4 = (1,1,1,-1)$.
\end{exercise}

\end{exercises}

\chapter{Differential Forms on Open Sets}

\section{$k$-Forms and the Exterior Derivative}

\begin{defn}{Differential $k$-Form}{diff_kform}
A \textbf{differential $k$-form} on an open set $U \subseteq \mathbb{R}^n$ is:
\[
\omega = \sum_{i_1 < \cdots < i_k} f_{i_1 \ldots i_k}(x) \, dx_{i_1} \wedge \cdots \wedge dx_{i_k}
\]
where the $f_{i_1 \ldots i_k}$ are smooth functions on $U$.
\end{defn}

The space of $k$-forms on $U$ is denoted $\Omega^k(U)$.

\begin{defn}{Exterior Derivative}{ext_deriv}
The \textbf{exterior derivative} $d: \Omega^k(U) \to \Omega^{k+1}(U)$ is defined by:
\[
d\omega = \sum_{i_1 < \cdots < i_k} df_{i_1 \ldots i_k} \wedge dx_{i_1} \wedge \cdots \wedge dx_{i_k}
\]
\end{defn}

\begin{example}
\textbf{(Computing $d$)}
Let $\omega = xy\,dx + z^2\,dy + e^z\,dz$ on $\mathbb{R}^3$. Then:
\begin{align*}
d\omega &= d(xy) \wedge dx + d(z^2) \wedge dy + d(e^z) \wedge dz \\
&= (y\,dx + x\,dy) \wedge dx + 2z\,dz \wedge dy + e^z\,dz \wedge dz \\
&= x\,dy \wedge dx + 2z\,dz \wedge dy \\
&= -x\,dx \wedge dy - 2z\,dy \wedge dz
\end{align*}

Now let $\eta = xyz\,dx \wedge dy$ on $\mathbb{R}^3$. Then:
\[d\eta = d(xyz) \wedge dx \wedge dy = (yz\,dx + xz\,dy + xy\,dz) \wedge dx \wedge dy = xy\,dz \wedge dx \wedge dy = xy\,dx \wedge dy \wedge dz\]
\end{example}

\begin{thm}{Properties of $d$}{d_props}
\begin{enumerate}
    \item $d(d\omega) = 0$ for all $\omega$
    \item $d(\omega \wedge \eta) = d\omega \wedge \eta + (-1)^k \omega \wedge d\eta$ (graded Leibniz rule)
    \item $d$ commutes with pullback: $d(f^*\omega) = f^*(d\omega)$
\end{enumerate}
\end{thm}

\begin{proof}
(1) It suffices to check on generators. For a 0-form (function) $f$:
\[d(df) = d\left(\sum_i \frac{\partial f}{\partial x_i} dx_i\right) = \sum_{i,j} \frac{\partial^2 f}{\partial x_j \partial x_i} dx_j \wedge dx_i\]
This vanishes because $\frac{\partial^2 f}{\partial x_j \partial x_i} = \frac{\partial^2 f}{\partial x_i \partial x_j}$ (symmetry) while $dx_j \wedge dx_i = -dx_i \wedge dx_j$ (antisymmetry).

(2) Write $\omega = f \, dx_I$ (multi-index notation). Then:
\begin{align*}
d(\omega \wedge \eta) &= d(f \, dx_I \wedge \eta) \\
&= df \wedge dx_I \wedge \eta + f \, d(dx_I \wedge \eta) \\
&= d\omega \wedge \eta + (-1)^{|I|} \omega \wedge d\eta
\end{align*}
using that $d(dx_I) = 0$ and $df \wedge dx_I = (-1)^{|I|} dx_I \wedge df$.

(3) For functions: $d(f^*g) = d(g \circ f) = \sum \frac{\partial(g \circ f)}{\partial x_i} dx_i = f^*(dg)$ by chain rule.

For higher forms, use that $f^*$ preserves wedge products and linearity.
\end{proof}

\begin{defn}{Closed and Exact Forms}{closed_exact_k}
A $k$-form $\omega$ is:
\begin{itemize}
    \item \textbf{Closed} if $d\omega = 0$
    \item \textbf{Exact} if $\omega = d\eta$ for some $(k-1)$-form $\eta$
\end{itemize}
Since $d^2 = 0$: Exact $\Rightarrow$ Closed.
\end{defn}

\section{Vector Calculus Operations}

In $\mathbb{R}^3$, the exterior derivative captures all classical operations:

\begin{itemize}
    \item \textbf{Gradient}: $df = \frac{\partial f}{\partial x}dx + \frac{\partial f}{\partial y}dy + \frac{\partial f}{\partial z}dz \leftrightarrow \nabla f$
    \item \textbf{Curl}: $d(P\,dx + Q\,dy + R\,dz) = \text{(curl components)} \leftrightarrow \nabla \times \mathbf{F}$
    \item \textbf{Divergence}: $d(P\,dy \wedge dz + Q\,dz \wedge dx + R\,dx \wedge dy) = (\nabla \cdot \mathbf{F}) dx \wedge dy \wedge dz$
\end{itemize}

The identities $\nabla \times (\nabla f) = 0$ and $\nabla \cdot (\nabla \times \mathbf{F}) = 0$ follow from $d^2 = 0$.

\section{The Hodge Star Operator}

\begin{defn}{Hodge Star}{hodge}
On an oriented Euclidean space $V$ of dimension $n$, the \textbf{Hodge star} $\star: \Lambda^k(V^*) \to \Lambda^{n-k}(V^*)$ satisfies:
\[
\alpha \wedge (\star\beta) = \langle \alpha, \beta \rangle \, \text{vol}
\]
where $\text{vol}$ is the volume form.
\end{defn}

\begin{example}
In $\mathbb{R}^3$: $\star dx = dy \wedge dz$, $\star dy = dz \wedge dx$, $\star dz = dx \wedge dy$, $\star(dx \wedge dy) = dz$, etc.
\end{example}

\begin{example}
\textbf{(Hodge Star in $\mathbb{R}^4$)}
With coordinates $(x, y, z, w)$ and volume form $dx \wedge dy \wedge dz \wedge dw$:
\begin{align*}
\star dx &= dy \wedge dz \wedge dw & \star(dx \wedge dy) &= dz \wedge dw \\
\star dy &= -dx \wedge dz \wedge dw & \star(dx \wedge dz) &= -dy \wedge dw \\
\star dz &= dx \wedge dy \wedge dw & \star(dx \wedge dy \wedge dz) &= dw
\end{align*}
Note: In $\mathbb{R}^4$, $\star\star = +1$ on 2-forms (since $2(4-2) = 4$ is even).
\end{example}

\begin{example}
\textbf{(Computing $\star\star$ in $\mathbb{R}^3$)}
Let $\omega = dx \wedge dy$ (a 2-form). Then:
\[\star\omega = dz, \quad \star\star\omega = \star(dz) = dx \wedge dy\]
But wait---check the sign: $k = 2$, $n = 3$, so $(-1)^{k(n-k)} = (-1)^{2 \cdot 1} = 1$ \checkmark.

For a 1-form $\alpha = dx$: $\star\alpha = dy \wedge dz$, $\star\star\alpha = \star(dy \wedge dz) = dx$.
Here $k = 1$, so $(-1)^{1 \cdot 2} = 1$ \checkmark.
\end{example}

\begin{thm}{Properties of Hodge Star}{hodge_props}
\begin{enumerate}
    \item $\star\star\omega = (-1)^{k(n-k)} \omega$ for $\omega \in \Lambda^k$
    \item $\star 1 = \text{vol}$, $\star\text{vol} = 1$
    \item If $\{e_1, \ldots, e_n\}$ is an oriented orthonormal basis with dual $\{x_1, \ldots, x_n\}$:
    \[\star(x_{i_1} \wedge \cdots \wedge x_{i_k}) = \text{sign}(\sigma) \, x_{j_1} \wedge \cdots \wedge x_{j_{n-k}}\]
    where $\{j_1, \ldots, j_{n-k}\} = \{1, \ldots, n\} \setminus \{i_1, \ldots, i_k\}$ and $\sigma$ is the permutation $(i_1, \ldots, i_k, j_1, \ldots, j_{n-k})$.
\end{enumerate}
\end{thm}

\subsection{Inner Product on Forms}

An inner product on $V$ induces an inner product on $\Lambda^k(V^*)$.

\begin{defn}{Inner Product on $k$-Forms}{inner_prod_forms}
For $\alpha, \beta \in \Lambda^k(V^*)$, define:
\[\langle \alpha, \beta \rangle = \star(\alpha \wedge \star\beta)\]
This is a scalar since $\alpha \wedge \star\beta \in \Lambda^n(V^*)$ and $\star$ takes $n$-forms to $0$-forms.
\end{defn}

\begin{thm}{Properties of Inner Product on Forms}{inner_prod_forms_props}
\begin{enumerate}
    \item The basis $\{x_{i_1} \wedge \cdots \wedge x_{i_k} : 1 \leq i_1 < \cdots < i_k \leq n\}$ is orthonormal.
    \item If $A: V \to V$ is a special orthogonal transformation, then $A^*: \Lambda^k(V^*) \to \Lambda^k(V^*)$ is also orthogonal.
    \item $\langle \alpha, \star\beta \rangle = \langle \star\alpha, \beta \rangle$ (up to sign).
\end{enumerate}
\end{thm}

\begin{thm}{Volume of Parallelepiped}{vol_parallel}
If $v_1, \ldots, v_k \in V$ and $\ell_i = D(v_i) \in V^*$, then:
\[\|\ell_1 \wedge \cdots \wedge \ell_k\| = \text{Vol}(P(v_1, \ldots, v_k))\]
where $P(v_1, \ldots, v_k)$ is the $k$-dimensional parallelepiped spanned by the vectors.
\end{thm}

\begin{proof}
If $v_1, \ldots, v_k$ are linearly dependent, both sides are 0. Otherwise, by Gram-Schmidt, we can reduce to the orthonormal case. For orthonormal $e_1, \ldots, e_k$: $\|x_1 \wedge \cdots \wedge x_k\| = 1$ (where $x_i = D(e_i)$), and $P(e_1, \ldots, e_k)$ has unit volume.
\end{proof}

\section{Interior Product (Contraction)}

\begin{defn}{Interior Product}{interior_prod}
For $v \in V$ and $\omega \in \Lambda^k(V^*)$, the \textbf{interior product} (or \textbf{contraction}) $\iota_v \omega \in \Lambda^{k-1}(V^*)$ is defined by:
\[(\iota_v \omega)(X_1, \ldots, X_{k-1}) = \omega(v, X_1, \ldots, X_{k-1})\]

Alternative notation: $v \lrcorner \omega$ or $i_v \omega$.
\end{defn}

\begin{thm}{Properties of Interior Product}{interior_props}
\begin{enumerate}
    \item $\iota_v$ is linear in $v$: $\iota_{av + bw} = a\iota_v + b\iota_w$
    \item $\iota_v \circ \iota_v = 0$
    \item \textbf{Leibniz rule:} $\iota_v(\alpha \wedge \beta) = (\iota_v \alpha) \wedge \beta + (-1)^{|\alpha|} \alpha \wedge (\iota_v \beta)$
    \item For $f$ a 0-form: $\iota_v(f) = 0$
    \item For $\ell \in V^*$: $\iota_v \ell = \ell(v)$
\end{enumerate}
\end{thm}

\begin{proof}
(2) $(\iota_v \circ \iota_v \omega)(X_1, \ldots, X_{k-2}) = \omega(v, v, X_1, \ldots, X_{k-2}) = 0$ by skew-symmetry.

(3) Follows from the definition by considering the permutations.
\end{proof}

\begin{thm}{Relation to Hodge Star}{interior_hodge}
For $V$ Euclidean with volume form $\text{vol}$ and $v \in V$:
\[\star D(v) = \iota_v(\text{vol})\]
where $D: V \to V^*$ is the isomorphism $D(v)(w) = \langle v, w \rangle$.
\end{thm}

\section{The Codifferential and Laplacian}

\begin{defn}{Codifferential}{codiff}
The \textbf{codifferential} $\delta: \Omega^k(M) \to \Omega^{k-1}(M)$ is defined by:
\[\delta = (-1)^{n(k+1)+1} \star d \star\]
On a compact oriented Riemannian manifold, $\delta$ is the formal adjoint of $d$ with respect to the $L^2$ inner product.
\end{defn}

\begin{thm}{Properties of Codifferential}{codiff_props}
\begin{enumerate}
    \item $\delta \circ \delta = 0$ (follows from $d \circ d = 0$)
    \item $\delta(\omega) = 0$ for 0-forms $\omega$
    \item $\langle d\alpha, \beta \rangle = \langle \alpha, \delta\beta \rangle$ (on compact manifolds)
\end{enumerate}
\end{thm}

\begin{defn}{Hodge Laplacian}{laplacian}
The \textbf{Hodge Laplacian} (or Laplace-de Rham operator) is:
\[\Delta = d\delta + \delta d\]
On functions, this coincides with the usual Laplacian (up to sign conventions).
\end{defn}

\begin{defn}{Harmonic Forms}{harmonic}
A form $\omega$ is \textbf{harmonic} if $\Delta \omega = 0$.

Equivalently, $\omega$ is harmonic iff $d\omega = 0$ and $\delta\omega = 0$ (i.e., closed and coclosed).
\end{defn}

\section{Vector Calculus via Differential Forms}

The classical vector calculus operations in $\mathbb{R}^3$ are special cases of the exterior derivative and Hodge star.

\begin{thm}{Gradient, Curl, Divergence via Forms}{vector_calc_forms}
In $\mathbb{R}^3$ with standard coordinates $(x, y, z)$:

\textbf{Gradient:} For a function $f$:
\[df = \frac{\partial f}{\partial x}dx + \frac{\partial f}{\partial y}dy + \frac{\partial f}{\partial z}dz \longleftrightarrow \nabla f\]

\textbf{Curl:} For a 1-form $\alpha = P\,dx + Q\,dy + R\,dz$:
\[\star d\alpha = \left(\frac{\partial R}{\partial y} - \frac{\partial Q}{\partial z}\right)dx + \left(\frac{\partial P}{\partial z} - \frac{\partial R}{\partial x}\right)dy + \left(\frac{\partial Q}{\partial x} - \frac{\partial P}{\partial y}\right)dz\]
which corresponds to $\nabla \times \mathbf{F}$.

\begin{example}
\textbf{(Curl Computation)}
Let $\mathbf{F} = (y^2 z, xz^2, xy^2)$, so $\alpha = y^2 z\,dx + xz^2\,dy + xy^2\,dz$.

\[d\alpha = (2yz\,dy + y^2\,dz) \wedge dx + (z^2\,dx + 2xz\,dz) \wedge dy + (y^2\,dx + 2xy\,dy) \wedge dz\]
\[= 2yz\,dy \wedge dx + y^2\,dz \wedge dx + z^2\,dx \wedge dy + 2xz\,dz \wedge dy + y^2\,dx \wedge dz + 2xy\,dy \wedge dz\]
\[= (2xy - 2xz)\,dy \wedge dz + (y^2 - y^2)\,dz \wedge dx + (z^2 - 2yz)\,dx \wedge dy\]

Then $\star d\alpha = 2x(y-z)\,dx + 0\,dy + (z^2 - 2yz)\,dz$, so $\nabla \times \mathbf{F} = (2x(y-z), 0, z^2 - 2yz)$.
\end{example}

\textbf{Divergence:} For a 2-form $\beta = P\,dy \wedge dz + Q\,dz \wedge dx + R\,dx \wedge dy$:
\[\star d\beta = \frac{\partial P}{\partial x} + \frac{\partial Q}{\partial y} + \frac{\partial R}{\partial z}\]
which is $\nabla \cdot \mathbf{F}$.
\end{thm}

\begin{remark}
The identities $\nabla \times (\nabla f) = 0$ and $\nabla \cdot (\nabla \times \mathbf{F}) = 0$ are both consequences of $d^2 = 0$:
\begin{itemize}
    \item $\text{curl}(\text{grad } f) = 0$: $\star d(df) = \star d^2 f = 0$
    \item $\text{div}(\text{curl } \mathbf{F}) = 0$: $\star d(\star d\alpha) = 0$ (since $d^2 = 0$ after unwinding)
\end{itemize}
\end{remark}

\begin{thm}{Laplacian on Functions}{laplacian_func}
On $\mathbb{R}^n$, for a function $f$:
\[\Delta f = -\delta df = -\star d \star df = -\sum_{i=1}^n \frac{\partial^2 f}{\partial x_i^2}\]
(with appropriate sign conventions).
\end{thm}

\section{Cartan's Magic Formula}

\begin{thm}{Cartan's Magic Formula}{cartan_magic}
Let $X$ be a vector field and $\mathcal{L}_X$ the Lie derivative. Then:
\[\mathcal{L}_X = d \circ \iota_X + \iota_X \circ d\]
In other words: $\mathcal{L}_X \omega = d(\iota_X \omega) + \iota_X(d\omega)$.
\end{thm}

\begin{proof}[Sketch]
Both sides are derivations of degree 0 on $\Omega^*(M)$, so it suffices to check on functions and exact 1-forms:
\begin{itemize}
    \item On functions $f$: $\mathcal{L}_X f = X(f) = df(X) = \iota_X(df)$ and $d(\iota_X f) = d(0) = 0$.
    \item On $dg$: Both sides equal $d(\mathcal{L}_X g) = d(X(g))$.
\end{itemize}
\end{proof}

\begin{cor}{Lie Derivative of Closed Forms}{lie_closed}
If $d\omega = 0$, then $\mathcal{L}_X \omega = d(\iota_X \omega)$. Thus $\mathcal{L}_X$ preserves closed forms.
\end{cor}

\section{Physical Interpretations}

Differential forms have natural physical interpretations that make vector calculus more transparent.

\subsection{Work and Circulation}

\begin{defn}{Work}{work_form}
If $\mathbf{F}$ is a force field and $\gamma$ is a path, the \textbf{work} done by $\mathbf{F}$ along $\gamma$ is:
\[\text{Work} = \int_\gamma \mathbf{F} \cdot d\mathbf{r} = \int_\gamma \alpha\]
where $\alpha = F_1\,dx + F_2\,dy + F_3\,dz$ is the 1-form associated to $\mathbf{F}$.
\end{defn}

If $\gamma$ is a closed loop, this integral is called the \textbf{circulation} of $\mathbf{F}$.

\subsection{Flux}

\begin{defn}{Flux}{flux_form}
If $\mathbf{F}$ is a vector field and $S$ is a surface, the \textbf{flux} of $\mathbf{F}$ through $S$ is:
\[\text{Flux} = \iint_S \mathbf{F} \cdot \mathbf{n}\, dS = \int_S \beta\]
where $\beta = F_1\,dy \wedge dz + F_2\,dz \wedge dx + F_3\,dx \wedge dy$ is the 2-form associated to $\mathbf{F}$.
\end{defn}

\begin{remark}
The relationship is: $\beta = \iota_{\mathbf{F}}(\text{vol}) = \star D(\mathbf{F})$, where $D$ is the isomorphism $V \to V^*$.
\end{remark}

\subsection{Physical Meaning of $d$}

\begin{itemize}
    \item $df$ represents the rate of change of a scalar field
    \item $d\alpha$ for a 1-form $\alpha$ measures the ``rotation'' (curl)
    \item $d\beta$ for a 2-form $\beta$ measures the ``source density'' (divergence)
\end{itemize}

\subsection{Conservation Laws}

\begin{thm}{Conservation via Stokes}{conservation}
If $\mathbf{J}$ is a current density with $\nabla \cdot \mathbf{J} = 0$ (divergence-free), then the flux through any closed surface is zero---no net flow in or out.

This is the content of $d\beta = 0$ for the associated 2-form.
\end{thm}

\subsection{Electromagnetism Preview}

The Maxwell equations can be written elegantly as:
\[dF = 0, \quad d\star F = J\]
where $F = E \wedge dt + B$ is the electromagnetic 2-form and $J$ is the current 3-form.

\section{Change of Variables}

\begin{thm}{Change of Variables Formula}{cov}
Let $\varphi: A \to B$ be a $C^1$ diffeomorphism between measurable sets. Then for integrable $f$:
\[
\int_B f(y) \, dy = \int_A (f \circ \varphi)(x) |\det(d\varphi_x)| \, dx
\]
\end{thm}

\begin{thm}{Fubini's Theorem}{fubini}
For $f: \mathbb{R}^{n+m} \to \mathbb{R}$ integrable:
\[
\int_{\mathbb{R}^{n+m}} f = \int_{\mathbb{R}^n} \left( \int_{\mathbb{R}^m} f(x, y) \, dy \right) dx
\]
\end{thm}

\begin{exercises}

\begin{exercise}
Let $h: \mathbb{R} \to \mathbb{R}$ be smooth with $\lim_{x \to \pm\infty} h(x) = \pm\infty$ and $h' > 0$. Prove $h$ is a diffeomorphism.
\end{exercise}

\begin{exercise}
Is $x \mapsto x^3$ a diffeomorphism from $\mathbb{R}$ to $\mathbb{R}$?
\end{exercise}

\begin{exercise}
Prove that $\phi(x) = \begin{cases} e^{-1/x^2} & x > 0 \\ 0 & x \leq 0 \end{cases}$ is $C^\infty(\mathbb{R})$.
\end{exercise}

\begin{exercise}
Compute $\int_\gamma f(x^2 + y^2 + z^2)(x\,dx + y\,dy + z\,dz)$ where $\gamma(t) = (5\sin^2(20t), 2-2\cos(t/3), \sin(t)/2)$ for $t \in [0, \pi]$ and $\int_0^1 f(t)\,dt = 1$.
\end{exercise}

\end{exercises}

\chapter{Manifolds}

\section{Submanifolds}

\begin{defn}{Submanifold}{submanifold}
A subset $M \subseteq V$ is a \textbf{$k$-dimensional submanifold} if for every $p \in M$, there exists a neighborhood $U_p$ of $p$ in $V$ and a diffeomorphism $\varphi: U_p \to \mathbb{R}^n$ such that:
\[
\varphi(M \cap U_p) = \{(y_1, \ldots, y_n) : y_{k+1} = \cdots = y_n = 0\}
\]
\end{defn}

Equivalently:
\begin{itemize}
    \item \textbf{As a graph}: Locally $M = \{(x, f(x)) : x \in U \subseteq \mathbb{R}^k\}$ for some smooth $f$.
    \item \textbf{As a level set}: Locally $M = \{F_1 = \cdots = F_{n-k} = 0\}$ where $dF_1, \ldots, dF_{n-k}$ are linearly independent.
    \item \textbf{As an image}: $M = \text{Im}(\varphi)$ for an embedding $\varphi: U \to V$.
\end{itemize}

\begin{example}
$S^{n-1} = \{x \in \mathbb{R}^n : \|x\|^2 = 1\}$ is a manifold (level set of $F(x) = \|x\|^2 - 1$).
\end{example}

\section{Tangent Space}

\begin{defn}{Tangent Space}{tangent_space}
For $M$ a manifold and $p \in M$:
\begin{itemize}
    \item If $M = \text{Im}(\varphi)$, then $T_p M = \text{Im}(d_{\varphi^{-1}(p)} \varphi)$.
    \item If $M = \{F = 0\}$, then $T_p M = \ker(d_p F)$.
\end{itemize}
\end{defn}

\begin{center}
\begin{tikzpicture}[scale=1.2]
    % Draw the manifold (a curved surface)
    \draw[thick, blue] plot[smooth, tension=0.8] coordinates {(-2,-0.5) (-1,0.5) (0,0.3) (1,0.8) (2,0.2)};
    \draw[thick, blue] plot[smooth, tension=0.8] coordinates {(-2,-0.8) (-1,0.2) (0,0) (1,0.5) (2,-0.1)};
    
    % Mark point p
    \fill (0,0.15) circle (2pt) node[below right] {$p$};
    
    % Draw tangent plane (as a parallelogram)
    \draw[thick, red, fill=red!10, opacity=0.5] (-1,-0.5) -- (1,-0.5) -- (1.3,0.8) -- (-0.7,0.8) -- cycle;
    
    % Tangent vectors
    \draw[->, thick, MagicTeal] (0,0.15) -- (0.8,0.15) node[right] {$v$};
    \draw[->, thick, MagicTeal] (0,0.15) -- (0.15,0.65) node[above] {$w$};
    
    % Labels
    \node at (2.5,0.3) {$M$};
    \node[red] at (1.8,0.9) {$T_pM$};
\end{tikzpicture}
\end{center}

The tangent space is independent of the parametrization/presentation.

\begin{defn}{Tangent Bundle}{tangent_bundle}
The \textbf{tangent bundle} is $TM = \{(p, v) : p \in M, v \in T_p M\} \subseteq V \times V$.
\end{defn}

$TM$ is a $2k$-dimensional manifold.

\section{Measure Zero and Measurable Sets}

\begin{defn}{Measure Zero}{measure_zero}
A set $S \subseteq \mathbb{R}^n$ has \textbf{measure zero} if for every $\varepsilon > 0$, there exist countably many rectangles $R_1, R_2, \ldots$ such that:
\begin{enumerate}
    \item $S \subseteq \bigcup_{i=1}^\infty R_i$
    \item $\sum_{i=1}^\infty \text{Vol}(R_i) < \varepsilon$
\end{enumerate}
\end{defn}

\begin{thm}{Properties of Measure Zero}{meas_zero_props}
\begin{enumerate}
    \item Countable sets have measure zero.
    \item Countable unions of measure zero sets have measure zero.
    \item Subsets of measure zero sets have measure zero.
    \item If $S \subset \mathbb{R}^n$ has measure zero and $f: \mathbb{R}^n \to \mathbb{R}^n$ is Lipschitz, then $f(S)$ has measure zero.
    \item If $\dim M < n$, any smooth map $g: M \to \mathbb{R}^n$ has image of measure zero.
\end{enumerate}
\end{thm}

\begin{proof}[Sketch of (5)]
Cover $M$ with countably many coordinate charts. Each chart has dimension $< n$, so its image lies in a ``thin'' set. The union is still measure zero.
\end{proof}

\section{Regular and Critical Values}

\begin{defn}{Regular and Critical Values}{reg_crit}
Let $f: M \to N$ be a smooth map between manifolds.
\begin{itemize}
    \item A point $p \in M$ is a \textbf{regular point} if $Df_p$ is surjective.
    \item A point $p$ is a \textbf{critical point} if $Df_p$ is not surjective (i.e., $Df_p$ is degenerate).
    \item A value $y \in N$ is a \textbf{regular value} if every $p \in f^{-1}(y)$ is regular (or if $f^{-1}(y) = \emptyset$).
    \item Otherwise $y$ is a \textbf{critical value}.
\end{itemize}
\end{defn}

\begin{defn}{Critical Set}{critical_set}
The \textbf{critical set} (or \textbf{singular locus}) of $f$ is:
\[\Sigma(f) = \{x \in M : d_x f \text{ is degenerate (not surjective)}\}\]
The set of critical values is $f(\Sigma(f))$.
\end{defn}

\begin{thm}{Sard's Theorem}{sard}
The set of critical values of a smooth map $f: M \to N$ has measure zero in $N$.
\end{thm}

\begin{proof}
We prove the case $f: U \to \mathbb{R}^n$ where $U \subseteq \mathbb{R}^m$ is open.

\textbf{Step 1: Reduce to a rectangle.}
Cover $U$ by countably many closed rectangles $R_1, R_2, \ldots$. Since countable unions of measure zero sets have measure zero, it suffices to show $f(\Sigma(f) \cap R)$ has measure zero for each rectangle $R$.

\textbf{Step 2: Subdivide the rectangle.}
Fix a rectangle $R$ of side length $\ell$. Subdivide $R$ into $N^m$ small cubes of side $\ell/N$.

\textbf{Step 3: Estimate on each small cube.}
Let $Q$ be a small cube of side $s = \ell/N$ containing a critical point $x_0$.

Since $f$ is $C^1$ on $R$, there exists $M > 0$ such that $\|Df_x\| \leq M$ for all $x \in R$.

For $x \in Q$: $\|f(x) - f(x_0)\| \leq M \cdot \|x - x_0\| \leq M \cdot \sqrt{m} \cdot s$.

So $f(Q)$ is contained in a ball of radius $Ms\sqrt{m}$ around $f(x_0)$.

\textbf{Step 4: Use degeneracy.}
Here's the key: at a critical point $x_0$, the image $Df_{x_0}(\mathbb{R}^m)$ has dimension $< n$.

By Taylor: $f(x) = f(x_0) + Df_{x_0}(x - x_0) + O(\|x - x_0\|^2)$.

The linear part $Df_{x_0}(x - x_0)$ lies in a subspace of dimension $< n$.

The error is $O(s^2)$. So $f(Q)$ lies within distance $O(s^2)$ of a $(n-1)$-dimensional subspace.

This means $f(Q)$ lies in a ``slab'' of thickness $O(s^2)$ and width $O(s)$ in the other directions.

$\text{Vol}(f(Q)) \leq C \cdot s^{n-1} \cdot s^2 = C \cdot s^{n+1}$

\textbf{Step 5: Sum over cubes.}
There are at most $N^m$ cubes, each with $s = \ell/N$. The total volume of critical values is at most:
\[\text{Vol}(f(\Sigma(f) \cap R)) \leq N^m \cdot C \cdot s^{n+1} = N^m \cdot C \cdot \left(\frac{\ell}{N}\right)^{n+1} = C\ell^{n+1} \cdot N^{m-n-1}\]

If $m \leq n$, then $m - n - 1 < 0$, so this goes to 0 as $N \to \infty$.

(For $m > n$, we need higher derivatives and a more refined argument, but the idea is similar.)
\end{proof}

\begin{cor}{Generic Regular Values}{generic_reg}
Almost every $y \in N$ is a regular value. In particular, regular values are dense.
\end{cor}

\begin{thm}{Preimage Theorem}{preimage}
If $y$ is a regular value of $f: M^m \to N^n$, then $f^{-1}(y)$ is a smooth submanifold of dimension $m - n$ (empty if $m < n$).
\end{thm}

\begin{proof}
At any $p \in f^{-1}(y)$, $Df_p: T_p M \to T_y N$ is surjective. By the Implicit Function Theorem, near $p$ we can write $f^{-1}(y)$ as a graph over $\ker(Df_p)$, which has dimension $m - n$.
\end{proof}

\begin{remark}
This explains why ``most'' level sets $f^{-1}(c)$ are smooth manifolds---the ``bad'' values $c$ are measure zero, so generically we're fine.
\end{remark}

\section{Manifolds with Boundary}

\begin{defn}{Manifold with Boundary}{mfld_bdy}
A $k$-dimensional \textbf{manifold with boundary} $M$ is locally diffeomorphic to either $\mathbb{R}^k$ or $\mathbb{R}^k_+ = \{x_k \geq 0\}$.

The \textbf{boundary} $\partial M$ consists of points mapping to $\{x_k = 0\}$; it is a $(k-1)$-dimensional manifold.
\end{defn}

\begin{example}
$D^n = \{x \in \mathbb{R}^n : \|x\| \leq 1\}$ is a manifold with boundary $\partial D^n = S^{n-1}$.
\end{example}

\section{Bump Functions and Partitions of Unity}

This is one of the most important technical tools in differential geometry!

\subsection{Bump Functions}

\begin{lem}{Existence of Smooth Bump Function}{bump_exist}
There exists a $C^\infty$ function $\rho: \mathbb{R} \to [0, \infty)$ such that:
\begin{itemize}
    \item $\rho(x) = 0$ for $|x| \geq 1$
    \item $\rho(x) > 0$ for $|x| < 1$
    \item $\rho(-x) = \rho(x)$
\end{itemize}
\end{lem}

\begin{proof}
Consider:
\[h(x) = \begin{cases} e^{-1/x^2} & x > 0 \\ 0 & x \leq 0 \end{cases}\]

All derivatives of $e^{-1/x^2}$ go to 0 as $x \to 0^+$ (since exponential decay beats any polynomial), so $h$ is $C^\infty$.

Then $\rho(x) = h(1+x) \cdot h(1-x)$ works.
\end{proof}

\begin{defn}{Bump Function in $\mathbb{R}^n$}{bump_rn}
For $a \in \mathbb{R}^n$ and $\delta > 0$, define:
\[\psi_{a,\delta}(x) = \rho\left(\frac{\|x - a\|^2}{\delta^2}\right)\]
This is $C^\infty$, supported in $\overline{B_\delta(a)}$, and positive on $B_\delta(a)$.
\end{defn}

\subsection{Cut-off Functions}

\begin{thm}{Existence of Cut-off Functions}{cutoff}
Let $C \subset \mathbb{R}^n$ be compact and $U \supset C$ open. Then there exists $\sigma \in C^\infty(\mathbb{R}^n)$ with:
\begin{itemize}
    \item $0 \leq \sigma \leq 1$
    \item $\sigma \equiv 1$ on $C$
    \item $\text{supp}(\sigma) \subset U$
\end{itemize}
\end{thm}

\begin{proof}
Choose $\varepsilon > 0$ so that the $\varepsilon$-neighborhood $U_\varepsilon(C) \subset U$.

By compactness, cover $C$ by finitely many balls $B_\varepsilon(z_1), \ldots, B_\varepsilon(z_N)$.

Let $\sigma_1 = \sum_{j=1}^N \psi_{z_j, \varepsilon/2}$. This is positive on $U_{\varepsilon/2}(C)$ and supported in $U$.

Similarly construct $\sigma_2$ supported on the complement of $C$ (in a bounded region).

Then $\sigma = \frac{\sigma_1}{\sigma_1 + \sigma_2}$ does the job.
\end{proof}

\subsection{Partitions of Unity}

\begin{defn}{Partition of Unity}{partition_unity}
Let $\{U_j\}_{j=1}^N$ be an open cover of a compact set $C$. A \textbf{partition of unity subordinate to this cover} is a collection of $C^\infty$ functions $\theta_1, \ldots, \theta_K$ such that:
\begin{enumerate}
    \item $\sum_{j=1}^K \theta_j(x) = 1$ for all $x \in C$
    \item Each $\text{supp}(\theta_j)$ is contained in some $U_i$
\end{enumerate}
\end{defn}

\begin{thm}{Existence of Partitions of Unity}{partition_exists}
For any compact $C$ and open cover $\{U_j\}$, a partition of unity subordinate to this cover exists.
\end{thm}

\begin{proof}
By compactness, find finitely many balls $B_\varepsilon(z_1), \ldots, B_\varepsilon(z_K)$ covering $C$, each contained in some $U_j$.

Let $\psi_j = \psi_{z_j, \varepsilon}$ be bump functions. Then $\sum_j \psi_j > 0$ on a neighborhood of $C$.

Define:
\[\theta_j = \frac{\psi_j \cdot \sigma}{\sum_k \psi_k}\]
where $\sigma$ is a cut-off for $C$.

Then $\sum_j \theta_j = \sigma = 1$ on $C$, and each $\theta_j$ is supported in some $U_i$.
\end{proof}

\subsection{Local to Global}

\begin{remark}
Why partitions of unity matter: they let us \textbf{globalize local constructions}.

For instance, to integrate a $k$-form $\omega$ over a manifold $M$:
\begin{enumerate}
    \item Cover $M$ by coordinate charts $\{U_j\}$
    \item Take a partition of unity $\{\theta_j\}$ subordinate to this cover
    \item Write $\omega = \sum_j \theta_j \omega$ (each $\theta_j \omega$ is supported in a chart)
    \item Define $\int_M \omega = \sum_j \int_{U_j} \theta_j \omega$
\end{enumerate}
This is well-defined (independent of choices)!
\end{remark}

\section{Stereographic Projection}

\begin{defn}{Stereographic Projection}{stereo}
Let $N = (0, \ldots, 0, 1)$ and $S = (0, \ldots, 0, -1)$ be the north and south poles of $S^n \subset \mathbb{R}^{n+1}$.

The \textbf{stereographic projection from the north pole} is:
\[p_+: S^n \setminus \{N\} \to \mathbb{R}^n, \quad p_+(x_1, \ldots, x_{n+1}) = \frac{1}{1 - x_{n+1}}(x_1, \ldots, x_n)\]

Similarly for $p_-: S^n \setminus \{S\} \to \mathbb{R}^n$ from the south pole.
\end{defn}

% Stereographic projection diagram
\begin{center}
\begin{tikzpicture}[scale=1.5]
    % Draw sphere (circle)
    \draw[thick] (0,0) circle (1);
    
    % North and south poles
    \fill (0,1) circle (1.5pt) node[above] {$N$};
    \fill (0,-1) circle (1.5pt) node[below] {$S$};
    
    % Point on sphere
    \fill (0.7,0.714) circle (1.5pt) node[above right] {$p$};
    
    % Projection line
    \draw[dashed] (0,1) -- (1.4,0);
    
    % Projection point
    \fill (1.4,0) circle (1.5pt) node[below] {$p_+(p)$};
    
    % Equatorial plane
    \draw[thick, MagicBlue] (-1.8,0) -- (1.8,0) node[right] {$\mathbb{R}^n$};
\end{tikzpicture}
\end{center}

\begin{thm}{Properties of Stereographic Projection}{stereo_props}
\begin{enumerate}
    \item $p_+$ and $p_-$ are diffeomorphisms onto $\mathbb{R}^n$
    \item On $S^n \setminus \{N, S\}$, we have $p_- \circ p_+^{-1}(x) = \frac{x}{\|x\|^2}$ (inversion)
    \item $\{(S^n \setminus \{N\}, p_+), (S^n \setminus \{S\}, p_-)\}$ forms an atlas for $S^n$
\end{enumerate}
\end{thm}

\begin{proof}
The inverse of $p_+$ is:
\[p_+^{-1}(y) = \left(\frac{2y_1}{1 + \|y\|^2}, \ldots, \frac{2y_n}{1 + \|y\|^2}, \frac{\|y\|^2 - 1}{\|y\|^2 + 1}\right)\]

For the transition map, direct computation gives $p_- \circ p_+^{-1}(x) = x/\|x\|^2$, which is smooth on $\mathbb{R}^n \setminus \{0\}$.
\end{proof}

\section{Vector Bundles}

\begin{defn}{Vector Bundle}{vec_bundle}
A \textbf{vector bundle of rank $r$} over a manifold $M$ is a collection of $r$-dimensional vector spaces $\{L_p\}_{p \in M}$ that ``varies smoothly'' with $p$.

More precisely: for each $p \in M$, there's a neighborhood $U$ and smooth linearly independent vector fields $v_1(x), \ldots, v_r(x) \in L_x$ for $x \in U$.
\end{defn}

\begin{example}
\begin{itemize}
    \item \textbf{Tangent bundle} $TM = \{T_p M\}_{p \in M}$ (rank $n$ for $n$-manifold)
    \item \textbf{Normal bundle} $NM = \{T_p M^\perp\}_{p \in M}$ for $M \subset \mathbb{R}^n$ (rank = codimension)
    \item \textbf{Cotangent bundle} $T^*M = \{T_p^* M\}_{p \in M}$
\end{itemize}
\end{example}

\begin{defn}{Trivial Bundle}{trivial_bundle}
A vector bundle $L$ over $M$ is \textbf{trivial} if there exist globally defined, linearly independent sections $v_1, \ldots, v_r: M \to L$.

Equivalently, $L \cong M \times \mathbb{R}^r$.
\end{defn}

\begin{example}
\begin{itemize}
    \item $TU$ for $U \subset \mathbb{R}^n$ is trivial (take coordinate vector fields)
    \item $TS^1$ is trivial (take the tangent vector field going around)
    \item $TS^2$ is \textbf{not trivial}! (``hairy ball theorem'': no nonvanishing vector field on $S^2$)
\end{itemize}
\end{example}

\begin{example}
\textbf{Trivial vs nontrivial:}
\begin{itemize}
    \item The \textbf{cylinder} $S^1 \times \mathbb{R}$ is a trivial rank-1 bundle over $S^1$
    \item The \textbf{Möbius strip} is a nontrivial rank-1 bundle over $S^1$ (no global nonvanishing section---the ``twist'' prevents it)
\end{itemize}
\end{example}

\begin{defn}{Codimension}{codim}
The \textbf{codimension} of a $k$-submanifold $M \subset \mathbb{R}^n$ is $\text{codim}(M) = n - k$.

This is the rank of the normal bundle.
\end{defn}

\section{Abstract Manifolds and Charts}

So far we've thought of manifolds as subsets of $\mathbb{R}^n$. But we can define them more abstractly!

\begin{defn}{Abstract Manifold}{abstract_mfld}
An $n$-dimensional \textbf{smooth manifold} is a set $M$ together with:
\begin{enumerate}
    \item A collection of subsets $\{U_\alpha\}_{\alpha \in A}$ covering $M$ (coordinate neighborhoods)
    \item For each $\alpha$, a bijection $\Phi_\alpha: U_\alpha \to G_\alpha \subset \mathbb{R}^n$ (coordinate map)
\end{enumerate}
such that:
\begin{itemize}
    \item Each $G_\alpha = \Phi_\alpha(U_\alpha)$ is open in $\mathbb{R}^n$
    \item The \textbf{transition maps} $\Phi_\beta \circ \Phi_\alpha^{-1}: \Phi_\alpha(U_\alpha \cap U_\beta) \to \Phi_\beta(U_\alpha \cap U_\beta)$ are smooth diffeomorphisms
\end{itemize}
The collection $\{(U_\alpha, \Phi_\alpha)\}$ is called an \textbf{atlas}.
\end{defn}

\begin{remark}
Two atlases define the same smooth structure if their union is still an atlas (i.e., transition maps between charts from different atlases are smooth).
\end{remark}

\begin{defn}{Coordinate Chart}{coord_chart}
A pair $(U, \Phi)$ where $U \subset M$ and $\Phi: U \to G \subset \mathbb{R}^n$ is a diffeomorphism is called a \textbf{coordinate chart} (or local coordinates, or patch).

The inverse $\Psi = \Phi^{-1}: G \to U$ is called a \textbf{parametrization}.
\end{defn}

\begin{example}
\textbf{Different ways to think about $S^1$:}
\begin{enumerate}
    \item As $\{x^2 + y^2 = 1\} \subset \mathbb{R}^2$ (level set)
    \item As $\{\theta \mapsto (\cos\theta, \sin\theta) : \theta \in [0, 2\pi)\}$ (parametrization)
    \item As $\mathbb{R}/\mathbb{Z}$ (quotient---identify $\theta \sim \theta + 2\pi$)
    \item Via two charts: $U_1 = S^1 \setminus \{(1,0)\}$, $U_2 = S^1 \setminus \{(-1,0)\}$ with angle coordinates
\end{enumerate}
\end{example}

\subsection{Gluing Construction}

\begin{remark}
You can build manifolds by ``gluing'' pieces of $\mathbb{R}^n$ together via transition maps:

Take copies $G_1, G_2, \ldots$ of open sets in $\mathbb{R}^n$. Specify gluing maps $h_{\alpha\beta}: G_\alpha \supset U_{\alpha\beta} \to U_{\beta\alpha} \subset G_\beta$ satisfying:
\begin{itemize}
    \item $h_{\beta\alpha} = h_{\alpha\beta}^{-1}$
    \item $h_{\gamma\beta} \circ h_{\beta\alpha} = h_{\gamma\alpha}$ (cocycle condition)
\end{itemize}
Then $M = \bigsqcup G_\alpha / \sim$ where $x \sim h_{\alpha\beta}(x)$.
\end{remark}

\begin{example}
$S^n$ via gluing: Take two copies of $\mathbb{R}^n$ and glue $\mathbb{R}^n \setminus \{0\}$ to itself via $x \mapsto x/\|x\|^2$ (inversion). This is the stereographic atlas!
\end{example}

\section{Whitney Embedding Theorem}

A natural question: can every abstract manifold be realized as a submanifold of some $\mathbb{R}^N$?

\begin{thm}{Whitney Embedding Theorem}{whitney}
Every smooth $n$-dimensional manifold can be smoothly embedded in $\mathbb{R}^{2n}$.

(And immersed in $\mathbb{R}^{2n-1}$ for $n > 1$.)
\end{thm}

\begin{remark}
So our ``submanifold of $\mathbb{R}^N$'' viewpoint isn't really a restriction---every manifold fits this picture!
\end{remark}

\begin{example}
\begin{itemize}
    \item $S^1$ embeds in $\mathbb{R}^2$ (Whitney bound: $2 \cdot 1 = 2$, achieved!)
    \item $\mathbb{RP}^2$ doesn't embed in $\mathbb{R}^3$, but does in $\mathbb{R}^4$ (Whitney bound: $2 \cdot 2 = 4$)
    \item The Klein bottle doesn't embed in $\mathbb{R}^3$ but does in $\mathbb{R}^4$
\end{itemize}
\end{example}

\section{Graphical Submanifolds}

\begin{defn}{Graphical Submanifold}{graphical}
A submanifold $M \subset V$ is \textbf{graphical with respect to a splitting $V = W_1 \oplus W_2$} if $M$ can locally be written as:
\[M = \{(w_1, f(w_1)) : w_1 \in U \subset W_1\}\]
for some smooth $f: U \to W_2$.
\end{defn}

\begin{thm}{Local Graph Theorem}{local_graph}
Every $k$-submanifold $M \subset \mathbb{R}^n$ is locally graphical. More precisely:

At any point $p \in M$, after possibly permuting coordinates, $M$ near $p$ is the graph of a function $f: U \subset \mathbb{R}^k \to \mathbb{R}^{n-k}$.
\end{thm}

\begin{proof}
Use the implicit function theorem! If $M = \{F_1 = \cdots = F_{n-k} = 0\}$ near $p$ with $dF_1, \ldots, dF_{n-k}$ independent, then some $(n-k) \times (n-k)$ minor of the Jacobian is nonzero. WLOG it's the last $n-k$ variables. Then IFT gives $(x_{k+1}, \ldots, x_n) = f(x_1, \ldots, x_k)$.
\end{proof}

\begin{defn}{Splitting Diffeomorphism}{split_diffeo}
A \textbf{splitting diffeomorphism} at $p \in M$ is a diffeomorphism $\Phi: U_p \to \mathbb{R}^n$ such that:
\[\Phi(M \cap U_p) = \{(y_1, \ldots, y_k, 0, \ldots, 0)\}\]
This ``straightens out'' $M$ to a coordinate plane.
\end{defn}

\begin{remark}
Being graphical with respect to a splitting is equivalent to saying the projection $\pi_{W_1}: M \to W_1$ is a local diffeomorphism.
\end{remark}

\section{Tubular Neighborhoods}

\begin{defn}{Tubular Neighborhood}{tubular}
A \textbf{tubular neighborhood} of a submanifold $M \subset V$ is a neighborhood $U \supset M$ together with a smooth retraction $r: U \to M$ such that each fiber $r^{-1}(p)$ is a ball in the normal space $N_p M$.
\end{defn}

\begin{thm}{Existence of Tubular Neighborhoods}{tubular_exist}
Every compact submanifold $M \subset \mathbb{R}^n$ has a tubular neighborhood.
\end{thm}

\begin{remark}
Intuitively: you can ``thicken'' $M$ to a neighborhood where every point has a unique closest point on $M$.
\end{remark}

\section{Orientation and Integration}

\begin{defn}{Orientation}{orientation_mfld}
An \textbf{orientation} on a manifold $M$ is a continuous choice of orientation on each tangent space $T_p M$.
\end{defn}

\begin{defn}{Integration of $k$-Forms}{int_kform}
For $M$ an oriented $k$-dimensional manifold and $\omega$ a $k$-form with compact support:
\[
\int_M \omega
\]
is defined using partitions of unity and local parametrizations.
\end{defn}

\begin{exercises}

\begin{exercise}
Let $SO(3) \subset \mathcal{M}_{3,3}$ be the group of orthogonal matrices with determinant $+1$. Let $UT(S^2) = \{(x,y) \in \mathbb{R}^3 \times \mathbb{R}^3 : \|x\| = \|y\| = 1, x \cdot y = 0\}$.
\begin{enumerate}
    \item Show that $SO(3)$ and $UT(S^2)$ are submanifolds of $\mathbb{R}^9$ and $\mathbb{R}^6$ respectively.
    \item Prove they are diffeomorphic.
\end{enumerate}
\end{exercise}

\begin{exercise}
Consider the torus $T = \{z^2 + (r-2)^2 = 1\}$ in cylindrical coordinates and $T' = S^1 \times S^1 \subset \mathbb{R}^4$.
\begin{enumerate}
    \item Prove $T$ is a submanifold of $\mathbb{R}^3$.
    \item Prove $T$ and $T'$ are diffeomorphic.
    \item Compute the surface areas of $T$ and $T'$.
\end{enumerate}
\end{exercise}

\begin{exercise}
Let $f: \mathbb{R}^n \to \mathbb{R}^{n(n+1)/2}$ map $x$ to all products $x_i x_j$. Let $M = f(S^{n-1})$ (the real projective space).
\begin{enumerate}
    \item Describe $f^{-1}(y) \cap S^{n-1}$ for $y \in M$.
    \item Prove $M$ is an $(n-1)$-dimensional manifold.
\end{enumerate}
\end{exercise}

\end{exercises}

\chapter{Stokes' Theorem and Cohomology}

\section{Stokes' Theorem}

\begin{thm}{Stokes' Theorem}{stokes}
Let $M$ be an oriented $k$-dimensional manifold with boundary, and let $\omega$ be a $(k-1)$-form with compact support. Then:
\[
\int_M d\omega = \int_{\partial M} \omega
\]
where $\partial M$ has the induced orientation.
\end{thm}

% Stokes' theorem visualization
\begin{center}
\begin{tikzpicture}[scale=0.9]
    % The manifold M (a region)
    \draw[thick, fill=MagicBlue!20] plot[smooth cycle, tension=0.8] coordinates {(-2,0) (-1.5,1.5) (0.5,2) (2,1) (2.5,-0.5) (1,-1.5) (-1,-1)};
    
    % Boundary orientation (arrows)
    \draw[->, thick, MagicPink] (-1.75,0.75) arc (150:100:0.5);
    \draw[->, thick, MagicPink] (-0.5,1.85) arc (140:90:0.5);
    \draw[->, thick, MagicPink] (1.5,1.5) arc (100:50:0.5);
    \draw[->, thick, MagicPink] (2.3,0.2) arc (60:10:0.5);
    \draw[->, thick, MagicPink] (1.8,-1.1) arc (330:280:0.5);
    \draw[->, thick, MagicPink] (0,-1.4) arc (250:200:0.5);
    \draw[->, thick, MagicPink] (-1.5,-0.5) arc (210:160:0.5);
    
    % Labels
    \node at (0,0) {$M$};
    \node[MagicPink] at (3,1) {$\partial M$};
    
    % The equation
    \node at (0,-2.3) {$\displaystyle\int_M d\omega = \int_{\partial M} \omega$};
\end{tikzpicture}
\end{center}

\begin{proof}
We prove Stokes' theorem in three steps.

\textbf{Step 1: The local case in $\mathbb{R}^k$ and $\mathbb{R}^k_+$.}

\textit{Case A: Interior (no boundary).} Let $\omega$ be a compactly supported $(k-1)$-form on $\mathbb{R}^k$. We can write:
\[\omega = \sum_{i=1}^k f_i(x) \, dx_1 \wedge \cdots \wedge \widehat{dx_i} \wedge \cdots \wedge dx_k\]
where $\widehat{dx_i}$ means omission. Then:
\[d\omega = \sum_{i=1}^k (-1)^{i-1} \frac{\partial f_i}{\partial x_i} \, dx_1 \wedge \cdots \wedge dx_k\]

Consider one term: $\int_{\mathbb{R}^k} \frac{\partial f_i}{\partial x_i} \, dx_1 \cdots dx_k$.

By Fubini, integrate first in $x_i$:
\[\int_{-\infty}^\infty \frac{\partial f_i}{\partial x_i} \, dx_i = f_i(x_1, \ldots, \infty, \ldots, x_k) - f_i(x_1, \ldots, -\infty, \ldots, x_k) = 0\]
since $f_i$ has compact support. So $\int_{\mathbb{R}^k} d\omega = 0 = \int_{\emptyset} \omega$.

\textit{Case B: Half-space $\mathbb{R}^k_+ = \{x_k \geq 0\}$.} For $i < k$, the same argument gives 0 (integrating in $x_i$ from $-\infty$ to $\infty$).

For $i = k$, the term $f_k \, dx_1 \wedge \cdots \wedge dx_{k-1}$ gives:
\begin{align*}
\int_{\mathbb{R}^k_+} d(f_k \, dx_1 \wedge \cdots \wedge dx_{k-1}) &= (-1)^{k-1} \int_{\mathbb{R}^{k-1}} \left[ \int_0^\infty \frac{\partial f_k}{\partial x_k} \, dx_k \right] dx_1 \cdots dx_{k-1} \\
&= (-1)^{k-1} \int_{\mathbb{R}^{k-1}} [f_k]_{x_k=0}^{x_k=\infty} \, dx_1 \cdots dx_{k-1} \\
&= (-1)^k \int_{\mathbb{R}^{k-1}} f_k(x_1, \ldots, x_{k-1}, 0) \, dx_1 \cdots dx_{k-1}
\end{align*}

The boundary $\partial \mathbb{R}^k_+ = \{x_k = 0\} \cong \mathbb{R}^{k-1}$ with the induced orientation gives:
\[\int_{\partial \mathbb{R}^k_+} \omega = (-1)^k \int_{\mathbb{R}^{k-1}} f_k|_{x_k=0} \, dx_1 \cdots dx_{k-1}\]
(The sign $(-1)^k$ comes from the induced boundary orientation: the outward normal is $-e_k$.)

\textbf{Step 2: Partition of unity.}

Cover $M$ by coordinate charts $\{U_\alpha\}$. Take a partition of unity $\{\theta_\alpha\}$ subordinate to this cover. Write:
\[\omega = \sum_\alpha \theta_\alpha \omega\]
Each $\theta_\alpha \omega$ is supported in a chart. Since $d$ and integration are linear:
\[\int_M d\omega = \sum_\alpha \int_M d(\theta_\alpha \omega), \quad \int_{\partial M} \omega = \sum_\alpha \int_{\partial M} \theta_\alpha \omega\]

So it suffices to prove $\int_M d(\theta_\alpha \omega) = \int_{\partial M} \theta_\alpha \omega$ for each $\alpha$.

\textbf{Step 3: Two types of charts.}

\textit{Interior charts:} If $U_\alpha$ doesn't meet $\partial M$, then $U_\alpha$ maps to $\mathbb{R}^k$ and $\theta_\alpha \omega$ has compact support. By Case A:
\[\int_{U_\alpha} d(\theta_\alpha \omega) = 0 = \int_{\emptyset} \theta_\alpha \omega\]

\textit{Boundary charts:} If $U_\alpha$ meets $\partial M$, then $U_\alpha$ maps to $\mathbb{R}^k_+$ with $\partial M \cap U_\alpha$ mapping to $\{x_k = 0\}$. By Case B:
\[\int_{U_\alpha} d(\theta_\alpha \omega) = \int_{\partial M \cap U_\alpha} \theta_\alpha \omega\]

Summing over all charts:
\[\int_M d\omega = \sum_{\alpha: U_\alpha \cap \partial M \neq \emptyset} \int_{\partial M \cap U_\alpha} \theta_\alpha \omega = \int_{\partial M} \omega\]
\end{proof}

\textbf{Special cases:}
\begin{itemize}
    \item $k = 1$ (curves): $\int_a^b df = f(b) - f(a)$ (Fundamental Theorem of Calculus)
    \item $k = 2$ in $\mathbb{R}^2$ (Green's Theorem): $\int_D (Q_x - P_y) \, dA = \oint_{\partial D} P\,dx + Q\,dy$
    \item $k = 2$ surfaces in $\mathbb{R}^3$ (Classical Stokes): $\int_S (\nabla \times \mathbf{F}) \cdot d\mathbf{S} = \oint_{\partial S} \mathbf{F} \cdot d\mathbf{r}$
    \item $k = 3$ in $\mathbb{R}^3$ (Divergence Theorem): $\int_V \nabla \cdot \mathbf{F} \, dV = \int_{\partial V} \mathbf{F} \cdot d\mathbf{S}$
\end{itemize}

\begin{example}
\textbf{(Using Green's Theorem to Compute Area)}
The area of a region $D$ can be computed as:
\[\text{Area}(D) = \frac{1}{2} \oint_{\partial D} x\,dy - y\,dx\]
since $d(\frac{1}{2}(x\,dy - y\,dx)) = dx \wedge dy$.

For the ellipse $\frac{x^2}{a^2} + \frac{y^2}{b^2} = 1$, parametrize $\gamma(t) = (a\cos t, b\sin t)$:
\[\text{Area} = \frac{1}{2} \int_0^{2\pi} [a\cos t \cdot b\cos t - b\sin t \cdot (-a\sin t)]\,dt = \frac{ab}{2} \int_0^{2\pi} dt = \pi ab\]
\end{example}

\begin{example}
\textbf{(Divergence Theorem Application)}
Let $\mathbf{F} = (x^3, y^3, z^3)$ and $V$ be the unit ball. Then $\nabla \cdot \mathbf{F} = 3(x^2 + y^2 + z^2) = 3r^2$.

\[\int_V \nabla \cdot \mathbf{F}\, dV = 3\int_0^{2\pi}\int_0^\pi\int_0^1 r^2 \cdot r^2 \sin\phi\, dr\,d\phi\,d\theta = 3 \cdot 2\pi \cdot 2 \cdot \frac{1}{5} = \frac{12\pi}{5}\]
\end{example}

\section{De Rham Cohomology}

\begin{defn}{De Rham Cohomology}{derham}
The \textbf{$k$th de Rham cohomology} of a manifold $M$ is:
\[
H^k(M; \mathbb{R}) = \frac{\{\text{closed } k\text{-forms}\}}{\{\text{exact } k\text{-forms}\}} = \frac{\ker(d: \Omega^k \to \Omega^{k+1})}{\text{Im}(d: \Omega^{k-1} \to \Omega^k)}
\]
\end{defn}

\begin{thm}{Poincaré Lemma}{poincare_cohom}
If $U$ is contractible, then $H^k(U; \mathbb{R}) = 0$ for $k \geq 1$.
\end{thm}

\begin{thm}{Homotopy Invariance}{homotopy_inv_cohom}
If $f, g: M \to N$ are homotopic, then $f^* = g^*: H^k(N) \to H^k(M)$.
\end{thm}

\begin{corollary}
Homotopy equivalent spaces have isomorphic cohomology.
\end{corollary}

\begin{example}
$H^1(S^1; \mathbb{R}) = \mathbb{R}$, generated by $d\theta$. This shows $S^1$ is not contractible.
\end{example}

\section{Applications}

\begin{thm}{No Retraction}{no_retract}
There is no continuous retraction from $D^n$ onto $S^{n-1}$.
\end{thm}

\begin{proof}
If $r: D^n \to S^{n-1}$ is a retraction with $r|_{S^{n-1}} = \text{id}$, let $\omega$ be a closed non-exact form on $S^{n-1}$. Then $r^*\omega$ is closed on $D^n$, hence exact: $r^*\omega = d\beta$. But $\int_{S^{n-1}} \omega = \int_{S^{n-1}} r^*\omega = \int_{D^n} d(r^*\omega) = 0$ by Stokes, contradicting $\omega$ being non-exact.
\end{proof}

\begin{thm}{Brouwer Fixed Point Theorem}{brouwer}
Every continuous map $f: D^n \to D^n$ has a fixed point.
\end{thm}

\begin{proof}
If $f(x) \neq x$ for all $x$, define $r(x)$ as the intersection of the ray from $f(x)$ through $x$ with $S^{n-1}$. This gives a retraction, contradicting the previous theorem.
\end{proof}

\begin{exercises}

\begin{exercise}
Compute $\int_C (x^2 - y)dx + (x - y^2)dy$ where $C$ is the ellipse $\frac{x^2}{a^2} + \frac{y^2}{b^2} = 1$ oriented counter-clockwise.
\end{exercise}

\begin{exercise}
Let $f: U \to \mathbb{R}$ be harmonic ($\Delta f = 0$) on an open set $U \subseteq \mathbb{R}^2$. Prove that $f(x,y)$ equals the average of $f$ on any circle around $(x,y)$ contained in $U$.
\end{exercise}

\begin{exercise}
For $\mathbf{v} = (x^2 + y^2 + z^2)(x\partial_x + y\partial_y + z\partial_z)$, compute $\text{Work}_\Gamma(\mathbf{v})$ where $\Gamma$ connects $A = (a_1, a_2, a_3)$ to $B = (b_1, b_2, b_3)$.
\end{exercise}

\begin{exercise}
Compute $\int_S \frac{dy \wedge dz}{x} + \frac{dz \wedge dx}{y} + \frac{dx \wedge dy}{z}$ where $S$ is the ellipsoid $\frac{x^2}{a^2} + \frac{y^2}{b^2} + \frac{z^2}{c^2} = 1$.
\end{exercise}

\begin{exercise}
Use Stokes' theorem to compute $\int_C (y+z)dx + (z+x)dy + (x+y)dz$ where $C$ is the curve $x = \sin^2 t$, $y = 2\sin t \cos t$, $z = \cos^2 t$ for $t \in [0, \pi]$.
\end{exercise}

\end{exercises}

%%%%%%%%%%%%%%%%%%%%%%%%%%%%%%%%%%%%%%%%%%%%%%%%%%%%%%%%%%%%%%%%%%%%%%%%%%%%%%%
% ADVANCED TOPICS FROM MATH 62 (Covered in lecture, on exam)
%%%%%%%%%%%%%%%%%%%%%%%%%%%%%%%%%%%%%%%%%%%%%%%%%%%%%%%%%%%%%%%%%%%%%%%%%%%%%%%

\chapter{Advanced Topics: Degree Theory and Linking Numbers}

\section{Degree of a Map}

The degree is a powerful topological invariant that counts (with signs) how many times a map ``wraps'' one manifold around another.

\begin{defn}{Degree}{degree62}
Let $f: M \to N$ be a smooth map between compact, connected, oriented $n$-manifolds without boundary. The \textbf{degree} of $f$ is:
\[
\deg(f) = \int_M f^*\omega \bigg/ \int_N \omega
\]
where $\omega$ is any $n$-form on $N$ with $\int_N \omega \neq 0$.
\end{defn}

\begin{thm}{Degree is an Integer}{deg_integer62}
The degree is always an integer.
\end{thm}

\begin{thm}{Local Formula for Degree}{deg_local62}
If $b \in N$ is a regular value of $f$, then:
\[
\deg(f) = \sum_{a \in f^{-1}(b)} \text{sign}(\det(d_a f))
\]
This is the signed count of preimages.
\end{thm}

\begin{thm}{Homotopy Invariance}{homotopy_inv_deg62}
If $f_0, f_1: M \to N$ are homotopic, then $\deg(f_0) = \deg(f_1)$.
\end{thm}

\begin{example}
For $f: S^1 \to S^1$ given by $f(e^{i\theta}) = e^{in\theta}$, we have $\deg(f) = n$. This is the ``winding number.''
\end{example}

\begin{example}
Let $A$ be a $3 \times 3$ matrix with $\det(A) \neq 0$. Define $F_A: S^2 \to S^2$ by $F_A(x) = \frac{Ax}{\|Ax\|}$.

Since $A$ is invertible, $F_A$ is a diffeomorphism onto its image. The degree depends on orientation: $\deg(F_A) = \text{sign}(\det(A))$.
\end{example}

\section{Winding and Linking Numbers}

\begin{defn}{Winding Number}{winding62}
Let $\gamma: [0, 2\pi] \to \mathbb{R}^2 \setminus \{p\}$ be a closed loop. The \textbf{winding number} of $\gamma$ around $p$ is:
\[
w(\gamma, p) = \frac{1}{2\pi} \oint_\gamma \frac{-(y-p_2)dx + (x-p_1)dy}{(x-p_1)^2 + (y-p_2)^2}
\]
\end{defn}

\begin{thm}{Winding as Degree}{winding_deg62}
The winding number equals the degree of the map $\gamma/|\gamma - p|: S^1 \to S^1$.
\end{thm}

\begin{defn}{Linking Number}{linking62}
Let $\phi, \psi: S^1 \to \mathbb{R}^3$ be two disjoint closed curves. The \textbf{linking number} $\text{lk}(\phi, \psi)$ measures how many times one curve winds around the other.

If $D$ is a disk bounded by $\phi$, then:
\[
\text{lk}(\phi, \psi) = \sum_{\psi(t) \in D} \text{sign}(\psi'(t) \cdot n)
\]
where $n$ is the normal to $D$.
\end{defn}

\begin{thm}{Gauss Linking Integral}{gauss_linking62}
\[
\text{lk}(\phi, \psi) = \frac{1}{4\pi} \oint_\phi \oint_\psi \frac{(\phi - \psi) \cdot (d\phi \times d\psi)}{|\phi - \psi|^3}
\]
\end{thm}

\begin{example}
\textbf{(Hopf Link)}
The Hopf link consists of two circles in $\mathbb{R}^3$ that are linked once. Their linking number is $\pm 1$ (depending on orientations).
\end{example}

\section{Higher Connectivity}

\begin{defn}{$k$-Connected}{k_connected62}
A space $A$ is \textbf{$k$-connected} if any continuous map $f: S^m \to A$ for $m \leq k$ is homotopic to a constant map.

Equivalently, $H^i(A; \mathbb{R}) = 0$ for all $i \leq k$.
\end{defn}

\begin{thm}{Connectivity of Spheres}{sphere_conn62}
$\mathbb{R}^{n+1} \setminus \{0\}$ (equivalently, $S^n$) is $(n-1)$-connected but not $n$-connected.
\end{thm}

\begin{proof}[Sketch]
$\mathbb{R}^{n+1} \setminus \{0\}$ is homotopy equivalent to $S^n$ via $x \mapsto x/\|x\|$.

For $m < n$: any map $S^m \to S^n$ can be deformed to miss a point, hence is homotopic to a constant (the target minus a point is contractible).

For $m = n$: the identity map $S^n \to S^n$ has degree 1, so it's not homotopic to a constant (which has degree 0).
\end{proof}

\begin{thm}{Cohomology of $T^2$}{cohom_torus62}
For the torus $T^2 = S^1 \times S^1$ with coordinates $\phi_1, \phi_2$:
\begin{itemize}
    \item $H^0(T^2) = \mathbb{R}$, generated by the constant 1
    \item $H^1(T^2) = \mathbb{R}^2$, generated by $[d\phi_1]$ and $[d\phi_2]$
    \item $H^2(T^2) = \mathbb{R}$, generated by $[d\phi_1 \wedge d\phi_2]$
\end{itemize}
\end{thm}

\section{Lie Derivatives (Brief Overview)}

The Lie derivative measures how a tensor field changes along the flow of a vector field.

\begin{defn}{Flow of a Vector Field}{flow}
Let $X$ be a vector field on $U \subseteq \mathbb{R}^n$. The \textbf{flow} of $X$ is the family of diffeomorphisms $\phi_t: U \to U$ satisfying:
\[\frac{d}{dt} \phi_t(p) = X(\phi_t(p)), \quad \phi_0 = \text{id}\]
Intuitively, $\phi_t(p)$ follows the integral curve of $X$ starting at $p$ for time $t$.
\end{defn}

\begin{defn}{Lie Derivative of a Function}{lie_fn}
For a smooth function $f$ and vector field $X$:
\[\mathcal{L}_X f = \frac{d}{dt}\Big|_{t=0} (\phi_t^* f) = \frac{d}{dt}\Big|_{t=0} f \circ \phi_t = X(f) = df(X)\]
This is just the directional derivative of $f$ along $X$.
\end{defn}

\begin{defn}{Lie Derivative of a Differential Form}{lie_form}
For a $k$-form $\omega$ and vector field $X$:
\[\mathcal{L}_X \omega = \frac{d}{dt}\Big|_{t=0} \phi_t^* \omega\]
This measures the ``rate of change'' of $\omega$ as we flow along $X$.
\end{defn}

\begin{thm}{Cartan's Magic Formula}{cartan}
For a vector field $X$ and $k$-form $\omega$:
\[\mathcal{L}_X \omega = d(\iota_X \omega) + \iota_X (d\omega)\]
where $\iota_X$ is the \textbf{interior product} (contraction): $(\iota_X \omega)(v_1, \ldots, v_{k-1}) = \omega(X, v_1, \ldots, v_{k-1})$.
\end{thm}

\begin{remark}
Cartan's formula is incredibly useful because it expresses the Lie derivative purely in terms of $d$ and $\iota_X$, without needing to compute the flow explicitly.
\end{remark}

\begin{defn}{Lie Bracket}{lie_bracket}
For vector fields $X, Y$, the \textbf{Lie bracket} $[X, Y]$ is the vector field defined by:
\[[X, Y](f) = X(Y(f)) - Y(X(f))\]
This measures the ``failure of commutativity'' of flows: if $[X, Y] = 0$, the flows of $X$ and $Y$ commute.
\end{defn}

\begin{thm}{Properties of Lie Derivative}{lie_props}
\begin{enumerate}
    \item $\mathcal{L}_X$ is a derivation: $\mathcal{L}_X(\omega \wedge \eta) = (\mathcal{L}_X \omega) \wedge \eta + \omega \wedge (\mathcal{L}_X \eta)$
    \item $\mathcal{L}_X$ commutes with $d$: $\mathcal{L}_X(d\omega) = d(\mathcal{L}_X \omega)$
    \item $[\mathcal{L}_X, \mathcal{L}_Y] = \mathcal{L}_{[X,Y]}$ (Jacobi identity)
\end{enumerate}
\end{thm}

\section{Phase Flow of Vector Fields}

\begin{defn}{Integral Curve}{integral_curve}
An \textbf{integral curve} of a vector field $X$ is a curve $\gamma: I \to M$ such that
\[\gamma'(t) = X(\gamma(t))\]
This is just an ODE! So by existence/uniqueness, through every point there's a unique maximal integral curve.
\end{defn}

\begin{defn}{Phase Flow}{phase_flow}
The \textbf{phase flow} of $X$ is the family of maps $\phi_t: M \to M$ given by
\[\phi_t(p) = \gamma_p(t)\]
where $\gamma_p$ is the integral curve starting at $p$, i.e., $\gamma_p(0) = p$.
\end{defn}

\begin{thm}{Properties of Flow}{flow_props}
\begin{enumerate}
    \item $\phi_0 = \text{id}$
    \item $\phi_s \circ \phi_t = \phi_{s+t}$ (group property)
    \item $\frac{d}{dt}\big|_{t=0} \phi_t(p) = X(p)$
    \item $\mathcal{L}_X \omega = \frac{d}{dt}\big|_{t=0} \phi_t^* \omega$
\end{enumerate}
\end{thm}

\begin{proof}[Sketch of (4)]
This is the key insight: the Lie derivative measures ``infinitesimal change'' along the flow:
\[\mathcal{L}_X \omega = \lim_{t \to 0} \frac{\phi_t^* \omega - \omega}{t} = \frac{d}{dt}\Big|_{t=0} \phi_t^* \omega\]
\end{proof}

\section{Divergence-Free Vector Fields}

\begin{defn}{Divergence-Free}{div_free}
A vector field $X$ is \textbf{divergence-free} (or \textbf{incompressible}) if $\mathcal{L}_X \omega = 0$ where $\omega$ is the volume form.

By Cartan: $\mathcal{L}_X \omega = d(\iota_X \omega) + \iota_X(d\omega) = d(\iota_X \omega)$ since $d\omega = 0$ (it's a top form). So divergence-free means $d(\iota_X \omega) = 0$.
\end{defn}

\begin{thm}{Volume Preservation}{vol_preserve}
$X$ is divergence-free if and only if its flow $\phi_t$ preserves volume:
\[\phi_t^* \omega = \omega\]
for all $t$.
\end{thm}

\begin{proof}
$\frac{d}{dt} \phi_t^* \omega = \phi_t^* (\mathcal{L}_X \omega)$. If $\mathcal{L}_X \omega = 0$, then $\phi_t^* \omega$ is constant, equal to $\phi_0^* \omega = \omega$.
\end{proof}

\begin{remark}
In $\mathbb{R}^3$: $X$ is divergence-free iff $\nabla \cdot X = 0$. This is the condition for incompressible fluid flow.
\end{remark}

\section{Contact Structures}

\begin{defn}{Contact Form}{contact_form}
A 1-form $\alpha$ on a $(2n+1)$-dimensional manifold is a \textbf{contact form} if
\[\alpha \wedge (d\alpha)^n \neq 0\]
everywhere. The hyperplane field $\xi = \ker \alpha$ is called a \textbf{contact structure}.
\end{defn}

\begin{example}
On $\mathbb{R}^3$ with coordinates $(x, y, z)$:
\[\alpha = dz - y\,dx\]
Then $d\alpha = -dy \wedge dx = dx \wedge dy$, and $\alpha \wedge d\alpha = dz \wedge dx \wedge dy \neq 0$.
\end{example}

\begin{remark}
Contact structures are ``maximally non-integrable'': $\alpha \wedge d\alpha \neq 0$ means the hyperplane field $\ker \alpha$ cannot be tangent to any surface. This is the opposite of what happens when $\alpha \wedge d\alpha = 0$. Contact geometry is the odd-dimensional analog of symplectic geometry.
\end{remark}

\begin{exercises}

\begin{exercise}
Prove that $\mathbb{R}^2 \setminus \{(1,0), (-1,0)\}$ is not diffeomorphic to $\mathbb{R}^2 \setminus \{(0,0)\}$.
\end{exercise}

\begin{exercise}
Show that there are no degree 1 maps $f: S^2 \to T^2$.
\end{exercise}

\begin{exercise}
Let $A$ be a $3 \times 3$ matrix with $\det(A) \neq 0$. Compute $\deg(F_A)$ where $F_A(x) = Ax/\|Ax\|$ for $x \in S^2$.
\end{exercise}

\begin{exercise}
Let $\gamma = (\gamma_1, \gamma_2): [0, 2\pi] \to \mathbb{R}^2 \setminus \{(-1,0), (1,0)\}$ be a loop with winding number 2 around $(-1, 0)$ and $-3$ around $(1, 0)$. Compute $\text{lk}(\phi, \psi)$ where $\phi(s) = (\cos s, \sin s, 0)$ and $\psi(t) = (\gamma_1(t), 0, \gamma_2(t))$.
\end{exercise}

\end{exercises}

%%%%%%%%%%%%%%%%%%%%%%%%%%%%%%%%%%%%%%%%%%%%%%%%%%%%%%%%%%%%%%%%%%%%%%%%%%%%%%%
% SOLUTIONS TO MATH 62 EXERCISES
%%%%%%%%%%%%%%%%%%%%%%%%%%%%%%%%%%%%%%%%%%%%%%%%%%%%%%%%%%%%%%%%%%%%%%%%%%%%%%%

\backmatter % Solutions should be unnumbered

\chapter{Solutions to Math 62 Exercises}

\section*{Chapter 1: Dual Spaces and Topology}

\begin{solution}[1.1]
Given a basis $\{\ell_1, \ldots, \ell_n\}$ of $V^*$, we want to construct vectors $w_1, \ldots, w_n \in V$ with $\ell_i(w_j) = \delta_{ij}$.

Define $\Phi: V \to \mathbb{R}^n$ by $\Phi(v) = (\ell_1(v), \ldots, \ell_n(v))$. We claim $\Phi$ is an isomorphism.

\textbf{Injective:} If $\Phi(v) = 0$, then $\ell_i(v) = 0$ for all $i$. Since $\{\ell_i\}$ is a basis for $V^*$, this means $f(v) = 0$ for all $f \in V^*$, so $v = 0$.

\textbf{Surjective:} $\dim V = \dim V^* = n = \dim \mathbb{R}^n$, and $\Phi$ is injective, so $\Phi$ is bijective.

Now let $e_1, \ldots, e_n$ be the standard basis of $\mathbb{R}^n$. Define $w_j = \Phi^{-1}(e_j)$. Then:
\[
\ell_i(w_j) = (\Phi(w_j))_i = (e_j)_i = \delta_{ij}
\]
\end{solution}

\begin{solution}[1.3]
\textbf{For $\mathcal{D}$:} We need $\ker \mathcal{D} = \{0\}$. If $\mathcal{D}(v) = 0$, then $\langle v, x \rangle = 0$ for all $x \in V$. Taking $x = v$ gives $\|v\|^2 = 0$, so $v = 0$. Hence $\mathcal{D}$ is injective.

\textbf{For $i_S$:} Suppose $i_S(v) = 0$, meaning $v^*(x) = 0$ for all $x$. Write $v = \sum a_k v_k$ in the basis. If $v \neq 0$, some $a_k \neq 0$, and $v^*(v_k) = a_k \neq 0$, contradiction. So $v = 0$.
\end{solution}

\begin{solution}[1.5]
Let $A$ be connected and suppose $\overline{A}$ is disconnected: $\overline{A} = U \cup V$ with $U, V$ open in $\overline{A}$, disjoint, and non-empty.

Since $A \subseteq \overline{A}$, we have $A = (A \cap U) \cup (A \cap V)$. These are relatively open in $A$ and disjoint. 

If both $A \cap U$ and $A \cap V$ are non-empty, then $A$ is disconnected, contradiction.

So WLOG, $A \subseteq U$. Then $\overline{A} \subseteq \overline{U}$. But $U$ is closed in $\overline{A}$ (complement of open $V$), so $\overline{U} \cap \overline{A} = U$. Hence $\overline{A} \subseteq U$, meaning $V = \emptyset$, contradiction.
\end{solution}

\begin{solution}[1.7]
\textbf{(a)} $[0,1]$ vs $[0,1)$: Remove two points from $[0,1]$ to disconnect it (remove interior points). But removing any two points from $[0,1)$: if we remove $0$ and another point $a$, we get $(0, a) \cup (a, 1)$ which is disconnected. If we remove two interior points, we get three components. The key is that $[0,1]$ has exactly 2 ``special'' endpoints that don't disconnect when removed, while $[0,1)$ has only 1.

$[0,1)$ vs $(0,1)$: Removing one point from $[0,1)$: if we remove $0$, it stays connected. If we remove an interior point, it disconnects. For $(0,1)$, removing \emph{any} point disconnects it. So they're not homeomorphic.

\textbf{(b)} Remove a point from $\mathbb{R}^n$ ($n > 1$): the result $\mathbb{R}^n \setminus \{p\}$ is still path-connected (go around the point). But removing a point from $\mathbb{R}$ gives two disconnected half-lines. A homeomorphism would preserve connectedness of point-removals.
\end{solution}

\section*{Chapter 2: Differential Forms}

\begin{solution}[2.1]
$f(x,y) = (e^x - 2y, x + y^2, \sin y + 1)$, $\alpha = z\,dx - 3\,dy + e^y\,dz$.

Using $f^*(dg) = d(g \circ f)$ and linearity:
\begin{align*}
f^*\alpha &= (\sin y + 1)d(e^x - 2y) - 3d(x + y^2) + e^{x+y^2}d(\sin y + 1) \\
&= (\sin y + 1)(e^x dx - 2dy) - 3(dx + 2y\,dy) + e^{x+y^2}\cos y\,dy \\
&= (e^x\sin y + e^x - 3)dx + (-2\sin y - 2 - 6y + e^{x+y^2}\cos y)dy
\end{align*}
\end{solution}

\begin{solution}[2.3]
The projection gives $x = y^2$, so parameterize as $(y^2, y, h(y))$. The plane field condition $dz - y\,dx = 0$ becomes:
\[
h'(y)\,dy - y \cdot 2y\,dy = 0 \implies h'(y) = 2y^2 \implies h(y) = \frac{2}{3}y^3 + C
\]
With $(0,0,0) \in \Gamma$: $h(0) = 0$, so $C = 0$.

The curve is $\Gamma(t) = (t^2, t, \frac{2}{3}t^3)$.
\end{solution}

\begin{solution}[2.5]
Consider $\alpha = \frac{-z\,dy + y\,dz}{y^2 + z^2}$ on $\mathbb{R}^3 \setminus P$.

Let $\Gamma = \{x = 0, y^2 + z^2 = 1\}$ parameterized by $\gamma(t) = (0, \cos t, \sin t)$.
\[
\int_\Gamma \alpha = \int_0^{2\pi} \frac{-\sin t(-\sin t) + \cos t \cdot \cos t}{1}\,dt = \int_0^{2\pi} 1\,dt = 2\pi \neq 0
\]
Since $\alpha$ is closed (check: $d\alpha = 0$) but $\oint \alpha \neq 0$, $\alpha$ is not exact, so the domain is not simply connected.
\end{solution}

\section*{Chapter 3: Multilinear Algebra}

\begin{solution}[3.1]
$B(X, Y) = \text{Tr}(XY)$.

\textbf{Symmetry:} $\text{Tr}(XY) = \sum_{i,j} X_{ij}Y_{ji} = \sum_{i,j} Y_{ji}X_{ij} = \text{Tr}(YX)$.

\textbf{Rank and Index:} $Q(X) = \text{Tr}(X^2) = \sum_{i,j} X_{ij}X_{ji}$.

For diagonal entries: $\sum_i X_{ii}^2$ (all positive contributions).

For off-diagonal: $\sum_{i<j} 2X_{ij}X_{ji}$. Using the identity $(X_{ij} + X_{ji})^2 - (X_{ij} - X_{ji})^2 = 4X_{ij}X_{ji}$, we get equal numbers of positive and negative contributions.

Result: $n_+ = n + \binom{n}{2} = \frac{n(n+1)}{2}$, $n_- = \binom{n}{2} = \frac{n(n-1)}{2}$, rank $= n^2$.
\end{solution}

\begin{solution}[3.3]
\textbf{(a)} Let $\dim V = n$ (odd) and $\omega$ a 2-form with matrix $A$. Since $\omega$ is skew-symmetric, $A^T = -A$. Then:
\[
\det A = \det(A^T) = \det(-A) = (-1)^n \det A = -\det A
\]
So $\det A = 0$, meaning $\omega$ is degenerate.

\textbf{(b)} The rank of a skew-symmetric bilinear form equals the dimension of the non-degenerate part. By (a), this must be even-dimensional (else it would be degenerate). Hence the rank is even.
\end{solution}

\section*{Chapter 5: Manifolds}

\begin{solution}[5.1]
\textbf{(a) $SO(3)$:} Define $F: \mathcal{M}_{3,3} \to \text{Sym}(3)$ by $F(A) = A^T A$. Then $O(3) = F^{-1}(I)$.

The derivative: $dF_A(H) = A^T H + H^T A$.

For $A \in O(3)$ and any symmetric $S$, take $H = \frac{1}{2}AS$. Then $A^T H + H^T A = \frac{1}{2}S + \frac{1}{2}S^T = S$.

So $dF_A$ is surjective, making $O(3)$ (and thus $SO(3)$) a 3-dimensional submanifold.

\textbf{$UT(S^2)$:} Define $G(x,y) = (x \cdot x, y \cdot y, x \cdot y)$. Then $UT(S^2) = G^{-1}(1,1,0)$.

$dG_{(x,y)} = \begin{pmatrix} 2x^T & 0 \\ 0 & 2y^T \\ y^T & x^T \end{pmatrix}$ has full rank when $x, y$ are unit vectors with $x \perp y$.

\textbf{(b)} Define $\phi: SO(3) \to UT(S^2)$ by $\phi([v_1 | v_2 | v_3]) = (v_1, v_2)$.

This is well-defined since columns of an orthogonal matrix are orthonormal.

Inverse: $\phi^{-1}(x, y) = [x | y | x \times y]$. The cross product gives a unit vector orthogonal to both, with the correct orientation (determinant = 1).

Both maps are smooth (linear/polynomial), so this is a diffeomorphism.
\end{solution}

\section*{Chapter 6: Stokes' Theorem}

\begin{solution}[6.1]
The ellipse is parameterized by $\gamma(t) = (a\cos t, b\sin t)$ for $t \in [0, 2\pi]$.

Using Green's theorem with $P = x^2 - y$, $Q = x - y^2$:
\[
\oint_C = \iint_D \left(\frac{\partial Q}{\partial x} - \frac{\partial P}{\partial y}\right) dA = \iint_D (1 - (-1))\,dA = 2 \cdot \text{Area}(D) = 2\pi ab
\]
\end{solution}

\begin{solution}[6.3]
The work is $\int_\Gamma \omega$ where $\omega = (x^2 + y^2 + z^2)(x\,dx + y\,dy + z\,dz)$.

Note that $x\,dx + y\,dy + z\,dz = \frac{1}{2}d(x^2 + y^2 + z^2)$.

Let $u = x^2 + y^2 + z^2$. Then $\omega = u \cdot \frac{1}{2}du = \frac{1}{2}d(\frac{u^2}{2}) = d(\frac{u^2}{4})$.

So $f = \frac{1}{4}(x^2 + y^2 + z^2)^2$ is a primitive, and:
\[
\text{Work} = f(B) - f(A) = \frac{1}{4}(b_1^2 + b_2^2 + b_3^2)^2 - \frac{1}{4}(a_1^2 + a_2^2 + a_3^2)^2
\]
\end{solution}

\begin{solution}[6.5]
Let $\omega = (y+z)dx + (z+x)dy + (x+y)dz$.

Compute: $d\omega = (dy + dz) \wedge dx + (dz + dx) \wedge dy + (dx + dy) \wedge dz$.

Expanding: $= dy \wedge dx + dz \wedge dx + dz \wedge dy + dx \wedge dy + dx \wedge dz + dy \wedge dz$

$= -dx \wedge dy + dx \wedge dy - dx \wedge dz + dx \wedge dz - dy \wedge dz + dy \wedge dz = 0$.

So $\omega$ is closed. The curve lies on the plane $x + z = 1$ (since $\sin^2 t + \cos^2 t = 1$), which is simply connected.

By Poincaré lemma, $\omega = d\beta$ for some $\beta$. By Stokes: $\oint_C \omega = \int_S d\omega = 0$.
\end{solution}

\section*{Chapter 7: Advanced Topics}

\begin{solution}[Degree of $F_A$]
Since $A$ is invertible, $F_A: S^2 \to S^2$ is smooth and surjective.

For a regular value $b \in S^2$, the preimage $F_A^{-1}(b)$ consists of the points $x \in S^2$ where $Ax$ is parallel to $b$. Since $A$ is invertible, there's exactly one such line direction, giving two points $\pm x_0$ on $S^2$.

At each preimage point, the sign of $\det(d_{x_0} F_A)$ depends on whether $A$ preserves or reverses orientation. Since both points contribute the same sign:
\[
\deg(F_A) = \text{sign}(\det(A))
\]
\end{solution}

\begin{solution}[$\mathbb{R}^2$ minus two points]
We show $\dim H^1(U) > 1$ where $U = \mathbb{R}^2 \setminus \{(\pm 1, 0)\}$.

Consider two 1-forms:
\[
\alpha_1 = \frac{-(y)dx + (x-1)dy}{(x-1)^2 + y^2}, \quad \alpha_2 = \frac{-(y)dx + (x+1)dy}{(x+1)^2 + y^2}
\]
Both are closed on $U$. Consider loops $\gamma_1$ around $(1, 0)$ only and $\gamma_2$ around $(-1, 0)$ only.

Then $\oint_{\gamma_1} \alpha_1 = 2\pi$, $\oint_{\gamma_1} \alpha_2 = 0$ (and similarly for $\gamma_2$).

So $[\alpha_1]$ and $[\alpha_2]$ are linearly independent in $H^1(U)$, giving $\dim H^1(U) \geq 2$.

But $\mathbb{R}^2 \setminus \{0\}$ has $\dim H^1 = 1$. Since cohomology is a diffeomorphism invariant, these spaces are not diffeomorphic.
\end{solution}

\begin{solution}[No degree 1 maps $S^2 \to T^2$]
Suppose $f: S^2 \to T^2$ has degree 1. Then $f^*: H^2(T^2) \to H^2(S^2)$ is an isomorphism.

Now $H^1(T^2) = \mathbb{R}^2$ with generators $[d\phi_1], [d\phi_2]$, and $[d\phi_1] \wedge [d\phi_2] = [\eta]$ generates $H^2(T^2)$.

For $f^*[\eta] \neq 0$, we need $f^*[d\phi_1] \wedge f^*[d\phi_2] \neq 0$ in $H^2(S^2)$.

But $f^*[d\phi_i] \in H^1(S^2) = 0$ (since $S^2$ is simply connected).

So $f^*[d\phi_i] = 0$, hence $f^*[\eta] = 0$, contradicting degree 1.
\end{solution}

\begin{solution}[Linking number]
The linking number $\text{lk}(\phi, \psi)$ counts signed intersections of $\psi$ with any disk bounded by $\phi$.

$\phi$ bounds the disk $D = \{(x, y, 0) : x^2 + y^2 \leq 1\}$.

$\psi$ intersects $D$ when $\gamma_2(t) = 0$ (the $z$-coordinate is 0) and $|\gamma_1(t)| < 1$ (inside the unit disk).

The sign at each intersection is $\text{sign}(\gamma_2'(t))$ (positive if crossing upward).

The winding number around $(-1, 0)$ counts signed crossings of the negative $x$-axis by $\gamma$, which is the sum of signs for $\gamma_1 < -1$.

The winding number around $(1, 0)$ counts signed crossings for $\gamma_1 > 1$.

The linking number counts crossings for $-1 < \gamma_1 < 1$:
\[
\text{lk}(\phi, \psi) = w(\gamma, (-1,0)) - w(\gamma, (1,0)) = 2 - (-3) = 5
\]
\end{solution}

\end{document}